\documentclass[11pt,a4paper]{article}
\usepackage[T1]{fontenc}
\usepackage[utf8]{inputenc}
\usepackage{lmodern,microtype}
\usepackage[margin=25mm]{geometry}
\usepackage{amsmath,amssymb,amsthm,mathtools,mathrsfs}
\pdfmapfile{+euler.map}
\pdfmapfile{+rsfs.map}
\usepackage{enumitem,xurl,xcolor}
\usepackage[colorlinks=true,linkcolor=blue!45!black,citecolor=blue!45!black,urlcolor=blue!45!black]{hyperref}
\allowdisplaybreaks[2]
\setlength{\emergencystretch}{2em}
\numberwithin{equation}{section}
\newtheorem{theorem}{Theorem}[section]
\newtheorem{proposition}[theorem]{Proposition}
\newtheorem{lemma}[theorem]{Lemma}
\newtheorem{corollary}[theorem]{Corollary}
\newcommand{\one}{\mathbf1}
\newcommand{\QQ}{\mathbb Q}
\newcommand{\lcm}{\operatorname{lcm}}
\newcommand{\sVec}{\mathrm{sVec}}
\newcommand{\Ffamily}{\mathscr F}
\hypersetup{pdftitle={Braided fusion categories with dimension dividing p6q6r4},pdfauthor={},bookmarksdepth=2}
\title{Braided fusion categories with dimension\\dividing \(p^6q^6r^4\)}
\author{}
\date{}
\begin{document}
\maketitle
\begin{abstract}
Let $p,q,r$ be distinct primes. Every braided fusion category over
$\mathbb C$ of Frobenius--Perron dimension a proper divisor of $p^6q^6r^4$ is weakly
group-theoretical. At the full dimension, a possible obstruction is
integral, perfect, fusion-simple and modular, and the quartic exponent
lies on the largest prime. The sole possible even dimension is $5400^2$.
In odd dimension every category in this family below $5^6\,7^6\,23^4$
is solvable. The proof uses local modular actions to find two simple
objects whose dimensions have prescribed prime divisors. Balancing
constrains their $S$-entry to a short positive sum of roots of unity.
At the first odd parameter allowed by the resulting scalar inequality,
the forced word has four phases whose six pairwise ratios are
distinct up to inversion.
\end{abstract}

\section{Introduction}\label{sec:introduction}

A fusion category is \emph{weakly group-theoretical} if it is
categorically Morita equivalent to a nilpotent fusion category.
A category is nilpotent when iterating the adjoint construction,
which takes the subcategory generated by $X\otimes X^*$ for simple $X$,
eventually gives $\mathrm{Vec}$. Weak group-theoreticity is preserved
by the extension, equivariantization and decomposition operations
that organize the structure theory of fusion categories
\cite{ENOsolvable,NataleWGT}. Etingof, Nikshych and Ostrik asked
whether every weakly integral fusion category is weakly group-theoretical
\cite[Section~11, Question~2]{ENOsolvable}. For braided categories,
Dong proved this for dimensions $p^a q^b d$, where $d$ is square-free
and coprime to $pq$, and for $p^2q^2r^2$ \cite{DongWeak}.
Here all three primes may occur to high powers. Their relative sizes
will matter as well as their exponents.

We consider braided categories over $\mathbb C$ whose
Frobenius--Perron dimensions lie in
\[
 \Ffamily=\{D\in\mathbb Z_{>0}:D\mid p^6q^6r^4
                    \text{ for distinct primes }p,q,r\}.
\]
No ordering of the primes is imposed in this definition. A category
is \emph{perfect} if the unit is its only invertible simple object,
and \emph{fusion-simple} if it has no nontrivial proper fusion
subcategory. Modular means that its symmetric center is trivial.

\begin{theorem}\label{thm:main}
Let $p,q,r$ be distinct primes. All categories below are braided fusion
categories over $\mathbb C$, and dimension means Frobenius--Perron
dimension. Let $\Ffamily$ be the set of positive integer divisors of
$a^6b^6c^4$ as $a,b,c$ range over distinct primes.
The following statements hold.
\begin{enumerate}[label=\textup{(\roman*)},leftmargin=*]
\item Every braided fusion category of dimension a proper divisor of
$p^6q^6r^4$ is weakly group-theoretical. If a category of the full
dimension is not weakly group-theoretical, then it is integral,
perfect, fusion-simple and modular, and $r>\max(p,q)$.
\item Every braided fusion category of even dimension in $\Ffamily$
is weakly group-theoretical, except possibly at
\[
                  2^6\,3^6\,5^4=5400^2=29\,160\,000.
\]
\item Every braided fusion category of odd dimension in $\Ffamily$
strictly below
\[
       D_0=5^6\,7^6\,23^4=514\,422\,090\,765\,625
                         =(22\,680\,875)^2
\]
is weakly group-theoretical.
\end{enumerate}
\end{theorem}

By Natale and Plavnik \cite[Proposition~7.1]{NPdegrees}, an odd-dimensional
weakly group-theoretical fusion category is solvable. This applies in
particular throughout the range of~\textup{(iii)}. Both thresholds concern
$\Ffamily$; the odd bound is not a bound for arbitrary dimensions.
Neither exceptional dimension is excluded or realized here.

\paragraph{Scope and provenance.}
This article is a self-contained extraction of the three-prime part of the
joint project of Burciu, Dong and Palcoux. The family bounds and the
conductor, local-action and support methods belong to that same project;
we make no separate priority claim for their appearance in this shorter
form. All proofs needed for these bounds are included here. The exposition
separates the established cyclotomic unimodularity theorem of
Beck--Ho\c{s}ten \cite[Theorem~13]{BH} from its use to recover the
\emph{actual} positive balancing support, and from the categorical
common-constituent and difference-descent arguments.

\paragraph{How the primes enter the proof.}
Choose a counterexample minimal under divisibility. It gives either an
ordinary modular table or the positive table obtained by identifying
a transparent fermion with the unit. In the ordinary case let $d_Y$
denote the simple dimension, $\chi_J$ the character at a column $J$,
and $P_p$ the set of labels whose characters reduce to the dimension
character at $p$. When the twist of $J$ has nontrivial $p$-part,
local modular invariance gives
\[
               \sum_{Y\in P_p}d_Y\chi_J(Y)=0 .
\]
The unit contributes $1$. Reducing at a different prime $q$ gives a
nonunit $Y$ with $q\nmid d_Y$ and $S_{YJ}\ne0$. Interchanging the
roles of two primes supplies complementary divisibility conditions.

For such a pair $X,Y$, divide its $S$-entry by
$\lcm(d_X,d_Y)$. Balancing expresses the result as a positive rational
sum of roots of unity of total weight $\gcd(d_X,d_Y)$. The two residue
conditions can leave very few possible phases. At the even boundary,
for instance, the common weight is forced to be four. The smaller
odd prime must be $3$, and if the larger prime $r$ is at least $7$,
the only possible nonscalar word is $1+3\xi$, with $\xi$ of nontrivial
$r$-power order. Its two positions are incompatible with the adjoint
products in a minimal table. This last exclusion stops at $r=5$.

Sections~\ref{sec:local-tools} and~\ref{sec:crossings} establish the
local and cyclotomic results. Sections~\ref{sec:odd-faces}
and~\ref{sec:even-faces} eliminate the proper divisors.
Section~\ref{sec:boundaries} treats the two remaining boundaries.
For odd $p<q<r$ it gives $3q\le\min(r,p^2)$, whose first possible
triple is $(5,7,23)$. At that triple an explicit additive support
attains equality in the scalar overlap estimate. A second application
of the adjoint argument excludes this equality case for the word
forced by the transfer theorem (Proposition~\ref{supp:eleven-mask}).

We use the standard basic properties of fusion categories and modular
data from \cite{EGNO}. We fix the orientation of $S$ and the
normalization of the modular representation in
Section~\ref{sec:local-tools}.
Write $d_X=\operatorname{FPdim}(X)$, let $0$ or $\one$ denote the
unit, and write $v_p$ for the $p$-adic valuation. A degree is
\emph{$p$-free} when it is not divisible by $p$. The house of an
algebraic number is the largest modulus of its Galois conjugates.
We occasionally abbreviate weakly group-theoretical to WGT.

\subsection{The form of a minimal obstruction}
Weakly integral categories carry their positive spherical
structure; a braiding then supplies the corresponding ribbon structure.

\begin{lemma}[Minimal obstruction]\label{loc:minimal}
Let $D$ be a positive integer. Suppose that every braided fusion category
whose dimension is a proper divisor of $D$ is weakly group-theoretical.
If a braided category
$\mathcal C$ of dimension $D$ is not weakly group-theoretical, then it is
integral, equals its adjoint subcategory, and has exactly one of the
following forms.
\begin{enumerate}
\item It is perfect, fusion-simple and modular.
\item Its symmetric center and pointed subcategory are both
$\langle g\rangle=\mathrm{sVec}$, and its only fusion subcategories are
$\mathrm{Vec},\mathrm{sVec},\mathcal C$.
\end{enumerate}
In the second case the positive quotient
$K_0(\mathcal C)/(g-1)$ is a perfect simple fusion ring.
\end{lemma}
\begin{proof}
A nontrivial universal grading has neutral component of dimension
$D/|U|$, a proper divisor of $D$. The neutral component is braided, and
weak group-theoreticity passes from it to a graded extension. Hence the
universal grading is trivial. The dimension grading of a weakly integral
category now makes $\mathcal C$ integral
\cite[Theorem~3.10]{GNnil}.

A nontrivial Tannakian subcategory $\mathrm{Rep}(G)$ would give a
$G$-crossed de-equivariantization. If $H\le G$ is its grading support,
the neutral braided component has dimension $D/(|G||H|)$, again a proper
divisor: this rational number is an algebraic integer. Extension and
equivariantization would make $\mathcal C$ weakly group-theoretical.
Thus no such subcategory exists. Likewise a proper nontrivial
nondegenerate subcategory would split off with its centralizer; both
factors have proper divisor dimensions and give the same contradiction.
Here we use the standard permanence and decomposition theorems
\cite[Proposition~4.1]{ENOsolvable}
\cite[Theorems~3.13 and~4.44]{DGNO}
\cite[Proposition~4.1]{NataleWGT}.

A symmetric fusion category without a nontrivial Tannakian subcategory
is $\mathrm{Vec}$ or $\mathrm{sVec}$. Every proper fusion subcategory
of $\mathcal C$ is integral and has proper divisor dimension, hence is
weakly group-theoretical. In the absence of Tannakian subcategories it
is pointed \cite[Theorem~3.1]{NataleCore}. If $\mathcal C$ is modular,
$\mathcal C=\mathcal C_{\rm ad}$ and
$(\mathcal C_{\rm ad})'=\mathcal C_{\rm pt}$ give
$\mathcal C_{\rm pt}=\mathrm{Vec}$. In the slightly degenerate case,
there is neither a nontrivial nondegenerate subcategory nor a nontrivial
Tannakian subcategory, so $\mathcal C_{\rm pt}=\mathrm{sVec}$ by
\cite[Lemma~7.1]{NataleWGT}. These are the two asserted subcategory
lattices.

Finally $\theta_{gX}=-\theta_X$ implies $gX\not\simeq X,X^*$.
Choose representatives of the fermion pairs compatibly with duality.
The quotient has coefficients $N_{XY}^Z+N_{XY}^{gZ}$ and unit
coefficient $\delta_{Y,X^*}$. Its only dimension-one basis element is
the unit. A proper based subring would pull back to a fusion subcategory
strictly between $\mathrm{sVec}$ and $\mathcal C$, which is impossible.
\end{proof}

\section{Character arithmetic and local modular actions}\label{sec:local-tools}

\subsection{Character arithmetic}\label{sec:arithmetic}

We work with integral categories of dimension $D$, using either the ordinary
modular table or the positive fermion-quotient table of Lemma~\ref{loc:minimal}.
Write $M_{X,Y}=c_{Y,X}c_{X,Y}$ for the double braiding on $X\otimes Y$.
We use the modular convention of \cite[\S1.3]{DLNCongruence}:
$S_{XY}$ is the pivotal trace of $M_{X^*,Y}$. Equivalently, it is the
trace of $M_{X,Y}^{-1}$. Indeed the trace with a dualized strand is the
complex conjugate of the forward trace, and the balancing eigenvalues
are roots of unity with positive real multiplicities weighted by dimension.
In particular $S_{0X}=d_X$. We use the actual ribbon twists $\theta_X$.
The table dimension is $\Delta=D$ for an ordinary modular category and
$\Delta=D/2$ for a slightly degenerate category after identification
of its fermion pairs. In either case the coefficient-of-unit trace
$\tau$ and the characters of the commutative based ring satisfy
\begin{equation}\label{arith:data}
 \chi_Y(X)=\frac{S_{XY}}{d_Y},\qquad
 f_Y=\frac{\Delta}{d_Y^2}\in\mathbb Z_{>0},\qquad
 \sum_X\chi_Y(X)\overline{\chi_Z(X)}=\delta_{YZ}f_Y.
\end{equation}
We call $v_p(d_Y)$ the $p$-height of $Y$ and $v_p(f_Y)$ its $p$-defect.
The two alternatives of Lemma~\ref{loc:minimal} will be called the
ordinary and positive reduced minimal tables, respectively.
Here $S$ is symmetric, its entries are cyclotomic integers after
division by either simple dimension, and $S^*S=\Delta I$.
These are the usual modular character identities; their positive
fermion-quotient form follows from \cite[Proposition~3.3 and
Corollary~3.4]{YuSlightlyDegenerate}. The primitive idempotents and trace are
\begin{equation}\label{arith:idempotents}
 f_Ye_Y=\sum_X\overline{\chi_Y(X)}X,
 \qquad \tau(b)=\sum_Y\frac{\chi_Y(b)}{f_Y}.
\end{equation}
The full character field is $K_Y=\mathbb Q(\chi_Y(X):X)$.
Galois conjugation permutes the characters and preserves their
rational squared norms, hence their dimensions. Symmetry gives
\begin{equation}\label{arith:galois}
 \sigma\chi_Y(X)=\chi_{\sigma Y}(X)
                 =\chi_Y(\sigma X),\qquad d_{\sigma X}=d_X.
\end{equation}
The label permutation is a based-ring automorphism: evaluating its
linear extension on a product gives the product of the transformed
character values, and the characters separate elements.
All field traces and all congruences below are taken in these
cyclotomic fields. A congruence at a rational prime means a
congruence at every prime above it unless a place is specified.

\begin{lemma}\label{arith:strict}
In a perfect simple ordinary or positive reduced table:
\begin{enumerate}[label=\textup{(\roman*)},leftmargin=*]
\item $S_{XY}/\operatorname{lcm}(d_X,d_Y)$ is integral, and
      $|\sigma\chi_Y(X)|<d_X$ for nonunit $X,Y$.
\item Coprime nonunit dimensions have zero crossing. No nonunit
      dimension is a prime power. Every nonunit row meets a nonunit
      degree prime to any prescribed dimension prime.
\item $\Delta$ is square-full. If $\Delta$ is odd, no nonunit label
      is self-dual.
\end{enumerate}
\end{lemma}
\begin{proof}
For $s=d_X,t=d_Y$, B\'ezout applied to the two integral numbers
$S_{XY}/s$ and $S_{XY}/t$ proves the first integrality assertion.
The labels on which a character has absolute value equal to the
dimension form a based subring: equality in the triangle inequality
propagates to every constituent of a product. Simplicity makes this
subring trivial or the whole ring. In the latter case its codegree
is $\Delta$, so its character label has dimension one and is the
unit. This proves strictness, including at every Galois conjugate.
If $(s,t)=1$, the integral number $S_{XY}/st$ has every conjugate
of modulus less than one and is therefore zero. The identity
\begin{equation}\label{arith:row-zero}
              1+\sum_{X\ne0}d_X\chi_Y(X)=0\qquad(Y\ne0)
\end{equation}
reduced at a dimension prime supplies a nonunit hit prime to it.
This excludes prime-power degrees.

The ordinary modular case of the following vanishing argument is the
Burnside property of \cite{BPBurnside}. The same proof treats the
positive reduced table. If a nonunit row had no zero, its
Galois-averaged norm and the norm bound for nonzero algebraic integers
would give, with $s=d_Y$,
\[
 \frac{\Delta}{s^2}\ge
 \sum_X\left(\frac{d_X}{\gcd(d_X,s)}\right)^2
 >\sum_X\frac{d_X^2}{s^2}=\frac{\Delta}{s^2},
\]
where the strict inequality already occurs at the unit. Thus
$\prod_X XX^*$ vanishes on every non-dimension character, and
\eqref{arith:idempotents} gives
$\prod_Xd_X^2/\Delta\in\mathbb Z$. Since $d_X^2\mid\Delta$,
no prime occurs to exponent one in $\Delta$.
If $\Delta$ is odd, all $f_Y$ are odd, and
$\tau(b)\equiv\sum_Y\chi_Y(b)=\operatorname{Tr}(N_b)\pmod2$.
For a basis element $X$, the elementary matrix congruence
$\operatorname{Tr}(N_X^2)\equiv\operatorname{Tr}(N_X)\pmod2$
therefore gives $\delta_{X,X^*}=\delta_{X,0}$.
\end{proof}

\begin{lemma}[Conductor bound]\label{arith:conductor}
For either table, the different of $K_Y$ divides $f_Y\mathcal O_{K_Y}$,
and the conductor of $K_Y$ divides $2f_Y$.
In particular, $K_Y$ is unramified at a prime not dividing $f_Y$;
at an odd prime its conductor exponent is at most $v_p(f_Y)$.
The field of $S_{XY}/\operatorname{lcm}(d_X,d_Y)$ lies in
$K_X\cap K_Y$.
\end{lemma}
\begin{proof}
Put $K=K_Y$ and $f=f_Y$. If $a\in\mathfrak D_K^{-1}$, then
\[
 b=\sum_X\operatorname{Tr}_{K/\mathbb Q}
                   (a\overline{\chi_Y(X)})X
   =f\sum_{\sigma\in\operatorname{Gal}(K/\mathbb Q)}
                   \sigma(a)e_{\sigma Y}
\]
belongs to the integral based ring. Its eigenvalue $fa$ is integral;
hence $f\mathfrak D_K^{-1}\subseteq\mathcal O_K$.

Here is the local conductor calculation. After a sufficiently large
finite unramified base extension $F/\mathbb Q_p$, let $E/F$ be the
ramified part of $K$ and let $L$ be the
smallest $p$-power cyclotomic field containing it. Write
$e=[E:F]$ and $d_E=v_E(\mathfrak D_{E/F})$. The preceding divisibility
gives $d_E/e\le v_p(f)$. For odd $p$, write $e=e_0p^s$,
$e_0\mid p-1$. Minimality gives conductor $p^{s+1}$ when $e>1$,
and $L/E$ is tame. The derivative of the cyclotomic polynomial and
transitivity of the different give
\begin{equation}\label{arith:local-different}
 \frac{d_E}{e}
 =s+1-\frac{1-p^{-s}}{p-1}-\frac1e.
\end{equation}
This lies strictly between $s$ and $s+1$ when $s>0$, and between
zero and one for a nontrivial tame extension. Thus the odd
conductor divides $f$.

At two, conductor four has normalized different one. For conductor
$2^m$, $m\ge3$, the subgroup fixing $E$ cannot contain
$1+2^{m-1}$. Every nontrivial subgroup of $\langle5\rangle$
contains that element, so $[L:E]\le2$. In the quadratic case its
generator is $-1$ or $2^{m-1}-1$, and the relative different
has exponent two. Consequently
$d_E/e\ge m-1-2/2^{m-1}>m-2$. Hence $2^{m-1}\mid f$,
which proves the factor-two bound. A conductor cannot have
two-adic valuation exactly one, so $f$ odd also implies
unramifiedness at two. The final assertion follows from symmetry
and rational rescaling of the two character values.
\end{proof}

\paragraph{Residue blocks.}
Fix a place above $p$. Characters are in the same full $p$-block if
they have the same reductions on every basis element. Write $P_p$
for the block containing the dimension character. Over a complete
unramified coefficient ring $\mathcal O$ containing the residue
fields, the block idempotents lift, so
\begin{equation}\label{arith:local-blocks}
 e_B=\sum_{Y\in B}e_Y\in R\otimes\mathcal O,
 \qquad \sum_{Y\in B}\frac{\chi_Y(b)}{f_Y}\in\mathcal O
       \quad(b\in R\otimes\mathcal O).
\end{equation}
By \eqref{arith:galois}, applying a Galois automorphism to
all values of a given row just permutes its columns. Thus equality
of two reductions on every column is independent of the chosen place.
In particular,
\begin{equation}\label{arith:all-places}
       Y\in P_p\quad\Longrightarrow\quad
       \chi_Y(X)\equiv d_X\pmod{\mathfrak p}
       \text{ at every }\mathfrak p\mid p.
\end{equation}
For an algebraic integer in a number field $K$ unramified at $p$,
vanishing at every place above $p$ gives divisibility by $p$ in
$\mathcal O_K$, since $p\mathcal O_K$ is radical. In any field,
a rational integral trace that vanishes at a place above $p$ lies in
$p\mathbb Z$.

\begin{lemma}[Principal support]\label{arith:principal-support}
If $f_Z$ is a power of $p$, every nonzero entry of its row is on
$P_p$. If also $p\nmid d_Z$, the support is exactly $P_p$.
\end{lemma}
\begin{proof}
The Galois-stable sum $a=\sum_{X\in P_p}|\chi_Z(X)|^2$ is an
integer, and symmetry gives $a/f_Z=\tau(e_{P_p}ZZ^*)\in\mathcal O$.
Thus $f_Z\mid a$; the unit contribution and the full row norm give
$1\le a\le f_Z$, hence equality. Conversely, if $X\in P_p$ and
$p\nmid d_Z$, then $\chi_X(Z)\equiv d_Z\ne0$, so symmetry makes
the crossing nonzero.
\end{proof}

\begin{lemma}[Defect one]\label{arith:defect-one}
Work with an integral ordinary modular or positive reduced table
as in \eqref{arith:data}. In particular its integral trace order is
symmetric, its characters are cyclotomic, and its formal codegrees
are rational integers. Let a character in a full $p$-block have
$v_p(f)=1$. Blocks and inertia orbits are understood after an
unramified coefficient extension containing their residue fields.
For odd $p$ the block consists of two tame inertia orbits, both of
defect one. At $p=3$ their sizes are one and two.
At two, either the block is a single ramified quadratic branch,
or an unramified defect-one branch is paired with a branch of
degree $e$, where $e$ is a power of two and the other defect is
$1+v_2(e)$. Thus an even-rank dyadic block containing a
defect-one character has rank two.
No such block contains an unramified character of defect at least two.
\end{lemma}
\begin{proof}
Let $A$ be the local symmetric $\mathcal O$-order of the block and
$M=\prod_i\mathcal O_i$ its normalization after a sufficiently large
unramified extension, so that all branches have the same residue
field $k$. The trace on its $i$th field is $\operatorname{Tr}/f_i$.
Self-duality identifies the conductor with
\[
                (A:M)=M^\#=\prod_i f_i\mathfrak D_i^{-1}.
\]
Its exponent on branch $i$ is $c_i=e_iv_p(f_i)-\delta_i$,
where $\delta_i$ is the different exponent.
If $c_i=0$, its idempotent belongs to the local ring, so it is
the whole block. If $c_i=1$, there must be another branch: otherwise
the maximal ideal of $\mathcal O_i$ and the common coefficient ring
already generate $\mathcal O_i$, making $c_i=0$. Projection to the
other factors gives $A=\mathcal O_i\times_k B$: the conductor contains the
maximal ideal on the first branch, and all residues agree.
Both factors are one-dimensional local rings with residue field $k$,
and each has nonzero maximal ideal. The fiber product is Gorenstein
because $A$ is a symmetric order. The Gorenstein fiber-product criterion
\cite[Corollary~2.7]{NSTV} therefore makes $B$ a discrete valuation
ring. Thus there are exactly two branches, both with conductor exponent one.

For odd $p$, the conductor bound makes a defect-one branch tame,
so $c_i=e_i-(e_i-1)=1$. On the other branch,
$v_p(f)=(d_E+1)/e$. By \eqref{arith:local-different} this is not
an integer for wild ramification, whereas it equals one in the
tame case. At three each tame orbit has size one or two.
The local integrality of $\sum e_i/f_i=\sum e_i d_{X_i}^{\,2}/\Delta$
forces the two sizes to sum to zero modulo three: cancel the common
power of three, after which both dimension squares are one modulo
three. Thus the sizes are one and two.

At two, a ramified defect-one branch satisfies
$2(e_i-1)\le\delta_i\le e_i$, so $e_i=\delta_i=2$ and $c_i=0$.
An unramified one has $c_i=1$. In the resulting two-branch case,
integrality of $1/f_1+e/f_2$ gives $v_2(f_2)=1+v_2(e)$.
The rank is $1+e$, which is even only when $e=1$.
These descriptions also prove the last assertion.
\end{proof}

\begin{corollary}\label{arith:principal-height}
If $v_p(\Delta)\ge2$, every $X\in P_p$ satisfies
\[
 v_p(f_X)\ge2,\qquad
 v_p(d_X)\le\left\lfloor\frac{v_p(\Delta)-2}{2}\right\rfloor.
\]
\end{corollary}
\begin{proof}
The dimension character is unramified with defect $v_p(\Delta)$,
so Lemma~\ref{arith:defect-one} excludes defect one in its block.
A defect-zero idempotent is locally integral and is a singleton
block. Now use $d_X^2f_X=\Delta$.
\end{proof}

\begin{lemma}[Orbit energies]\label{arith:trace-energy}
For a subgroup $H$ of the cyclotomic Galois group, put
$e=|H\cdot Z|$ and $A_H(x)=|H|^{-1}\sum_{\sigma\in H}\sigma(x)$.
Then
\[
 \frac{f_Z}{e}=\sum_X|A_H\chi_Z(X)|^2,
 \qquad 0=\sum_Xd_XA_H\chi_Z(X)\quad(Z\ne0).
\]
In particular, writing $F_X=\mathbb Q(\chi_Z(X))$,
$T_X=\operatorname{Tr}_{F_X/\mathbb Q}\chi_Z(X)$ and summing over
full Galois orbits of $X$, one has
\begin{equation}\label{arith:trace-identities}
 f_Z=\sum_{[X]}[K_Z:F_X][K_X:F_X]T_X^2,
 \qquad 0=\sum_{[X]}d_X[K_X:F_X]T_X.
\end{equation}
If the normalized entries $z_X=S_{XZ}/\operatorname{lcm}(d_X,d_Z)$
are fixed by $H$ on a set $A$ of nonunit labels, then
\begin{equation}\label{arith:orbit-bound}
 e\left(d_Z^2+\sum_{X\in A}\operatorname{lcm}(d_X,d_Z)^2
                         \mathsf M(z_X)\right)\le\Delta,
\end{equation}
where $\mathsf M(z)$ is the mean of $|\sigma z|^2$ over its conjugates.
\end{lemma}
\begin{proof}
The $e$ conjugate rows are orthogonal and have squared norm $f_Z$.
Their average gives the first identity and remains orthogonal to
the dimension row. For the full Galois group its value at $X$ is
$T_X/[F_X:\mathbb Q]$; grouping the $[K_X:\mathbb Q]$ columns
in its orbit gives \eqref{arith:trace-identities}. Retaining the
unit and $A$, multiplying by $d_Z^2$ and taking the full Galois
mean proves \eqref{arith:orbit-bound}.
\end{proof}

\begin{lemma}\label{arith:prime-square-quartic}
If a nonunit row has $f_Z=p^2$, its nonzero absolute traces occur
only at degrees dividing $d_Z/\gcd(p,d_Z)$, and those nonunit
degrees have gcd one. In a simple perfect three-prime table this
is impossible when $v_p(\Delta)=2$ or $4$.
\end{lemma}
\begin{proof}
The norm identity \eqref{arith:trace-identities} reserves a positive
unit term, so a nonzero trace cannot be divisible by $p$.
Since $\Delta=p^2d_Z^2$, the ratio
$d_X/\gcd(d_X,d_Z)$ is either one or $p$; in the latter case
LCM integrality makes the trace divisible by $p$. Hence $d_X\mid d_Z$.
Principal support gives $X\in P_p$ whenever the trace is nonzero.
If $v_p(d_X)=v_p(d_Z)>0$, the identity
$\chi_Z(X)=(d_X/d_Z)\chi_X(Z)$ and \eqref{arith:all-places}
again make its trace divisible by $p$. This proves the divisibility.
The second identity in \eqref{arith:trace-identities} expresses
$-1$ as an integral combination of these nonunit degrees.
When $v_p(\Delta)=2$ or $4$, their common ambient divisor has
only the other two primes. Every nonunit degree contains both,
contradicting gcd one.
\end{proof}

\begin{lemma}\label{arith:small-codegrees}
In a nontrivial ordinary perfect simple integral modular category, every
formal codegree is at least five. For each simple $Z$, one also has
\begin{equation}\label{arith:sts-bound}
 \sqrt{f_Z}\,d_Z\le
                1+\sum_{X\ne0}d_X|\chi_Z(X)|.
\end{equation}
\end{lemma}
\begin{proof}
The $(Z,0)$ entry of $STS=p^+T^{-1}ST^{-1}$, with
$|p^+|=\sqrt\Delta$, gives \eqref{arith:sts-bound}.
Suppose $f_Z\le4$, put $n=d_Z$, and let $H$ be its nonunit
support. Row orthogonality gives
$\sum_{H}d_X\chi_Z(X)=-1$ and
$\sum_H|\chi_Z(X)|^2=f_Z-1$.
LCM integrality and Galois averaging show that $|H|\le f_Z-1$
and every $d_X$ on $H$ divides $n$; a failure of divisibility
alone would spend at least four in the latter sum. The degrees
on $H$ have gcd one, so there are at least two of them.
For two terms both are proper divisors of $n$, and Cauchy's
inequality gives
$\sqrt{f_Z}\,n\le1+(n/2)\sqrt{2(f_Z-1)}$, impossible for
$f_Z=3,4$ and $n\ge2$.
For three terms $f_Z=4$, each entry is a root of unity by the
norm bound. At most one degree is $n$, and the other two cannot
both be $n/2$. Thus $2n\le1+11n/6$, giving $n\le6$.
Lemma~\ref{arith:strict} leaves only $n=6$, whose only divisor
with two prime factors is itself, again contradicting gcd one.
The unit codegree is $\Delta>4$.
\end{proof}

\subsection{Residue blocks and local modular action}

In this section $\mathcal C$ is an ordinary integral modular category.
Write $D=\dim\mathcal C$, $s=S/\sqrt D$, and $T=\operatorname{diag}(\theta_X)$.
A label is \emph{$p$-active} if its $p$-primary twist differs from one.
The aim is to relate activity to the residue blocks of
Section~\ref{sec:arithmetic}. In particular, the vacuum entry of block
orthogonality will give a vanishing sum on the principal block $P_p$.
Reducing this sum at another prime forces nonzero entries at degrees
prime to that prime.

\subsubsection{Normalization and rotations}

The congruence and anomaly theorems give a root of unity $\gamma$ of
two-power order, with $\gamma^3=p^+/\sqrt D$, for which
\[
 \rho\!\begin{pmatrix}0&-1\\1&0\end{pmatrix}=s,
 \qquad
 \rho\!\begin{pmatrix}1&1\\0&1\end{pmatrix}=T/\gamma
\]
factors through $\mathrm{SL}_2(\mathbb Z/M)$, where
$M=\operatorname{ord}(T/\gamma)$. In the integral case one can take
$\gamma^8=1$; in the weakly integral case $\gamma^{16}=1$ suffices
\cite[Theorem~II, Proposition~6.7 and Theorem~6.10]{DLNCongruence}.
The Chinese remainder theorem gives commuting local factors. Write
$W_p,U_p,H_{a,p}$ for their Weyl, upper-unipotent and diagonal torus
operators, and use a subscript $J$ for a product of factors. Thus
$U_J=T_J/\gamma_J$, where the subscript also denotes primary parts.
The matrix identities we shall use are
\begin{equation}\label{loc:bruhat}
 U_J(W_JU_JW_J^{-1})U_J=W_J^{-1},\qquad
 \sigma_a(W_J)=H_{a,J}W_J=W_JH_{a,J}^{-1}.
\end{equation}
The second follows from the Galois covariance of $\rho$.
For integral categories its tori have the particularly simple form
\begin{equation}\label{loc:torus}
 H_a=\frac{\sigma_a(\sqrt D)}{\sqrt D}\,P_a,
\end{equation}
where $(P_a v)_X=v_{\sigma_a X}$ and $\sigma_a$ fixes the dimension
character. Indeed, apply \eqref{arith:galois} to the normalized entries
$s_{XY}=d_Y\chi_Y(X)/\sqrt D$. In square dimension all these tori are unsigned.
Put $M_J=\prod_{p\in J}p^{v_p(M)}$. The entries of $W_J$ belong to
$\mathbb Q(\zeta_{M_J})$, hence are unramified at other primes.
Indeed the congruence theorem places $\rho$ in $\mathbb Q(\zeta_M)$,
and the identities $\sigma_a(t)=t^a$, $\sigma_a(s)=H_as$ give,
first on the generators and then on every group element,
\[
 \sigma_a(\rho(g))=
 \rho\!\left(\operatorname{diag}(a,1)\,g\,
                         \operatorname{diag}(a^{-1},1)\right).
\]
If $a\equiv1\pmod{M_J}$, this conjugation fixes the
Chinese-remainder element defining $W_J$, which is Weyl on $J$ and
identity elsewhere. Thus its entries are fixed by
$\operatorname{Gal}(\mathbb Q(\zeta_M)/\mathbb Q(\zeta_{M_J}))$.
Also $W_J$ is symmetric, $\overline{W_J}=W_J^{-1}$, and
$W_J^2=H_{-1,J}$.
At odd primes the actual twists have no anomaly scalar, so
\begin{equation}\label{loc:twist-square}
                 H_{a,p}T_pH_{a,p}^{-1}=T_p^{a^2}.
\end{equation}
Thus inertia permutes their phases by unit squares.

We give the rotation argument needed for block invariance explicitly.
Let $\mathcal U_{ij}$ correspond to $X_i\boxtimes X_j^*$ under
$\mathcal C\boxtimes\mathcal C^{\mathrm{rev}}\simeq Z(\mathcal C)$.
Its underlying object is $X_i\otimes X_j^*$ and its twist is
$\theta_i/\theta_j$. The additive functional
$L_m^{ij}(V)=\nu_{m,1}^{\mathcal U_{ij}}(V)$ is the trace of rotation
on $\operatorname{Hom}(X_i\otimes X_j^*,V^{\otimes m})$.
Grouping $q$ adjacent tensor factors replaces rotation by its $q$th power;
all power traces belong to $\QQ(\zeta_N)$ for $N=\operatorname{ord}(T)$.
In the center the two factors of the normalized $S$-entry at
$(i,j^*),(A,B)$ are $s_{iA}$ and $\overline{s_{j^*B}}=s_{jB}$.
The twist at $(A,B)$ is $\theta_A/\theta_B$ and its underlying
object is $A\otimes B$. Thus \cite[Proposition~6.1]{NSind} gives
\begin{equation}\label{loc:indicator}
 L_m^{ij}(V)=\sum_{A,B}s_{iA}s_{jB}
                     (\theta_A/\theta_B)^mN_{AB}^{V}.
\end{equation}
For the rotation identities see \cite[Lemmas~2.5 and~2.7,
Propositions~2.8 and~5.7]{NSind}.
For example, on $\mathbb Z/3$ with $\theta_x=\omega^{x^2}$,
$\omega=e^{2\pi i/3}$, our convention gives
$s_{xy}=\omega^{-2xy}/\sqrt3$. Completing the square gives
$sTs=iT^{-1}sT^{-1}$ and $p^+/\sqrt3=i$.
For $(i,j,V)=(1,0,X_1)$, formula~\eqref{loc:indicator} at $m=1$ gives
$\frac13\sum_{a\bmod 3}\omega^{-2a+a^2-(1-a)^2}=\omega^{-1}$,
as required for rotation on $\operatorname{Hom}(X_1,X_1)$.

\begin{lemma}[The two rotation congruences]\label{loc:rotation}
Extend $L_m^{ij}$ linearly to cyclotomic integral coefficients. If
$p\mid m$, then at every place over $p$,
\begin{equation}\label{loc:rotation-mod}
                 L_m^{ij}(y)^p\equiv L_{m/p}^{ij}(y^p).
\end{equation}
For the primitive idempotent $e_Z$ of the fusion algebra and
$A_m=sT^ms^{-1}$ one has
\begin{equation}\label{loc:rotation-entry}
 L_m^{ij}(f_Ze_Z)=f_Z(A_m)_{iZ}\overline{(A_m)_{jZ}},
 \qquad L_1^{ii}(V)=N_{X_iX_i^*}^{V}.
\end{equation}
Finally choose $b\ge\max(1,v_p(N))$ with
$p^b\equiv1\pmod{N_{p'}}$. For a genuine object $V$ and $h\ge2b$,
\begin{equation}\label{loc:rotation-large}
 L_{N_{p'}}^{ii}(V^{p^h})
 \equiv L_{N_{p'}p^b}^{ii}(V^{p^{h-b}})
       \pmod{p^{h-2b+1}\overline{\mathbb Z}_p}.
\end{equation}
\end{lemma}
\begin{proof}
In characteristic $p$, the $p$th power of a rotation trace is the
trace of its $p$th power. Grouping $p$ factors proves
\eqref{loc:rotation-mod} for objects. Additivity and Frobenius in the
commutative residue fusion algebra prove it for integral coefficients.

Substitute
$f_Ze_Z=\sum_V\overline{\chi_Z(V)}[V]$ in
\eqref{loc:indicator}. Multiplicativity separates the two sums into
$f_Z(A_m)_{iZ}(A_{-m})_{jZ}$. Since $s^{-1}=Cs$, with charge
conjugation $C$ commuting with $s,T$, the matrix $A_m$ is symmetric
and unitary. This gives the first identity. At $m=1$ the modular
relation gives $|(A_1)_{iZ}|=|s_{iZ}|$; Verlinde and reciprocity give
the adjoint multiplicity in the second identity.

For \eqref{loc:rotation-large}, both quantities are traces of the
powers $p^h$ and $p^{h-b}$ of the same rotation on
$N_{p'}p^h$ strands. Its central insertion has twist one, so that
rotation has order dividing $N_{p'}p^h$. Its eigenvalue multiplicities
are constant under Galois automorphisms fixing $\mathbb Q(\zeta_N)$,
because all power traces belong to that field. The two powers agree
on the prime-to-$p$ parts and on eigenvalues of $p$-order at most
$p^{h-b}$. On eigenvalues of order $p^t$ with $t>h-b$, the units
congruent to one modulo $p^b$ have free orbits of size
$p^{t-b}\ge p^{h-2b+1}$. Both powers are constant on each such
orbit, since $t\le h$. Each orbit contributes a multiple of that
integer to the difference, proving the congruence.
\end{proof}

\begin{theorem}[Local invariance of residue blocks]\label{loc:block}
The coordinate space of every full $p$-block is invariant under the
$p$-local factor. The same is true after fixing all complementary
primary twists. Consequently
\begin{equation}\label{loc:block-orthogonality}
 \sum_{Z\in B}\overline{s_{Zi}}s_{Zj}=0
           \quad\text{if }\theta_{i,p}\ne\theta_{j,p}.
\end{equation}
In particular, if $E=P_p$ or
$E=P_p\cap\{Y:\theta_{Y,p'}=1\}$, then
\begin{equation}\label{loc:principal-orthogonality}
 \sum_{Y\in E}d_Y\chi_X(Y)=0
                 \quad\text{whenever }\theta_{X,p}\ne1.
\end{equation}
\end{theorem}
\begin{proof}
In the finite residue algebra choose a maximal ideal with residue
field $\mathbb F_{p^a}$ and let $B_0,\ldots,B_{a-1}$ be its geometric
character blocks. For any $c\in\mathbb F_{p^a}^{\times}$, the
Chinese remainder theorem gives an element whose residue is $c$ at
this maximal ideal and zero at the others. Choose a lift $V$ with
nonnegative integral coefficients in the simple basis; it is a genuine
object.

Take $b$ in Lemma~\ref{loc:rotation} also divisible by $a$, and let
$k$ tend to infinity through multiples of $a$. Both $V^{p^k}$ and
$V^{p^{k-b}}$ converge in the split $p$-adic algebra to
$\sum_j[c]^{p^j}e_{B_j}$, where $[c]$ is the Teichm\"uller lift.
Equation~\eqref{loc:rotation-large} makes
$L_{N_{p'}}^{ii}$ and $L_{N_{p'}p^b}^{ii}$ equal on this sum.
The functions $c\mapsto[c]^{p^j}$ are distinct characters of the
finite multiplicative group, so they are linearly independent. The
two functionals are therefore equal on every $e_{B_j}$. As
$N\mid N_{p'}p^b$, \eqref{loc:rotation-entry} gives
\[
 \sum_{Z\in B}|(sT^{N_{p'}}s^{-1})_{iZ}|^2
                    =\mathbf1_{i\in B}.
\]
The rows inside $B$ have their full squared norm on $B$, and those
outside have zero norm there. Thus $sT^{N_{p'}}s^{-1}$ commutes with
the block projector. Since $T^{N_{p'}}=T_p^{N_{p'}}$,
this operator is a scalar multiple of the $N_{p'}$th power of
$sU_ps^{-1}$. The latter has $p$-power order and $p\nmid N_{p'}$,
so its $N_{p'}$th power generates the same cyclic group.
The block projector therefore commutes with the lower unipotent
$sU_ps^{-1}=W_pU_pW_p^{-1}$ and with the diagonal upper $U_p$.
They generate the local factor. This argument also covers the case in
which the dyadic anomaly enlarges the normalized level.
Complementary twist projectors commute with this local factor, so
intersecting with their fibers preserves invariance. Conjugating the
commuting projectors by $s$ gives \eqref{loc:block-orthogonality}; its
vacuum entry gives \eqref{loc:principal-orthogonality}.
\end{proof}

\subsubsection{Support and separated primes}

\begin{lemma}[A coordinate line and its codegree]\label{loc:weyl-line}
If $W_p$ preserves the line of a label $X$, then $p\nmid f_X$.
For ordinary integral modular data, $p\nmid f_X$ holds exactly when
the row of $X$ misses every $p$-active column, and exactly when the
$p$-local group preserves the line of $X$.
Any of these conditions implies $\theta_{X,p}=1$.
\end{lemma}
\begin{proof}
Choose $m$ zero on the $p$-level and one on the other levels, with
$p\mid m$. The Bruhat identity gives
\begin{equation}\label{loc:crt-rotation}
 sT^ms^{-1}
       =\gamma_{p'}^3T_{p'}^{-1}W_{p'}^{-1}T_{p'}^{-1}.
\end{equation}
If $W_p$ preserves the line of $X$, then
$W_{p'}^{-1}=s^{-1}W_p$ gives
$|(W_{p'}^{-1})_{0X}|=1/\sqrt{f_X}$.
Thus $L_m^{00}(f_Xe_X)=1$. But
$(f_Xe_X)^p=f_X^{p-1}(f_Xe_X)$, so
\eqref{loc:rotation-mod} would make this value zero in residue
characteristic $p$ if $p\mid f_X$.

If $p\nmid f_X$, the expression for $e_X$ is locally integral and
$\{X\}$ is a full block. Theorem~\ref{loc:block} and the nonzero
vacuum entry show that $\theta_{X,p}=1$ and the asserted row vanishes.
Conversely, if the row misses all active columns, the lower
unipotent preserves its coordinate line with the same scalar as it
has on the vacuum. The diagonal upper unipotent preserves the line
too. Since they generate the local group and are conjugate there,
their scalars agree; hence $\theta_{X,p}=1$, and the line is
$W_p$-invariant. This proves the equivalences.
\end{proof}

\begin{corollary}[Odd twist orders]\label{loc:odd-twist}
In ordinary integral data, the $p$-part of the twist order at $X$
divides the $p$-part of $f_X$, for every odd prime $p$.
\end{corollary}
\begin{proof}
The case $p\nmid f_X$ is Lemma~\ref{loc:weyl-line}. Otherwise put
$a=v_p(f_X)>0$. The conductor bound
(Lemma~\ref{arith:conductor}) makes local inertia congruent to
one modulo $p^a$ fix $X$. Its square action on twists and the unit
$1+p^a$ imply that the twist order divides $p^a$.
\end{proof}

\begin{proposition}[Separation]\label{loc:separation}
Suppose no actual twist is active at both distinct primes $p,q$.
Then every $p$-active label has $q$-unit codegree and every
$q$-active label has $p$-unit codegree.
If both are dimension primes, their valuations in $D$ are even.
If $p,q$ are odd, no character field ramifies at both. The field
assertion also allows a dyadic prime when $D$ is square and
$\gamma=\pm1$.
\end{proposition}
\begin{proof}
The $q$-local group preserves every $p$-twist fiber. On an active
$p$-fiber its actual $q$-twist is one, so the upper unipotent is
scalar. Its lower conjugate and, by Bruhat, its Weyl are scalar as
well. Lemma~\ref{loc:weyl-line} gives the codegree assertion.
Categorical Cauchy supplies active labels at each dimension prime
\cite[Theorem~3.9]{BNRWFinite}; hence
$v_q(D)=2v_q(d_X)$ for a $p$-active $X$, and $v_p(D)=2v_p(d_Y)$ for a $q$-active $Y$.

At odd primes there is no anomaly scalar, and the just-proved even
valuations make the two local tori unsigned. Write
$A_p=(V^{G_p})^\perp$ for the orthogonal complement of the fixed space
of the local group $G_p$ in the unitary representation $V$. It is
spanned by the $G_p$-translates of the $p$-active coordinate vectors:
the orthogonal complement of that span is invariant and is fixed by
the upper unipotent, hence also by all its conjugates. The $q$-local
group fixes the $p$-active vectors and commutes with $G_p$, so it
fixes $A_p$ pointwise. Consequently
\[
                  (H_{a,p}-I)(H_{b,q}-I)=0.
\]
For commuting unsigned permutations $P,Q$ which both move $X$,
the coefficient of $\mathbf e_X$ in $(P-I)(Q-I)\mathbf e_X$ is
$1+\mathbf1_{PQX=X}$. Applying this observation to the two tori
shows that one of the full inertia groups fixes $X$. This says exactly that
$K_X$ is unramified at that prime. When $\gamma=-1$, tensor the
dyadic factor with the sign character of
$\mathrm{SL}_2(\mathbb F_2)$; its upper unipotent becomes the actual
twist and its tori are unchanged. The same proof covers that case.
\end{proof}

\subsubsection{Integral local columns and isolation}

\begin{lemma}[Local columns]\label{loc:columns}
If $D=n^2$, then $w=W_J\mathbf e_0=W_J^{-1}\mathbf e_0$ is rational
and
\begin{equation}\label{loc:column-integer}
                     \frac{n w_X}{d_X}\in\mathbb Z.
\end{equation}
For any integral modular category and $p\notin J$, put
$a=v_p(f_X)$. Every nonzero entry of the $X$-column of $W_J^{\pm1}$
satisfies, at each place above $p$,
\begin{equation}\label{loc:foreign-denominator}
 v_p((W_J^{\pm1})_{YX})\ge
 \begin{cases}0,&a=0,\\-\lfloor(a-1)/2\rfloor,&a>0.
 \end{cases}
\end{equation}
If all $J$-active labels in an invariant coordinate set have
$v_p(f_X)\le2$, the restricted Weyl matrices are integral at $p$.
\end{lemma}
\begin{proof}
In square dimension the unsigned tori fix the vacuum. Galois
covariance in \eqref{loc:bruhat} makes $w$ rational and gives
$W_J^{-1}\mathbf e_0=W_J\mathbf e_0$. Choose $m$ one on $J$ and
zero on the other levels. The analogue of
\eqref{loc:crt-rotation} and \eqref{loc:rotation-entry} give
$L_m^{00}(f_Xe_X)=f_X|w_X|^2=(nw_X/d_X)^2$.
This is an algebraic integer; its rational square root is an integer.

For the general bound write $W_J^2=\varepsilon P$, where $P$ is
an involutive permutation and $\varepsilon=\pm1$. Symmetry and
unitarity give
$\overline{(W_J^{-1})_{Pi,X}}=\varepsilon(W_J^{-1})_{iX}$.
The mixed identity \eqref{loc:rotation-entry} therefore makes
$f_X(W_J^{-1})_{iX}^2$ integral, up to a root of unity. If $a>0$,
choose the same $m$ divisible by $p$; the rotation congruence makes
its valuation strictly positive. The entries are unramified at $p$,
so their normalized valuations are integers. Thus
$a+2v_p((W_J^{-1})_{iX})\ge0$, strictly if $a>0$, which is
\eqref{loc:foreign-denominator}. Multiplication by $\varepsilon P$
treats $W_J$.

Finally the active columns $u_X=W_J\mathbf e_X$ and their adjoints
are integral at $p$ under the stated bound. On an invariant coordinate
space $E$,
\[
 W_JT_JW_J^{-1}|_E
   =I_E+\sum_{\substack{X\in E\\X\ J\text{-active}}}
                      (\theta_{X,J}-1)u_Xu_X^*.
\]
The right side is integral. The remaining factors in the Bruhat
identity are diagonal units or scalar roots of unity, so that
identity proves integrality of $W_J^{-1}|_E$, and then of its adjoint.
\end{proof}

\begin{proposition}[An isolated prime]\label{loc:isolation}
Assume that $D$ has at least two prime divisors and that a dimension
prime $p$ is isolated: no actual twist is active both at $p$ and
at another prime. Then $D=n^2$ and $\gamma=\pm1$. Write
$n=p^a m$, $(p,m)=1$, and set
$\mathcal W_t=\gamma_t^3W_t$, $b=p^a\mathcal W_p\mathbf e_0$.
Every $p$-active label has degree $p^jm$, $0\le j<a$, and a
$p$-primary character field. Moreover
\begin{equation}\label{loc:isolated-vector}
 b\in\mathbb Z^I,\qquad \|b\|^2=p^{2a},\qquad
 b_X\equiv\frac{\gamma^3d_X}{m}\pmod{p^a\mathbb Z_p}.
\end{equation}
Every label with $v_p(d_X)<a$ belongs to $P_p$ and has
$v_p(b_X)=v_p(d_X)$. Finally
\begin{equation}\label{loc:isolated-vacuum}
 p^a-b_0
 =p^{-a}\sum_{X\ p\text{-active}}
                      p^{2j_X}(1-\theta_{X,p})\in\mathbb Z.
\end{equation}
For odd $p$, $\mathcal W_p=W_p$.
\end{proposition}
\begin{proof}
Separation makes every valuation of $D$ even. An active label has
codegree a positive power of $p$, so its degree and conductor have
the stated forms. There is a nonempty fiber of the odd twists on
which the actual dyadic twist is one: use a $p$-active label when
$p$ is odd, and a label active at an odd prime when $p=2$.
On that fiber the dyadic upper unipotent is $\gamma^{-1}I$, so
Bruhat makes its Weyl $\gamma^3I$. Its square is an unsigned torus,
whence $\gamma^6=1$. The anomaly is a two-power root, giving
$\gamma=\pm1$.

If necessary apply the sign-character correction from
Proposition~\ref{loc:separation}. The corrected local uppers are
the actual twists and their Weyls are $\mathcal W_t$, whose product
is $\gamma^3s$. Every complementary factor fixes the moving space
of the $p$-local group; that group fixes its orthogonal complement.
Consequently
\begin{equation}\label{loc:isolated-split}
       \gamma^3s=\mathcal W_{p'}+(\mathcal W_p-I).
\end{equation}
All $p$-active labels have unit codegree at foreign primes, and all
complementarily active labels have $p$-unit codegree. The last part
of Lemma~\ref{loc:columns} makes $W_p$ integral away from $p$ and
$W_{p'}$ integral at $p$. Rationality and
\eqref{loc:column-integer} now give $b\in\mathbb Z^I$.
Multiply the vacuum column of \eqref{loc:isolated-split} by $p^a$;
its reduction modulo $p^a$ is \eqref{loc:isolated-vector}.
The asserted valuations follow, and nonzero entries of a local
vacuum column lie in $P_p$ by Theorem~\ref{loc:block}.
Unitarity gives its norm. The vacuum entry of the Bruhat identity is
$D(\mathcal W_p)_{00}=\sum_Xd_X^2\theta_{X,p}$, which gives
\eqref{loc:isolated-vacuum} after subtracting the inactive mass.
\end{proof}

\begin{corollary}[No isolated quartic prime]\label{loc:isolated-quartic}
In an ordinary perfect fusion-simple three-prime table, a quartic
dimension prime cannot be isolated if two distinct labels have height
one at that prime.
\end{corollary}
\begin{proof}
Write $n=p^2m$. An active label has degree $m$ or $pm$;
the latter has forbidden codegree $p^2$
(Lemma~\ref{arith:prime-square-quartic}). Suppose all active
degrees are $m$, and let their number be $L>0$. A height-one label
$X$ belongs to $P_p$ by Proposition~\ref{loc:isolation}, so it hits
each active label, whose dimension is a $p$-unit. Their LCM energy
and the positive vacuum contribution give $L<p^2$. Put
$\delta=\sum_{Z\ p\text{-active}}(1-\theta_{Z,p})$.
Equation~\eqref{loc:isolated-vacuum} makes $1-\delta/p^2$ an
algebraic integer. It is a strict convex combination of one and
nonidentity roots at every embedding, so its house is less than one.
It is zero; hence $b_0=p^2-1$. The two height-one labels have
$|b_X|\ge p$, contradicting
\[
              p^4=\|b\|^2\ge(p^2-1)^2+2p^2=p^4+1.
\]
\end{proof}

\begin{lemma}[Duality in odd dimension]\label{loc:odd-duality}
In an integral modular category of odd dimension, every nonunit
simple is non-self-dual. An active $p$-character has a Galois orbit
of size at least $p-1$. More precisely, at twist order $p^a$ each
primitive square class has constant, even phase multiplicity within
any Galois-stable set of labels. The same conclusions apply to
$p$-inertia orbits when the character field is $p$-primary.
\end{lemma}
\begin{proof}
The trace argument in Lemma~\ref{arith:strict} excludes
nonunit self-dual labels. Galois covariance acts on primary twists by
squares, so a primitive square class has constant multiplicity.
Complex conjugation sends a character label to its dual, a distinct
label with the same twist. Each multiplicity is therefore even.
There are $\varphi(p^a)/2$ phases in a primitive square class,
giving at least $\varphi(p^a)$ labels. If the field is $p$-primary,
full Galois conjugation is already supplied by $p$-inertia.
\end{proof}

\subsection{The fermion quotient}\label{loc:fermion-section}

When the symmetric center is generated by a fermion $g$, the positive
quotient identifies $X$ and $gX$. To recover the local action we embed
the category in an integral modular extension. There are three
dimensions to distinguish:
\[
 D_{\mathcal F}=\dim\mathcal F,\qquad
 \Delta=D_{\mathcal F}/2\quad\text{(positive table)},\qquad
 N=2D_{\mathcal F}=4\Delta\quad\text{(extension)}.
\]

\begin{lemma}[An integral minimal extension]\label{loc:extension}
Let $\mathcal F$ be an integral slightly degenerate category with
$\dim\mathcal F>2$, pointed subcategory
$\langle g\rangle=\mathrm{sVec}$, and only the subcategories
$\mathrm{Vec},\langle g\rangle=\mathrm{sVec},\mathcal F$.
It has an integral minimal modular extension $\mathcal M$.
Every minimal extension has pointed subcategory $\mathrm{sVec}$,
adjoint subcategory $\mathcal F$, and no other proper nontrivial
subcategory besides $\mathrm{sVec},\mathcal F$.
In an integral extension every pair of noninvertible simples $X,Y$
satisfies $|\sigma S_{XY}|<d_Xd_Y$ at every embedding $\sigma$.
In particular, $S_{XY}=0$ when their degrees are coprime.
\end{lemma}
\begin{proof}
A minimal modular extension exists by
\cite[Corollary~1.3]{JFR}; it has dimension $2\dim\mathcal F$.
Its fermion-monodromy grading has neutral component $\mathcal F$.
An odd invertible would generate with $g$ a nondegenerate pointed
subcategory of dimension four. Indeed its square lies in
$\{1,g\}$ and its bicharacter with $g$ is $-1$; in the cyclic
case its self-pairing has square $-1$. Splitting off this subcategory
would express $\mathcal F$ as
$\mathrm{sVec}\boxtimes\mathcal N$ with $\dim\mathcal N>1$, a
contradiction. Thus $\mathcal M_{\rm pt}=\mathrm{sVec}$, and
centralizers give $\mathcal M_{\rm ad}=\mathcal F$.

Any proper subcategory crossing the grading has dimension twice
that of its intersection with $\mathcal F$, hence dimension two or
four. Dimension two supplies an odd invertible. At dimension four
it is pointed or Ising: its odd component has dimension two, so
consists either of two invertibles or one object of dimension
$\sqrt2$. An Ising subcategory is nondegenerate and splits off,
again giving the forbidden decomposition of $\mathcal F$.

It remains to obtain an integral extension. Every minimal extension is
weakly integral. For an element $\sigma_a$ of inertia at an odd prime,
the Galois symmetry of its normalized modular data
\cite[Theorem~II]{DLNCongruence} gives
$d_{\sigma_a 0}^2=1$ (apply it to $s_{00}^2=1/N$)
and $\theta_{\sigma_a 0}=1$ (the anomaly normalization $\gamma$
has two-power order). The only invertible with twist one is the unit,
so $\sigma_a0=0$. The common sign in the Galois action on the vacuum
column cancels in $s_{i0}/s_{00}=d_i$. Thus the field of each $d_i$
is unramified at odd primes. Since $d_i^2$ is an integer, the
square-free part of $d_i^2$ is $1$ or $2$.

If the extension has nonintegral simples, its integral subcategory
is exactly $\mathcal F$, by the subcategory lattice just proved.
Every odd simple therefore has degree $a\sqrt2$, $a\in\mathbb Z_{>0}$.
Take the product with an Ising category with fermion $\psi$ and
degree-$\sqrt2$ simple $\sigma$. For the boson $b=g\boxtimes\psi$
form the category of local modules over $\one\oplus b$. This is the
neutral component of the corresponding de-equivariantization
\cite[Proposition~4.56(i),(ii)]{DGNO}. It is modular of dimension $4N/2^2=N$ and contains $\mathcal F$
by $X\mapsto (X\boxtimes\one)\otimes(\one\oplus b)$.
These free modules are local, and induction--restriction shows
that the embedding is fully faithful. The even sector is this
copy of $\mathcal F$. The odd local simples arise from
$Y=X\boxtimes\sigma$ with $X$ odd, so $d_Y=2a$. If $bY\ne Y$,
the free module on $Y$ is simple of module dimension $2a$.
If $bY=Y$, the two choices of the stabilizer character give
two module structures on $Y$, each of module dimension $d_Y/2=a$.
This supplies the integral minimal extension.

For any noninvertible $X$ in an integral extension, the constituents
of $XX^*$ generate $\mathcal F$. Otherwise they lie in
$\mathrm{sVec}$ and $d_X^2\le2$, since each of $\one,g$ occurs
at most once. Equality $|S_{XY}|=d_Xd_Y$ would force the monodromy
of $X,Y$ to be scalar by positivity in balancing. Its scalar on $X$
cancels the inverse scalar on $X^*$, so $Y$ centralizes $XX^*$
and hence $\mathcal F$. But $\mathcal F'=\mathrm{sVec}$.
This proves strictness for noninvertible $X,Y$, also at the conjugate
labels since their dimensions agree. If $(d_X,d_Y)=1$,
$S_{XY}/(d_Xd_Y)$ is an algebraic integer of house less than one,
so it vanishes.
\end{proof}

\begin{lemma}[The two integral orders]\label{loc:fermion-orders}
Let $\mathcal M$ be integral modular with fermion $g$, let
$\mathcal F=\langle g\rangle'$, and put
$\dim\mathcal M=4\Delta$. On duality-compatible even fermion-pair
representatives form the orders
\[
                 R_+=K_0(\mathcal F)/(g-1),\qquad
                 R_-=K_0(\mathcal F)/(g+1).
\]
Both have the orthonormal unit-coefficient trace, and their
multiplication tables agree modulo two. Full modular characters
restrict to $R_+$ on even labels and to $R_-$ on odd labels.
At two, the principal residue character in either order is the
common reduction of the dimension character of $R_+$.
Restriction identifies exactly fermion orbits. Their codegrees are
\begin{equation}\label{loc:restricted-codegrees}
 f_X^\pm=\begin{cases}f_X/2,&gX=X,\\f_X/4,&gX\ne X,
 \end{cases}\qquad f_X=4\Delta/d_X^2.
\end{equation}
At two, a restricted character is principal if and only if its
full character is principal, and corresponding positive and signed
blocks have equal rank. The conductor-node conclusions of
Lemma~\ref{arith:defect-one} apply to both orders.
\end{lemma}
\begin{proof}
Opposite twists make the even fermion action free and prevent it
from identifying a label with its dual. The two quotient structure
constants differ only by the sign of the $gZ$ term, and the
unit-coefficient pairing is $\delta_{XY}$ in either order.
If $\varepsilon_Z$ is the monodromy sign,
$S_{gX,Z}=\varepsilon_ZS_{XZ}$. Orthogonality, followed by division
of the even label set into fermion pairs, gives
\[
 \sum_{Z\ {\rm even}}\chi_X(Z)\overline{\chi_Y(Z)}
       =\frac{f_X}{2}(\delta_{XY}+\delta_{X,gY}),
\]
and hence \eqref{loc:restricted-codegrees} and the assertion on
restriction fibers. Every character of the even subalgebra extends
to the full semisimple algebra, so none are omitted.

In characteristic two the residue on the even subring determines
the entire residue character: for an odd $Y$ its square $Y^2$ is even,
and squaring is injective in the residue field. In particular the
dimension residue lifts exactly the full principal block. Moreover,
the positive and signed orders have identical special fibers.
Corresponding local summands have the same residue dimension,
which equals the rank of each lifted block. Finally the conductor-node argument uses a symmetric
integral trace order, cyclotomic characters and rational codegrees;
it does not require nonnegative multiplication. It therefore
applies to $R_-$ as well as $R_+$.
\end{proof}

\begin{lemma}[Lifting a positive principal pair]\label{loc:principal-lift}
At an odd prime a positive principal fermion pair has a full
principal representative. Both representatives are principal when
all odd degrees vanish in the residue field; otherwise exactly
one is principal.
\end{lemma}
\begin{proof}
For a representative $X$ let $a_Y,b_Y$ be the residues of
$\chi_X(Y),d_Y$. They agree on even objects.
For odd $Y,Z$ their products are even, so
$a_Y^2=b_Y^2$ and $a_Ya_Z=b_Yb_Z$. If every odd $b_Y$ vanishes,
so does every odd $a_Y$. Otherwise choose $b_Z\ne0$; all odd
values satisfy $a_Y=\epsilon b_Y$ with a single sign
$\epsilon=\pm1$. Replacing $X$ by $gX$ reverses exactly the odd
evaluations, proving the assertion.
\end{proof}

\subsubsection{Fermion sectors}\label{loc:fermion-dictionary}
On the full coordinate space put $F\mathbf e_X=\mathbf e_{gX}$
and $G\mathbf e_X=\varepsilon_X\mathbf e_X$, with $\varepsilon_X$
the fermion-monodromy sign. A free pair contributes the unit vectors
$(\mathbf e_X\pm\mathbf e_{gX})/\sqrt2$. The four joint eigenspaces
and their dyadic transitions are as follows.
\[
\begin{array}{c|c|c|c|c}
     & (F,G) & \text{coordinates} & W_2 & U_2 \\ \hline
 A & (+,+)& \text{even pair sums} & A & B\\
 B & (-,+)& \text{even pair differences} & C & A\\
 C & (+,-)& \text{odd pair sums and fixed labels} & B & C\\
 H & (-,-)& \text{odd pair differences} & H & H
\end{array}
\]
Indeed $s$ exchanges $F,G$, while every odd-primary upper
unipotent commutes with both (mates have the same odd twist).
The same follows for its lower conjugate by $s$, hence for the
odd local factors. This gives the $W_2$ column from $s=W_2W_{2'}$.
Balancing with the fermion gives $\theta_{gX}=-\varepsilon_X\theta_X$:
even mates have opposite dyadic spins, odd mates have equal dyadic spins.
Thus $T_2FT_2^{-1}=-FG$ and $T_2G=GT_2$, giving the $U_2$ column. In particular the positive normalized matrix is
$\widehat s=s|_A=(W_2|_A)(W_{2'}|_A)$.
On free pairs the $B,C$ and $H,H$ entries of $s$ are
$2S_{XY}/\sqrt N$. By Theorem~\ref{loc:block} each lifted two-block
is dyadically invariant; its intersection with $A$ is the
corresponding positive block.

\begin{proposition}[Local action in the positive quotient]\label{loc:reduced}
Let $\mathcal F$ satisfy the hypotheses of Lemma~\ref{loc:extension}.
For a full $p$-block $B$ in its positive table,
\[
 \sum_{Z\in B}\overline{\widehat s_{Zi}}\widehat s_{Zj}=0
\]
if the odd $p$-primary twists of $i,j$ differ, or, at $p=2$, if
their squared dyadic twists differ. At odd $p$, the block and its
complementary twist fibers, taken modulo the fermion sign, are
invariant under the full local factor. Every odd dimension prime
$p$ has an active even label. For odd $p$, the order of
$\theta_{X,p}$ divides the $p$-part of $\Delta/d_X^2$, and a
$p$-unit reduced codegree forces the row to vanish on every
$p$-active column.
\end{proposition}
\begin{proof}
Take an integral minimal modular extension. Select all full labels whose
characters on the even subring have the prescribed residue of
$B$. Its coordinate projector $P$ is a sum of ordinary full-block
projectors. It commutes with the fermion permutation, preserves
$A$, and restricts there to $P_B$. The normalized $s$ preserves
$A$ and restricts to $\widehat s$. Odd primary twists preserve
$A$; dyadic twists do so after squaring. Restrict the commuting
projector identities of Theorem~\ref{loc:block}, obtaining the
orthogonality and invariance assertions.

If every even label were $p$-inactive, the odd local upper
unipotent would fix $A$ pointwise, as would its lower conjugate.
The fixed space of the full local factor is invariant under the
modular representation. It contains
$\mathbf e_0+\mathbf e_g$ and hence, after applying $T$,
$\mathbf e_0-\mathbf e_g$. It therefore contains $s\mathbf e_0$,
whose coordinates are all nonzero. All labels would be
$p$-inactive, contradicting categorical Cauchy in the extension.

An even representative is free under $g$, so its full codegree is
four times its reduced codegree. The odd twist-order bound follows from
Corollary~\ref{loc:odd-twist} in the extension.
A reduced $p$-unit character is a singleton block, so the
orthogonality just proved and its nonzero vacuum entry give the
vanishing assertion.
\end{proof}

\subsection{Phase masses and principal witnesses}

Use $\Delta=D$ in ordinary data and $\Delta=\dim\mathcal F/2$ in
positive reduced data. Complementary twists in the latter setting
are understood modulo the fermion sign. All positive quotients below
come from the minimal fermion alternative of Lemma~\ref{loc:minimal}.

\begin{proposition}[An odd valuation forces activity]\label{loc:odd-slices}
Let $p$ and $v_p(\Delta)$ be odd. Let $B$ be a full $p$-block or
the whole label set. If $x,y$ have trivial $p$-twists and $y$ is
fixed by $p$-inertia, then for every complementary twist fiber
\[
 M_j=\sum_{\substack{i\in B\text{ in the fiber}\\
                  \operatorname{ord}(\theta_{i,p})=p^j}}
                s_{xi}\overline{s_{yi}}
       \quad\text{satisfies}\quad M_1=(p-1)M_0.
\]
In reduced data replace $s$ by $\widehat s$. At $x=y=0$, the
squared-dimension mass at order $p$ is $p-1$ times the inactive
mass. In particular the active mass is at least $1-1/p$ of the
fiber's total mass.
\end{proposition}
\begin{proof}
Let $E$ be the projector onto the indicated block and fiber.
Both $E$ and $Q=sEs^{-1}$ commute with the local factor. Write
$u,w,h$ for its upper unipotent, Weyl and torus; at odd $p$ the
upper unipotent is the actual primary twist. For a unit $n$ at
the local level, Bruhat gives the weighted twist moment
\[
 (Qwu^nw^{-1})_{xy}
  =(Qu^{-1/n}wh_{-n}u^{-1/n})_{xy}
  =\left(\frac{-n}{p}\right)(Qw)_{xy}.
\]
Here $wu^nw^{-1}=su^ns^{-1}$ because the other local factors
commute with $u$. Thus the left side is
$\sum_{i\in E}s_{xi}\overline{s_{yi}}\theta_{i,p}^n$,
where $i\in E$ means that the coordinate $i$ belongs to the chosen
block and fiber. On the right the outer unipotents fix $x,y$. The inertia-fixed label $y$
has torus sign $({-n}/p)$ by \eqref{loc:torus}, since
$v_p(\Delta)$ is odd. Averaging over units makes the right side
zero. A root of order $1,p$, or $p^j$ with $j\ge2$ has average
$1,-1/(p-1)$, or zero respectively; this gives the identity.
The positive quotient has the same odd local action and torus
sign by Proposition~\ref{loc:reduced}. At the vacuum its weights
are $d_i^2/\Delta$, giving the mass assertion.
\end{proof}

\begin{corollary}[A defect-one block]\label{loc:defect-one}
Let $p$ be odd and suppose every character in a full $p$-block
has $v_p(f_X)=1$. The block has exactly two inertia orbits: an
inactive orbit of length $a_0$ and degree $d_0$, and an orbit
of primitive $p$-twists of length $a_1$ and degree $d_1$.
Both lengths divide $p-1$, the active phases form one quadratic
class, and
\begin{equation}\label{loc:defect-mass}
                       a_1d_1^2=(p-1)a_0d_0^2.
\end{equation}
\end{corollary}
\begin{proof}
Lemma~\ref{arith:defect-one} gives two tame inertia
orbits. Their dimensions are constant and their lengths divide
$p-1$. Corollary~\ref{loc:odd-twist} bounds their twists by
$p$. Since $v_p(\Delta)=1+2v_p(d_X)$ is odd, Proposition~\ref{loc:odd-slices}
shows that both phase orders occur and gives their mass ratio.
Square Galois covariance makes the active orbit cover exactly
one quadratic class.
\end{proof}

\begin{lemma}[A common visible quadratic class]\label{loc:visible-class}
Let $t,s$ be distinct odd primes with $s\nmid t-1$. In an ordinary
integral table let $\mathcal U$ be a nonempty set of $s$-free
labels of $t$-defect one, and let $\mathcal V$ be a nonempty set
of $t$-defect-one labels in $P_s$. There is a single quadratic class of primitive $t$th roots
such that every column met by any of these rows has $t$-phase
one or in that class.
\end{lemma}
\begin{proof}
The two orbits in each defect-one block have equal $s$-valuation
of their dimensions, by \eqref{loc:defect-mass} and the fact that
their lengths and $t-1$ are $s$-units. Thus the active labels in
every block belonging to $\mathcal U$ have $s$-unit degree and
are hit by each principal row in $\mathcal V$. A row sees only
phases present in its own full $t$-block: $W_t$ preserves the
block by Theorem~\ref{loc:block}, and the complementary Weyl
preserves $t$-twist fibers.
The quadratic classes of each such pair of blocks must agree.
\end{proof}

\begin{lemma}[Activity in a fiber]\label{loc:fiber-activity}
If $p$ and $v_p(\Delta)$ are odd, every nonempty invariant
coordinate space contains a $p$-active label. This applies to
complementary twist fibers in either full fermion sector, and to
positive fibers defined using squared dyadic twists.
\end{lemma}
\begin{proof}
An entirely inactive invariant space would be fixed by the upper
and lower unipotents, hence by the whole local group. But a
nonsquare torus acts on its coordinate basis by minus a permutation,
by \eqref{loc:torus}; it cannot be the identity. In the positive
quotient, fermion-pair coordinates give the same sign. The stated
fibers are invariant because complementary factors and the fermion
grading commute with this odd local group.
\end{proof}

\begin{proposition}[Prime avoidance]\label{loc:prime-avoidance}
Let $E=P_p$ in ordinary or positive reduced data, with $p$ odd
in the reduced case. In ordinary data one may also take
$E=P_p\cap\{Y:\theta_{Y,p'}=1\}$. Assume the unit is its only
invertible label. For a prime $q\ne p$ and a label $X$, suppose
either that $X$ is $p$-active, or that $p,v_p(\Delta)$ are odd
and $q\nmid p-1$. At every place above $q$ there is a nonunit
$Y\in E$ such that $d_Y\chi_X(Y)$ has nonzero residue.
In particular $q\nmid d_Y$ and $S_{XY}\ne0$.
\end{proposition}
\begin{proof}
For active $X$, \eqref{loc:principal-orthogonality} and
Proposition~\ref{loc:reduced} give
$\sum_{Y\in E}d_Y\chi_X(Y)=0$. For inactive $X$ under the
second hypothesis, Proposition~\ref{loc:odd-slices} gives
\[
 \sum_{\substack{Y\in E\\\operatorname{ord}(\theta_{Y,p})=p}}
 d_Y\chi_X(Y)
 =(p-1)\sum_{\substack{Y\in E\\\theta_{Y,p}=1}}d_Y\chi_X(Y).
\]
The vacuum term is one, or $p-1$, a nonzero residue at $q$.
It cannot be the only term to survive. The argument also works with several invertibles in $E$ whenever
their weighted total survives at $q$. In a full fermion extension an
even row takes value one at both $\one$ and $g$, so their unweighted
contributions sum to two whenever both belong to $E$.
\end{proof}

\section{Short cyclotomic words and complementary degrees}\label{sec:crossings}

\subsection{Positive words and their coefficients}\label{sec:words}

For a physical product $X\otimes Y$ of degrees $s,t$, the inverse
double braiding acts on each simple summand $V$ by
$\theta_X\theta_Y/\theta_V$. Taking the positive pivotal trace in
the convention of Section~\ref{sec:arithmetic} gives
\begin{equation}\label{word:balancing}
 \frac{S_{XY}}{\operatorname{lcm}(s,t)}
 =\sum_V\frac{N_{XY}^Vd_V}{\operatorname{lcm}(s,t)}
                     \frac{\theta_X\theta_Y}{\theta_V}.
\end{equation}
The coefficients are positive on every constituent and sum to
$\gcd(s,t)$. In the positive fermion quotient choose even representatives and sum
over their actual product; replacing one representative by its mate
changes the signs of both its twist and the corresponding constituent
twists. We call such a positive rational expression a word and its
coefficient sum its mass. A different expression for the same
cyclotomic integer can have different support. The next two lemmas
identify when information about integral words determines the phases of
the given balancing expression.

\begin{lemma}[Integral words]\label{word:integrality}
For roots supported on at most two odd primes, with any dyadic
factor allowed, every integral value admitting a positive rational
word has an integral word of no greater mass. The same holds for
minimizing a positive linear cost on signed words; every rational
optimum is a convex combination of integral optima.
At a prime-power level, relations are exactly the constant
coefficient vectors on the cosets of the prime-order roots.
\end{lemma}
\begin{proof}
The prime-power assertion follows from
$\Phi_{p^a}(x)=\sum_{j=0}^{p-1}x^{jp^{a-1}}$; omitting one
root in each coset gives an integral basis. It remains true over
a coprime cyclotomic integer ring and after reduction at another
prime.

At square-free two-prime level $uv$, index roots by the edges of
$K_{u,v}$. Their relation lattice consists of the vertex stars,
with one relation between the stars. Reversing the signs on one
vertex class gives the oriented incidence matrix. This matrix is
totally unimodular: expand a square minor at a column with at most
one nonzero entry; if every column has two, its row sum is zero.
After deleting one row, a spanning-tree minor has determinant
$\pm1$. Complementary minors give a unimodular configuration for
the quotient lattice of roots. Higher prime powers split it into
copies of the same configuration; a dyadic factor gives signed
copies, since its roots occur in opposite pairs. This is also the
cyclotomic unimodularity of Beck--Ho\c{s}ten
\cite[Theorem~13]{BH}.

For a unimodular integer configuration, each vertex of $Ax=z$,
$x\ge0$, with $z$ integral is integral: extend its independent
support columns to a basis and use the unit determinant.
Positive costs give a bounded optimal face. Its vertices prove
both assertions, also after adjoining negative columns.
\end{proof}

Two positive expressions can have the same value and different supports:
the sum of the roots in any coset of $\mu_p$ is zero. To recover the support of a balancing expression, we use the integral
minimizers above together with the following contraction. A complementary
phase that would lower the cost is excluded from the given expression.
The remaining coefficient vector is a convex combination of integral
minimizers. If these have a unique coefficient vector, that vector gives
the actual coefficients after equal phases are combined.

\begin{lemma}[Removing complementary phases]\label{word:contraction}
Retain at most two primes in a level $M$ and put the complementary
primes in a coprime level $N$. Prescribe integer residues at the
retained primes. Suppose every signed integral $M$-word with those
residues has cost at least $g$, where a positive letter costs one
and a negative letter costs $c\ge1$. If an integral value with
those residues has a positive rational word of mass at most $g$,
then:
\begin{enumerate}[label=\textup{(\roman*)},leftmargin=*]
\item for $c=1$, all complementary components of its actual phases
      are signs;
\item if $c<\ell-1$, where $\ell$ is the least complementary prime,
      all complementary components are one.
\end{enumerate}
After absorbing the permitted signs, the given coefficient vector
is a convex combination of integral optima on the retained roots.
\end{lemma}
\begin{proof}
Use the tensor product
$\mathbb Z[\mu_{MN}]=\mathbb Z[\mu_M]\otimes\mathbb Z[\mu_N]$.
An integral functional on the second factor fixing one preserves
the stipulated residues: the radical at a retained prime is
generated by its primary uniformizer. Lemma~\ref{word:integrality}
then gives the same cost lower bound for a contracted rational word.

For (i), take a product of prime-power root bases containing $1$.
At each odd prime omit one root per prime-order coset, choosing the
omissions to retain the components of a specified nonsign root as well
as $1$. At two take one representative per opposite pair. Up to sign,
the specified root is a basis element different from $1$. The
coefficient-of-one functional takes only $0,1,-1$ on roots and kills
the specified root.
Its positive occurrence would lower the cost below $g$.
For (ii), on a complementary $u$-power factor choose a functional
which is one at one, minus one at a uniformly chosen nonidentity
$u$th root, and zero elsewhere. Its sum on every prime-order
coset is zero, so it is an integral functional. Tensor these
independent choices. A nontrivial complementary component has
expected signed cost at most $c/(u-1)<1$ for some such $u$;
every other component has expected cost at most one. Its positive
occurrence again gives a contraction of cost less than $g$.
The remaining expression lies in the bounded optimal face, so
Lemma~\ref{word:integrality} applies.
\end{proof}

\begin{lemma}[A large positive residue]\label{word:large-residue}
Suppose an integral cyclotomic value has a positive rational word
of mass at most $g\in\mathbb Z_{>0}$ and residue $g$ at a prime
$r>2g$. Every phase of that word is $r$-primary, and its combined
coefficients are nonnegative integers of sum $g$. The resulting
word is unique at every common $r$-power level.
\end{lemma}
\begin{proof}
A signed integral $r$-word of length at most $g$ has augmentation
in $[-g,g]$, congruent to $g$ modulo $r$. Thus it has exactly
$g$ positive letters. Contraction with both signs of cost one
places the actual expression in the convex hull of such optima.
Two positive words of length $g<r$ with the same value are equal:
their difference is constant on each prime-order coset, and a
nonzero constant requires at least $r$ letters on one side.
\end{proof}

\begin{lemma}[Two positive residues]\label{word:two-residues}
Let $q<r$ be primes. Suppose $g<r$ is a positive integer and an
integral cyclotomic value $z$ has a positive rational word of mass
at most $g$ and residue $g$ at both $q$ and $r$. Then
\begin{equation}\label{word:unsigned}
 z=h+q\sum_{j=1}^m\xi_j,
 \qquad h+qm=g,\quad h,m\ge0,
\end{equation}
with nonidentity $r$-primary roots $\xi_j$, except possibly when
$r=g+1$ and $q\mid g-1$. In that case one may also have
\begin{equation}\label{word:signed}
 z=-\sum_{t=1}^{r-1}\eta_t\zeta_r^t,
\end{equation}
where the $q$-primary roots $\eta_t$ consist of one copy of one
and $(g-1)/q$ complete prime-order cosets.
If the exception is unavailable, every actual word has exactly
the combined multiplicities in \eqref{word:unsigned}.
In particular, $z=g$ for $g<q$, and $|z|=g$ for $g=q$.
\end{lemma}
\begin{proof}
First consider a signed integral word on the $q,r$ alphabet of
length at most $g$ having the prescribed residues. Erasing a primary
component means replacing it by $1$ in each letter. This preserves
reduction at that prime. Erase the $r$-components to obtain a
$q$-primary value $w$. It satisfies $w\equiv g\pmod r$ in a
field unramified at $r$, so $1+(w-g)/r$ is integral. Its conjugates
lie in the disk with center $1-g/r$ and radius $g/r$, which touches
the unit circle only at one. The norm bound gives
\begin{equation}\label{word:erasure}
                         w=g\quad\text{or}\quad w=g-r.
\end{equation}
Now erase the $q$-components and write the signed $r$-coefficients
as $c_t$. Modulo $q$, the coefficients $c_t-g\delta_{t0}$ are
constant on each prime-order coset. If every coset has a vacant
position after marking zero, those constants vanish. The signed
augmentation is then $g$ modulo $q$, excluding $w=g-r$.
Equality $w=g$ saturates the length bound: after absorbing a sign
into a dyadic $q$-component when $q=2$, every letter is positive
with $q$-component one. Every off-identity multiplicity is divisible
by $q$, giving \eqref{word:unsigned}.

Without such a vacancy the length bound forces $g=r-1$ and one
letter at each nonidentity $r$th root. All signs are congruent to
$-g$ modulo $q$. A positive sign would give $q\mid r$; for $q=2$
neither sign has the required parity. Hence $q$ is odd, all signs
are negative, and $q\mid g-1$. Equation~\eqref{word:erasure}
now gives $\sum_t\eta_t=1$. The prime-power coset relations
give one identity plus complete $q$-cycles, as stated.
The same argument excludes every word of length less than $g$.

Contraction of cost one eliminates nonsign complementary phases
and puts the actual expression in the convex hull of integral
optima. Outside the exception these are positive $r$-words of
length $g<r$, with a unique coefficient vector. This proves the
assertion about actual phases and the two final consequences.
\end{proof}

\begin{lemma}[Tame orbit bounds]\label{word:tame-orbit}
If a positive integral word of $g<s$ roots of $s$-power order
has house less than $g$, its tame inertia orbit has size at least
\[
 \frac{s-1}{\max\{a:a\mid s-1,\ a\le g\}}.
\]
\label{word:first-range}
For the words of Lemma~\ref{word:two-residues}, if $q\le g<2q$
and the house is less than $g$, the tame $r$-orbit has size $r-1$,
also in the signed exception.
\end{lemma}
\begin{proof}
The positive coefficient vector of length less than $s$ is unique.
Strictness gives a nonidentity root, on which a tame stabilizer
acts freely. Its order therefore divides $s-1$ and is at most $g$.
In the unsigned first range, strictness forces $m=1$ in
\eqref{word:unsigned}, so the stabilizer fixes a nonidentity root
and is trivial. In the signed case, the nonidentity $r$th roots
are a basis over the $q$-primary integer ring. A stabilizer of
order $a$ makes the coefficients $\eta_t$ constant on its free
orbits. Their sum one is then $a$ times an algebraic integer,
forcing $a=1$.
\end{proof}

\begin{lemma}[One residue in an odd alphabet]\label{word:one-residue}
Suppose all primes in a positive rational word are at least
$\ell\ge3$. If its integral value has mass $g\in\mathbb Z_{>0}$,
residue $g$ at a prime $s>g$, and
$(s-g)(\ell-1)>g$, then every actual phase is $s$-primary and
the combined coefficients are nonnegative integers of sum $g$.
\end{lemma}
\begin{proof}
Choose $\max\{1,g/(s-g)\}<c<\ell-1$. A signed integral
$s$-word has augmentation $g$ modulo $s$. A negative augmentation
costs at least $c(s-g)>g$, and a nonnegative one of cost at most
$g$ requires exactly $g$ positive letters. Contraction removes
the complementary phases; uniqueness of the reduced positive
$s$-word gives the actual coefficients.
\end{proof}

\begin{lemma}[Prime mass and a zero residue]\label{word:prime-mass}
Let $q<r$ be odd primes. An integral value represented by a positive
odd-order word of mass $r$, of house less than $r$, and with
residues $r,0$ at $q,r$ has the form
\[
 z=\eta\left(h+q\sum_{j=1}^m\xi_j\right),
 \qquad h+qm=r,\quad h,m\ge1,
\]
where $\eta$ is $q$-primary and the $\xi_j$ are nonidentity
$r$-primary roots. Every actual word has these combined coefficients.
\end{lemma}
\begin{proof}
Give negative letters cost $3/2$. A signed word of cost at most
$r$ with a negative letter, or of cost less than $r$, has length
at most $r-1$. Erasing its $r$-components gives a $q$-primary
integer $w$ divisible by $r$ and of house less than $r$, hence zero.
Erasing its $q$-components gives coefficients whose deviations
from $r\delta_0$ are constant modulo $q$ on prime-order cosets.
A vacant position after marking zero would make the augmentation
both zero and $r$ modulo $q$. Thus all $r-1$ nonidentity
positions of the identity coset occur once. Their signs are
congruent to $-r$ modulo $q$, hence are all equal since $q$ is odd.
If they were positive, $q\mid r+1$; but the erased value $w=0$
would be a positive vanishing $q$-power word of length $r-1$,
forcing $q\mid r-1$, a contradiction. They are therefore negative,
whose cost $3(r-1)/2$ exceeds $r$. Contraction removes
all complementary phases and expresses the actual word as a
convex combination of unsigned integral optima of length $r$.

For such an optimum, erasing the $r$-components gives a positive
$q$-word $w$ of length $r$. It cannot vanish, because a positive
vanishing $q$-power word has length divisible by $q$. Thus $w/r$
is a nonzero algebraic integer of house at most one. Equality
in the triangle inequality makes all $q$-components the same
root $\eta$.
The remaining positive $r$-word has no full coset, which would
give zero. It also cannot occupy all nonidentity positions of
the identity coset. Placing the remaining letter at zero would
make coefficient constancy give $q\mid r$; repeating a nonidentity
position would give $1\equiv2\pmod q$; and placing it in another
coset would give $1\equiv0\pmod q$.
Hence every coset constant is zero, giving the formula.
Here $h>0$ since $q\nmid r$, and $m>0$ by strictness.
For two optima the ratio of their $q$-components lies in both
primary fields and is rational, hence is one. Their reduced
$r$-coefficient vectors are also equal. Convexity proves the
claim for the original rational expression.
\end{proof}

\subsection{Short character words and their adjoints}
\label{supp:section}

Fix a nonunit character row $Y$. A \emph{full positive word} for
$\chi_Y(X)$ is an expression $\sum_i c_i\xi_i$, with $c_i>0$
rational, roots of unity $\xi_i$, and $\sum_i c_i=d_X$.
Balancing supplies such words. Products, complex conjugation and
Galois conjugation preserve this property, with the corresponding
product object, dual or Galois-conjugate label. Fix a common finite group of roots of unity containing all the actual
balancing phases and the letters of the initial word. All subsequent
words are taken in this group.
An $r$-primary word up to a scalar has all ratios $\xi_i/\xi_j$
of $r$-power order.

We work in the ordinary minimal table or in its positive fermion
counterpart, where a nonunit label is represented by a noninvertible
simple object. Passing to $XX^*$ removes the scalar and replaces the
support by the set of phase ratios. In additive notation a two-position
support gives $\{0,\pm x\}$ for the adjoint. If this set and a Galois
dilate meet only at zero, the next two lemmas force a nonunit
constituent with a one-phase full word, contrary to strictness.

\begin{lemma}[Exposing a primary support]\label{supp:exposure}
Let $r$ be an odd prime, and suppose every other prime in the root alphabet
is at least $\ell\ge2$. A set $D$ of $r$-primary roots is exposed
by a rational linear functional, taking value $1$ exactly on $D$
and less than $1$ on every other root in the alphabet, if
\[
             \ell|D\cap C|<(\ell-1)r
       \qquad\text{for every coset $C$ of $\mu_r$}.
\]
\end{lemma}
\begin{proof}
On a coset containing $a$ marked roots assign value $1$ to them
and $-a/(r-a)$ to the others. These values sum to zero, so they
define a rational linear functional on the primary cyclotomic
field: the relations among the roots are precisely the sums
of the disjoint $\mu_r$-cosets. Tensor this functional with
normalized trace on the complementary cyclotomic field. A
complementary root has trace at least $-1/(\ell-1)$, and its
trace is less than $1$ unless it is $1$. Indeed, the local
traces are $1$, $0$ or $-1/(s-1)$ at a prime $s$.
The only possible new positive value at an unmarked position
is at most $a/((r-a)(\ell-1))<1$.
\end{proof}

\begin{lemma}[A common nonunit constituent]\label{supp:common}
Let $U,V$ be nonunit labels. Suppose full words for
$\chi_Y(UU^*)$ and $\chi_Y(VV^*)$ are supported on exposed
sets $D,F$. Then a common nonunit constituent of $UU^*$ and
$VV^*$ has a full word at $Y$ supported on $D\cap F$.
\end{lemma}
\begin{proof}
The positive based-ring inner product satisfies
\[
           \langle UU^*,VV^*\rangle=\langle UV,UV\rangle.
\]
The unit occurs once in each adjoint. If it were their only
common constituent, $UV$ would be simple. In the fermion
quotient this says that the tensor product of the chosen
representatives is one simple object, one of a fermion pair:
its nonnegative multiplicities have quotient norm one.
In either case its double braiding is scalar, so
$|S_{UV}|=d_Ud_V$, contrary to the strict character bound
of Lemma~\ref{arith:strict}.

Choose a common nonunit constituent $W$. Apply an exposing
functional to the decomposition of $UU^*$ into the actual
positive balancing words of its constituents. Its total
value equals its total mass $d_U^2$. Every phase has value
at most $1$, so every phase in the word for $W$ has value
$1$ and belongs to $D$. Applying the other functional to
the same word puts its phases in $F$ as well.
\end{proof}

\begin{corollary}[Two positions]\label{supp:two-position}
In a full two-position word for a nonunit at a nonunit row, if the
phase ratio has odd prime-power order, that order is $3$ or $5$.
In an odd alphabet only order $3$ is possible. In particular, if
$r\ge7$ is prime, a full word that is $r$-primary up to a scalar
requires at least three positions.
\end{corollary}
\begin{proof}
A single phase contradicts strictness. For two phases let
their ratio $\xi$ have order $r^a$. The adjoint support is
$D=\{1,\xi,\xi^{-1}\}$. If $a\ge2$, these three roots
lie in different $\mu_r$-cosets and are exposed even for
$\ell=2$. If $a=1$, exposure follows from $6<r$, or
from $9<2r$ in an odd alphabet. Whenever $r^a\ge5$,
$D\cap D^2=\{1\}$. In each of the exposed cases a Galois unit with
primary component $2$ and complementary component $1$, followed by
Lemma~\ref{supp:common}, would therefore give a nonunit
full word with one phase, again contradicting strictness.
\end{proof}

\begin{theorem}[Difference descent]\label{supp:descent}
Work in one of the ordinary or positive reduced minimal tables of
Lemma~\ref{loc:minimal}, with actual positive balancing words in a
fixed finite root alphabet. Let $r$ be an odd prime.
Suppose a nonunit label has a full
$r$-primary word, up to a scalar root, on at most $m+1$ positions
at the nonunit row $Y$.
Put
\[
 E=\frac{m(m+1)}2,\qquad H=\frac{r-1}2,\qquad
 u(h)=\min\{k\ge1:k(k+1)\ge h\}.
\]
If $\ell(2E+1)<(\ell-1)r$, where $\ell$ is a lower bound
for the complementary primes, then some $1\le h\le E$ satisfies
\begin{equation}\label{supp:eq:scalar}
                     (H-1)u(h)\le h(h-1).
\end{equation}
\end{theorem}
\begin{proof}
Keep the row $Y$ and the alphabet fixed. Among \emph{all} full words
of nonunit labels at $Y$ whose phase ratios are $r$-primary, allow
any number of positions and choose one minimizing
$h=(|S-S|-1)/2$. Here $S\subseteq\mathbb Z/r^a$ is its support in
additive notation after dividing by one occupied phase. Write $U$
for its label. Strictness gives $h>0$. The initial word has at most
$\binom{m+1}{2}$ opposite pairs of nonzero differences, so $h\le E$.
Lemma~\ref{supp:exposure} exposes the adjoint support $D=S-S$ of
size $2h+1$.

Choose the tame subgroup of $(\mathbb Z/r^a)^\times$, of
order $r-1$, extend its elements by the identity on complementary
roots, and take its quotient by signs. For each class
$b$ put $k_b=|(D\cap bD)\setminus\{0\}|/2$.
Galois covariance supplies $U^b$ at the same row, with
adjoint support $bD$. Lemma~\ref{supp:common} supplies a
nonunit word on
$D\cap bD=\{0,\pm x_1,\ldots,\pm x_{k_b}\}$.
Its nonzero differences up to sign occur among
$x_i,2x_i$ and $x_i\pm x_j$ for $i<j$, at most $k_b(k_b+1)$
in all. This word is eligible for the minimum even when it has
more than $m+1$ positions. Hence $h\le k_b(k_b+1)$ and
$k_b\ge u(h)$.

The tame subgroup acts freely on nonzero opposite pairs:
if $bx=\pm x$, then $b\equiv\pm1\pmod r$, and a tame
element is uniquely determined by its residue. Thus each
ordered pair of the $h$ opposite pairs contributes to at
most one $k_b$. Consequently $\sum_b k_b\le h^2$.
The identity term is $h$, and the other $H-1$ terms are at
least $u(h)$, proving \eqref{supp:eq:scalar}.
\end{proof}

\begin{corollary}[Three positions]\label{supp:three-position}
In an odd root alphabet, a nonunit at a nonunit row cannot
have a full positive $r$-primary word on at most three
positions when $r\ge11$.
\end{corollary}
\begin{proof}
Apply Theorem~\ref{supp:descent} with $\ell=3$ and $E=3$.
Exposure holds since $21<2r$. For $h=1,2,3$ the values
$h(h-1)/u(h)$ are $0,2,3$, whereas $H-1\ge4$.
\end{proof}

\begin{lemma}[Automatic exposure]\label{supp:auto-exposure}
For odd primes $p<q<r$, put
\[
 m=\left\lfloor\frac{\min(r,p^2)}q\right\rfloor,
 \qquad E=m(m+1)/2.
\]
If $m\ge1$, then $m\le p-2$ and
$p(2E+1)<(p-1)r$.
\end{lemma}
\begin{proof}
Since $q\ge p+2$, one has $(p-1)q>p^2$ and $m\le p-2$.
Also $r>mq$. Finally
\[
 \frac{p-1}{p}m(p+2)-\bigl(m(m+1)+1\bigr)
 =m\left(p-m-\frac2p\right)-1>0,
\]
because $m\ge1$ and $p-m\ge2$.
\end{proof}

\subsection{Complementary principal degrees}
\label{red:faces}

In either minimal table, principal residues hold at every place above
the prime.

\begin{lemma}[A residue forces ramification]
\label{red:residue-inertia}
Let $z$ be an algebraic integer of strict house below a positive integer $g$.
If $z\equiv g$ at distinct primes $a,b$ and $b\ge g$, then
$b$ ramifies in $\mathbb Q(z)$. Also, if $z\ne0$ and it vanishes
at a prime $b\ge g$, then $b$ ramifies. Congruences are required
at every place above the indicated prime.
\end{lemma}
\begin{proof}
If $b$ were unramified, $(z-g)/b$ would be integral. In the first
case every conjugate of $1+(z-g)/b$ has modulus less than one,
so this number is zero. Thus $z=g-b$, and reduction at $a$ would
give $a\mid b$. In the second case $z/b$ would be a nonzero
algebraic integer of house less than one. Both are impossible.
\end{proof}

\begin{lemma}[Complementary principal degrees at three primes]\label{red:complementary-degrees}
Let a perfect fusion-simple ordinary or positive reduced table have
three dimension primes $t,u,v$. In the reduced case assume $u$ odd.
The principal block $P_u$ contains $v$-free nonunits. Every such
degree is $t^iu^j$ with $i,j\ge1$ and
$2j+2\le v_u(\Delta)$. Thus at a quartic $u$ both complementary
degrees $ut^i,uv^k$ occur in $P_u$, each of defect two.

When $u<v$ are odd, the $t$-parts of all $v$-free nonunits in $P_u$
and all $u$-free nonunits in $P_v$ are at least $u$.
\end{lemma}
\begin{proof}
Prime avoidance (Proposition~\ref{loc:prime-avoidance}) for a
$u$-active column produces a $v$-free nonunit. It cannot have
prime-power degree, and Corollary~\ref{arith:principal-height}
bounds its $u$-height. At $v_u(\Delta)=4$ interchange $t,v$.
For $u<v$ odd pair opposite labels $X,Y$ of degrees
$t^iu^j,t^kv^l$. The lcm-normalized entry has weight
$g=t^{\min(i,k)}$ by balancing and strict house below $g$.
Principality gives residue $g$ over both $u,v$, on all conjugates
by Galois permutation of labels and symmetry. Lemma~\ref{word:two-residues}
therefore gives $g\ge u$.
\end{proof}

\begin{corollary}[The smallest exponents of a minimal table]
\label{red:small-exponents}\label{red:reduced-height}
A nontrivial minimal table has at least three dimension primes.
If it has exactly three, every ordinary valuation is at least four.
In the reduced case every odd valuation is at least four, and for
category dimension $2^aq^br^c$, $q<r$ odd,
\[
                    2^{\lfloor(a-1)/2\rfloor}\ge q,
                    \qquad a\ge5.
\]
\end{corollary}
\begin{proof}
With at most two primes, absence of prime-power degrees would make
every nonunit degree divisible by both primes, contradicting the
prime-free hit in Lemma~\ref{arith:strict}. This includes a reduced
quotient in which the prime two disappears. At three primes,
Lemma~\ref{red:complementary-degrees} gives the ordinary and odd
valuation bounds. In a reduced table the two-valuation is $a-1$;
square divisibility bounds the common two-part of opposite witnesses
by $2^{\lfloor(a-1)/2\rfloor}$, whereas their common part is at
least $q$.
\end{proof}

\subsection{A word forced by the quartic prime}

\begin{proposition}[The quartic seed]\label{supp:quartic-seed}
Let an ordinary perfect fusion-simple integral modular category
have odd dimension $D$ with precisely three prime divisors $p<q<r$
and $v_r(D)=4$.
Put $m=\lfloor r/q\rfloor$. If $m+1<p$, some nonunit
character at a nonunit row has a full positive $r$-primary word,
up to a scalar root, on at most $m+1$ positions.
\end{proposition}
\begin{proof}
Lemma~\ref{red:complementary-degrees} gives
$X\in P_r$ of degree $p^ir$, and $p$-free nonunits
$B\in P_q$ of degree $q^jr^v$, where $v\in\{1,2\}$.
If some $B$ has $v=1$, its crossing with $X$, divided by
the LCM of the degrees, has mass $r$, strict house less
than $r$, and residues $r,0$ at $q,r$.
Lemma~\ref{word:prime-mass} gives the required word; rescaling
from the crossing to character mass does not change its support.

Suppose all such $B$ have $r$-height two. Their codegrees
are $r$-units, so they miss every $r$-active column.
Prime avoidance (Proposition~\ref{loc:prime-avoidance})
on $P_q$ modulo $p$ therefore shows
that $q$ and $r$ have no jointly active column. The quartic
prime $r$ cannot be isolated by Corollary~\ref{loc:isolated-quartic}:
$P_r$ contains the two distinct complementary height-one degrees
$p^ir$ and $q^jr$. Choose a jointly $p,r$-active
column $J$, and apply avoidance on $P_p$ in the trivial
prime-to-$p$ twist fiber, modulo $q$. It gives a $q$-free
nonunit $V$ meeting $J$, with $p$-primary twist and degree
$p^er$: height two at $r$ would again force an $r$-unit
codegree and prevent the hit.

For any $p$-free $A\in P_r$, of degree $q^ir$, the
normalized $V,A$ crossing has mass $r$, nonzero value
by $V\in P_p$, and zero residue at every $r$-place by
$A\in P_r$. If unramified at $r$, its quotient by $r$
would be a nonzero algebraic integer of house less than
one. Thus this crossing ramifies at $r$. In particular
$K_V$ ramifies there and is unramified at $q$ by
Proposition~\ref{loc:separation}.
For every above $B$, the normalized $V,B$ crossing $z_B$
lies in $K_V\cap K_B$ and hence in a $p$-primary field.
It has mass $r$, strict house less than $r$, and residue
$r$ at both $p$ and $q$.

Average its actual positive word over the complementary
$q,r$ factor. Each negative trace has absolute value at
most $1/(q-1)$; replace a negative primary coefficient
by positive coefficients on the other $p-1$ roots of
its $\mu_p$-coset. The resulting positive rational
$p$-word has mass at most $r$. Subtract its minimum
coefficient on every coset. The reduced coefficients
$c_i$ are nonnegative integers by Lemma~\ref{word:integrality},
with a vacancy in every coset, and their sum $L$ satisfies
\[
 0\le L\le r,\quad L\equiv r\pmod p,\qquad
 c_i-r\mathbf1_{i=0}\ \text{is constant modulo $q$ on each coset}.
\]
Write $r=mq+\rho$, $0<\rho<q$. All nonidentity cosets
have coefficients divisible by $q$. If the identity
coset had only its zero position vacant, its other
coefficients would be $q-\rho$ modulo $q$, giving
\[
 L=(p-1)(q-\rho)+qn,\qquad 0\le n\le m.
\]
Reduction modulo $p$ gives $n\equiv m+1\pmod p$,
impossible because $m+1<p$. The identity coset therefore
has a nonzero vacant position as well. All its constants
vanish, so $L\equiv r\pmod{pq}$. Since $r<pq$, one has
$L=r$.

A nontrivial complementary phase would strictly decrease the mass:
its positive normalized trace is less than $1$, and a negative trace
costs at most $(p-1)/(q-1)<1$ after conversion. Equality $L=r$
therefore forces every letter of the actual word to be $p$-primary.
For each $W$ with $N_{VB}^W>0$, the letter in
\eqref{word:balancing} is $\theta_V\theta_B/\theta_W$.
Its $r$-part is trivial. The twist of $V$ is $p$-primary,
and $B$ is $r$-inactive by its $r$-unit codegree. Thus
$\theta_{W,r}=1$ for every constituent of $VB$. Its fusion
vector $[VB]$ and its Fourier transform
\[
              (s[VB])_T
              =\frac{d_T\chi_T(V)\chi_T(B)}{\sqrt D}
\]
are both supported on $r$-inactive labels; the latter
assertion follows from the $r$-unit codegree of $B$.
Hence both local unipotents fix $[VB]$, and so does
their torus (Theorem~\ref{loc:block}). Since $v_r(D)=4$,
\eqref{loc:torus} makes this the unsigned inertia permutation.
Unsigned Galois covariance and completeness of the characters
make that permutation a based-ring automorphism. As $B$ is
unramified at $r$, this gives
\begin{equation}\label{supp:eq:inertia-product}
                    BV^\sigma=BV\qquad(\sigma\in I_r).
\end{equation}

For each above $A$, its crossing with $V$ ramifies at
$r$, so some $\sigma\in I_r$ changes $\chi_A(V)$.
Evaluating \eqref{supp:eq:inertia-product} at $A$ gives
$S_{AB}=0$ for every $p$-free nonunit $B\in P_q$.
Moreover $A$ is $q$-inactive: separation makes every
$q$-active label have $r$-unit codegree, whereas $r\mid f_A$.
Choose a $q$-active label $T$. Equation~\eqref{loc:block-orthogonality}
gives
\[
                 0=\sum_{W\in P_q}
                       \chi_A(W)\overline{\chi_T(W)}.
\]
All $p$-free nonunit terms vanish. At the other nonunits,
LCM integrality and $p\nmid d_A$ make $\chi_A(W)$
divisible by $p$. The unit contributes $1$, a contradiction.
\end{proof}

\begin{corollary}[The first quartic ratio]\label{supp:first-ratio}
If an ordinary perfect fusion-simple integral modular category has
odd dimension $D$ with precisely three prime divisors $p<q<r$
and $v_r(D)=4$, then
$r>2q$.
\end{corollary}
\begin{proof}
If $r<2q$, then $m=1$ and $m+1<p$. The seed has at most
two positions, contrary to Corollary~\ref{supp:two-position},
since $r\ge7$. Equality $r=2q$ is impossible for primes.
\end{proof}

\section{The quartic position and the odd proper divisors}\label{sec:odd-faces}

\subsection{A quartic exponent below a larger prime}

\begin{lemma}[Propagation from visible phases]\label{red:visible-fiber}
Let $a$ be a prime and $\Phi$ a set of primary phases at a different prime. For each $a$-active column in a specified sector suppose there is a nonzero entry on a witness row, all of whose hits in that sector have phase in $\Phi$. Every fiber in the sector outside $\Phi$ is then $a$-inactive. If its coordinate span is invariant under the full $a$-local group of an ordinary modular table, every label on it has $a$-unit codegree. The same conclusion holds on a positive even fiber for odd $a$ in an integral minimal extension. In the latter assertion the fiber is defined on pairs, using the squared dyadic twist if necessary.
\end{lemma}
\begin{proof}
An active column outside $\Phi$ would contradict its witness. On such a fiber $T_a=I$, so the normalized upper unipotent acts by a scalar, possibly nontrivial at two. By invariance every local conjugate is scalar there; the upper and lower elementary subgroups generate the local group over $\mathbb Z/a^e$. Thus $W_a$ preserves each coordinate line, and Lemma~\ref{loc:weyl-line} gives the ordinary assertion. For a positive even fiber at odd $a$ the action can first be restricted to $A$. If $W_a$ preserves the symmetrized line of $Y,gY$, it also preserves both full lines: it commutes with the dyadic upper unipotent, whose eigenvalues on those two summands are opposite. Apply the ordinary assertion in the extension and divide the full codegree by four. The invariance required here holds for any fiber of a complementary local twist operator, intersected, at odd $a$, with a fermion sector.
\end{proof}

\begin{lemma}[Separation at a quartic prime]\label{red:quartic-separation}
Let an ordinary perfect fusion-simple integral modular table have dimension $t^a u^4v^b$, $u<v$, with three repeated primes. No twist order is divisible by $uv$. The same holds in a reduced minimal table for odd $u,v$.
\end{lemma}
\begin{proof}
Lemma~\ref{red:complementary-degrees} gives $X\in P_u$ of degree $ut^i$. Every nonunit $t$-free $Y\in P_v$ has degree $u^hv^j$, $h=1,2$. If $h=1$, $S_{XY}/\lcm(d_X,d_Y)$ has weight $u$, strict house below $u$ and residue $u$ at both $u,v$ by principality (at $u$ both are zero). This contradicts Lemma~\ref{word:two-residues}, also at $u=2$. Thus all such $Y$ have $u$-unit codegree and vanish on $u$-active columns. For a joint $u,v$-active column, reduction of its orthogonality on $P_v$ modulo $t$ would leave only the vacuum. At odd $u,v$ the same argument uses positive balancing and reduced block orthogonality.
\end{proof}

\begin{theorem}[A quartic below a larger prime]\label{red:ordered-quartic}
Let $p<q<r$ and $0\le a,b,c\le6$. Every braided category of dimension dividing
$p^a q^b r^c$ is WGT when
$\min(a,b)\le4$.
\end{theorem}
\begin{proof}
By descent and Corollary~\ref{red:small-exponents}, only an exact quartic $h$
below a larger prime $k$ remains. An ordinary table with $h$ least
is already impossible: Lemma~\ref{red:quartic-separation} separates $h$
from both larger primes, contrary to Corollary~\ref{loc:isolated-quartic}.
Corollary~\ref{red:reduced-height} disposes of a
reduced table at $h=2$. Thus suppose $h,k$ are odd, and denote the
third prime by $t$. In a reduced table $t=2$, since a category containing $\sVec$ has even dimension.

We work in the ordinary table or its integral minimal extension.
In the latter, separation of $h,k$ in the even sector extends to all
labels. Indeed the joint active space is invariant under the dyadic
Weyl. This operator interchanges the fermion permutation $F$ and the
monodromy sign $G$. A space supported on odd labels contains no
$G$-fixed vector, but contains the $F$-fixed vector $v_X+v_{gX}$
for each of its fermion pairs. Thus $h,k$ are separated by
Lemma~\ref{red:quartic-separation} and no character field ramifies at both
(Proposition~\ref{loc:separation}).

There is a joint $t,h$-active column $J$: in the ordinary case
otherwise $h$ is isolated; in the extension choose an even $h$-active
pair, whose opposite twists allow a dyadically active member.
Choose $A\in P_h$ of degree $hk^j$ and $B,C\in P_k$ respectively of
degrees $h^v k^l,t^w k^m$, with all powers positive, by Lemma~\ref{red:complementary-degrees} and Proposition~\ref{loc:prime-avoidance}. Choose full principal representatives in the extension by Lemma~\ref{loc:principal-lift}. Orthogonality on $P_t$ modulo $k$ now gives a $k$-free
noninvertible hit $V$ on $J$: the invertibles contribute one, or two
in the extension since $J$ is even. Defect-zero vanishing bounds its
$h$-height by one. It hits both $A$ and $C$ by principality, so
coprime vanishing gives $d_V=t^i h$, $i>0$.

The two normalized entries $z_Y=S_{VY}/\lcm(d_V,d_Y)$,
$Y=A,B$, belong to $K_V$ and have strict house less than $h$ by
Lemma~\ref{arith:strict} in the ordinary case and the strict bound
of Lemma~\ref{loc:extension} in the extension.
At $A$ one has $t^i z_A=\chi_A(V)\equiv d_V\equiv0$ at every place above $h$: Galois covariance permutes $V$ through labels of the same degree (Equation~\eqref{arith:all-places}), and $t$ is an $h$-unit. Yet
$z_A\ne0$ since $V\in P_t$ and $t\nmid d_A$. At $B$ the
principalities of $V,B$ give residue $h$ at both $t$ and $k$.
The zero-residue and two-residue assertions of
Lemma~\ref{red:residue-inertia} force ramification of $K_V$
at $h$ and at $k>h$, respectively, contradicting separation.
\end{proof}

\subsection{Odd quintic exponents}

Suppose one of the two smaller primes in an odd minimal table has
exponent five. Complementary principal degrees will give a row whose
character field is unramified at the largest prime. A short crossing
on that row will force ramification. We first construct the required
unramified row.

\begin{lemma}[A missing odd phase forces an invertible]
\label{red:all-prime-phases}
Let a perfect integral modular category have odd dimension $D$, and
let $v_p(D)$ be odd. If its $p$-primary twist level is $p$, then every
$p$th root occurs as a primary twist.
\end{lemma}
\begin{proof}
Put $U=U_p$, $W=W_p$, $w=W_{p'}e_0$, and let $P$ project onto
$\ker(U-I)$. Then $Pw=w$ and $Ww=se_0$. The local torus acts on
$w$ by the Legendre character. Bruhat therefore gives
\[
 PWU^tW^{-1}w=\left(\frac{-t}{p}\right)PWw
 \qquad(t\in\mathbb F_p^\times).
\]
If the nontrivial phase $\zeta_p^k$ were missing, its Fourier
projector would vanish. Conjugating by $W$ and applying $P$ gives
$w=-G Pse_0$, where
$G=\sum_{t\ne0}(\frac{-t}{p})\zeta_p^{-kt}$ has squared modulus
$p$. Thus $|w_j|^2=p d_j^2/D$ on the inactive labels and is zero
elsewhere. Choose an exponent $m$ zero on the $p$-level and one on
the complementary levels. By \eqref{loc:crt-rotation},
$|(sT^ms^{-1})_{0j}|=|w_j|$. Applying \eqref{loc:rotation-entry}
to the idempotent expansion of $V^*$ gives
\[
 L_m^{00}(V^*)=\sum_j\overline{\chi_j(V)}|w_j|^2
        = p\sum_{\theta_{j,p}=1}\frac{\overline{\chi_j(V)}}{f_j}.
\]
These are the coefficients of $b=p\sum_{\theta_{j,p}=1}e_j$.
They are algebraic integers as rotation traces and rational by Galois
stability of the inactive set. Thus $b\in R$, $b=b^*$ and $b^2=pb$.
The unit coefficient is $\sum_j|w_j|^2=1$, so $\|b\|^2=p$;
evaluation at the dimension character gives $\dim b=p$.

The coefficient vector $p sPs^{-1}e_0$ is fixed by the local lower
unipotent and transforms under the torus by the Legendre character.
It has an active coordinate $Y$: otherwise both unipotents would
fix it, contrary to the nontrivial torus character. Twist covariance
puts $\mathbb Q(\zeta_p)^+\subseteq K_Y$. Since $Y\ne Y^*$ in odd
dimension, $K_Y$ is nonreal, and its full orbit has at least $p-1$
labels. Their common nonzero integer coefficient exhausts the
nonvacuum norm $p-1$. Hence
$p=\dim b=1\pm(p-1)d_Y$, impossible for a nonunit in a perfect
category.
\end{proof}

\begin{lemma}[Exceptional mass at an odd quintic]
\label{red:quintic-mass}
Let an integral modular category have dimension $t^5n^2$, where $t$
is an odd prime and $t\nmid n$. Fix a quadratic class of primitive $t$th
roots, and suppose every label outside that class and phase one has
degree $t^i n$. Let $N_0>0$ count the exceptional labels of degree
$n$. Put $c=2$ in odd dimension and $c=1$ otherwise. Then
\[
             2N_0\ge c(t-1)t^2.
\]
If $U\in P_t$ has $t$-height $j$ and an inertia orbit of length $e$
at a prime $s\mid n$, then
\[
                   c(t-1)e t^{2j}<2t^3.
\]
\end{lemma}
\begin{proof}
Let $L_i$ count the degree-$t^in$ exceptional labels with $t$-twist
order greater than $t$, and $M_i$ those at one fixed primitive phase
in the other quadratic class. Square covariance makes the multiplicity
constant within a class, and gives $t\mid L_i$: every deeper square
orbit is a union of complete $\mu_t$-cosets. In odd dimension these
counts are even, since a nonunit and its distinct dual have the same
twist. Put $L=\sum_i t^{2i}L_i$, $M=\sum_i t^{2i}M_i$, and
$A=(t^5-L)/t$. The slice identity of Proposition~\ref{loc:odd-slices}
makes the inactive mass $n^2A$. The masses per primitive phase in the
two quadratic classes are $n^2M$ and $n^2(2A-M)$.

The two sums of the primitive phases over the quadratic classes add to
$-1$ and differ by a quadratic Gauss sum of squared modulus $t$. Hence the primary Gauss sum has squared modulus
$tn^4(A-M)^2$. Its regular-object indicator is
$n^2(A-M)^2/t^4$, a rational algebraic integer by
\eqref{loc:indicator}. Thus $t^3\mid L+tM$. All terms with $i>0$
are already divisible by $t^3$, so
$t^3\mid L_0+tM_0$. Positivity and the stated parity give
$L_0+tM_0\ge ct^3$, and therefore
\[
 N_0=L_0+\frac{t-1}{2}M_0
       \ge\frac{t-1}{2t}(L_0+tM_0)
       \ge\frac c2(t-1)t^2.
\]
Each counted label has codegree $t^5$, so its field is unramified
at $s$. Its degree $n$ is a $t$-unit, and principality gives a
nonzero entry on the $U$-row. These entries are fixed under
$s$-inertia. Write $d_U=t^ju$ with $u\mid n$. Applying the orbit
bound \eqref{arith:orbit-bound} to the $e$ conjugates of $U$ gives $eN_0(n/u)^2<t^5n^2/(t^{2j}u^2)$,
the inequality being strict because the vacuum also contributes.
Combine this with the preceding bound.
\end{proof}

\begin{lemma}[A trace on the trivial twist fiber]\label{red:inactive-trace}
Let an ordinary odd minimal table at three primes have dimension
$D=t^5n^2$, with $t$ prime and $t\nmid n$. Suppose the $t$-primary twist level
is $t$ and one quadratic class of primitive $t$th roots occurs
only on rows of codegree $t^3$ and field $\QQ(\zeta_t)$.
If that class occurs, then for each prime $\ell\mid n$ at least one
$\ell$-free nonunit in $P_t$ is fixed by $\ell$-inertia.
\end{lemma}
\begin{proof}
Fix a label $Z$ in that class and an outer prime $\ell\mid n$.
Its degree is $tn$. On a twist fiber in the chosen class the two
other dimension-prime groups act trivially. They preserve that
fiber, and their upper unipotents fix it pointwise by the
unit-codegree criterion. Their conjugates generate the local groups.
Any dyadic normalization factor acts by scalars, since $D$ and
the actual twist level are odd. Thus the $t$-local Weyl $W$
agrees with the normalized matrix $s$ up to a root of unity on each
such coordinate vector.
The same holds for the torus translates of $Z$.

Put $U=U_t$. For $a\in\mathbb F_t$ set
\[
 M_a=\sum_V d_V\overline{\chi_Z(V)}\theta_{V,t}^{a}
     =\frac{D}{d_Z}(sU^a s^{-1})_{0Z}.
\]
Orthogonality gives $M_0=0$. For $a\ne0$ use the local Bruhat identity
\[
               WU^aW^{-1}=U^{-1/a}WH_{-a,t}U^{-1/a}.
\]
It gives $M_a$ equal to a root of unity times $\sqrt D$. Indeed
the diagonal factors contribute phases, and the vacuum entry of
$W$ at the torus translate has modulus $d_Z/\sqrt D$ and
agrees with that of $s$ up to a root of unity. Since
$\sqrt D=t^2 n\sqrt t$ and $\ell\mid n$, summing over $a$ gives
\[
           \sum_{\theta_{V,t}=1} d_V\overline{\chi_Z(V)}
                      \equiv0\pmod {\ell\overline{\mathbb Z}} .
\]
This set of labels is Galois stable. Its vacuum contribution is
one. Grouping by Galois orbit therefore supplies an $\ell$-free
nonunit $X$ with trivial $t$-twist and nonzero absolute trace of
$x=\chi_Z(X)$. Its row meets $Z$, so $X\in P_t$ by the
prime-power codegree support theorem (Lemma~\ref{arith:principal-support}).

Suppose every such $\ell$-free nonunit is moved by $\ell$-inertia.
In the modular representation space $\mathcal V$, use the moving subspace
$A=(\mathcal V^{G_\ell})^\perp$ of the $\ell$-local group, as in the
proof of Proposition~\ref{loc:separation}. The $\ell$-active
coordinate vectors and their $G_\ell$-translates span $A$.
Since $G_\ell$ commutes with the $t$-local group, $A$ is
$t$-locally invariant. The chosen quadratic class is absent
from $A$: its coordinates are $\ell$-inactive, and $G_\ell$
preserves the $t$-twist fibers.
Let $P$ project onto the $U$-fixed
space. For a nonzero torus eigenvector $v\in PA$ with
eigencharacter $\lambda$, the Fourier projector at a missing
phase $\zeta_t^c$ vanishes on $W^{-1}v$. The same Bruhat identity
gives
\[
 0=\sum_{a\bmod t}\zeta_t^{-ca}P WU^aW^{-1}v
   =v+\left(\sum_{a\ne0}\zeta_t^{-ca}\lambda(-a)\right)PWv .
\]
Consequently $PWv\ne0$, and the parenthesized sum has the same
nonzero value for every $c$ in that class. Replacing $c$ by
$cu^2$ multiplies it by $\lambda(u^2)^{-1}$. Thus every square
torus acts trivially on $PA$. On $P A^\perp$ the $\ell$-group acts
trivially, so on the entire $t$-inactive space
\[
                (H_{u^2,t}-1)(H_{b,\ell}-1)=0 .
\]
Both tori act by unsigned Galois permutations since their
automorphisms fix $\sqrt D$. Choose $b$ moving $X$. The permutation observation in the
proof of Proposition~\ref{loc:separation} now shows that all
square $t$-units fix $X$.
The entry field $F=\QQ(x)\subseteq K_X\cap K_Z$ has degree at
most two, since $K_Z=\QQ(\zeta_t)$. Also $[K_X:F]\ge2$, since
$\ell$-inertia fixes $K_Z$ and moves $X$.

Write $u'$ for the other outer prime. The $\ell$-free nonunit $X$
has degree $t(u')^i$: the principal valuation bound gives
$t$-height at most one and a nonunit prime-power degree is
impossible. The $t$-valuation of $d_Z$ is also one.
Since $X\in P_t$, symmetry gives
$x=(d_X/d_Z)\chi_X(Z)\equiv0$ above $t$. Its nonzero absolute
field trace is an integer multiple of $t$. In the trace-norm
identity (Lemma~\ref{arith:trace-energy}) this orbit contributes
\[
 [K_Z:F][K_X:F]\operatorname{Tr}_{F/\QQ}(x)^2
                         \ \ge\ (t-1)t^2 .
\]
Reducing the ordinary dimension--trace identity at $u'$ supplies
a second, distinct orbit, of degree $t\ell^j$, with a nonzero trace
also divisible by $t$. It contributes at least $t^2$.
Together with the positive vacuum term this exceeds $t^3=f_Z$.
\end{proof}

\begin{proposition}[An odd quintic with two complementary witnesses]\label{red:quintic-fixed-witness}
Let an ordinary odd minimal table have three dimension primes
$t,u,s$, where $s>\max(t,u)$ and $v_t(D)=5$. Suppose the
$s$-free nonunits of $P_u$ and the $u$-free nonunits of $P_s$
all have $t$-height two. Some $s$-free nonunit of $P_t$ is fixed
by $s$-inertia.
\end{proposition}
\begin{proof}
Suppose $s$-inertia moves every $s$-free nonunit of $P_t$. Call the two
sets in the hypotheses the witnesses. They exist by prime avoidance.
Each witness has $t$-defect one. Lemma~\ref{loc:visible-class}
gives a single visible quadratic class for both sets. Each $s$-active or $u$-active label hits
a corresponding witness by principal-block orthogonality modulo
the omitted prime (Proposition~\ref{loc:prime-avoidance}). All their
$t$-phases are therefore $1$ or in the chosen class. There is a
label outside these phases: either a deeper twist occurs or
Lemma~\ref{red:all-prime-phases} gives the opposite primitive class.
Lemma~\ref{red:visible-fiber} applied at $s,u$ gives every label $Z$ of such a fiber pure $t$-power codegree. It follows that $D=t^5n^2$ and the
possible codegrees $t,t^3,t^5$ have degrees $t^2n,tn,n$,
respectively. Each exceptional Galois orbit contains at least
$t-1$ labels: square covariance gives at least $(t-1)/2$
distinct phases, and both duals occur at each phase in odd
dimension.

The row of $Z$ is supported on $P_t$. The dimension--trace
equation reduced at each outer prime gives separate nonunit
orbits of degrees $tu^i$ and $ts^j$ with nonzero entry trace.
Each contributes at least one to the trace norm, while the vacuum
contributes $[K_Z:\QQ]\ge t-1$. Hence $f_Z\ne t$.
If an exceptional label had degree $n$, take an $s$-free nonunit
$U\in P_t$. It has degree $tu^i$ and inertia orbit length $e\ge2$
by assumption. Lemma~\ref{red:quintic-mass}, with $j=1$ and $c=2$, would give $(t-1)e<t$. Every exceptional row therefore has
codegree $t^3$.

To determine its field, put $A=(\mathcal V^{G_s})^\perp$ as in
Lemma~\ref{red:inactive-trace}.
On $A$ the only $t$-phases have order at most $t$. Thus $U_t(t)$
and its conjugates are trivial there. The upper--lower
factorization
\[
 U_t(1)L_t(b)=L_t\bigl(b/(1+b)\bigr)H_{1+b,t}U_t\bigl(1/(1+b)\bigr)
\]
for the lower unipotent $L_t(b)$ at $b=tj$ gives $H_{1+tj,t}=1$ there:
$U_t(1/(1+tj))$ agrees with $U_t(1)$ and the lower parameters
are multiples of $t$. On $A^\perp$ the $s$-group is trivial.
Thus $(H_{1+tj,t}-I)(H_{b,s}-I)=0$ on the whole representation.
Both tori fix $\sqrt D$, so the same permutation observation
shows that every label moved by $s$-inertia is fixed by wild
$t$-inertia. Apply this to
the $s$-free trace-contributing orbit above. Its entry field $F$
is tame at $t$, and its nonzero trace is divisible by $t$ since
$d_Z=tn$. If $K_Z$ were wild at $t$, then $t\mid[K_Z:F]$ and
this orbit alone would cost $t^3$, before the vacuum term.
The codegree makes $K_Z$ a field of $t$-power conductor
(Lemma~\ref{arith:conductor}). Tameness now gives
$K_Z\subseteq\QQ(\zeta_t)$,
with equality by its degree bound. The local Galois unit $1+t$
fixes $Z$, and its square fixes the twist of $Z$. Hence that
twist has $t$-part of order at most $t$. This holds on every
exceptional row, and the other rows already have $t$-phase of
order at most $t$. The global $t$-primary level is thus $t$,
and Lemma~\ref{red:inactive-trace} gives the required contradiction.
\end{proof}

\begin{proposition}[The least quintic face]\label{red:least-quintic}
There is no ordinary odd minimal table with three ordered primes
\(p<q<r\) and \(v_p(D)=5\).
\end{proposition}
\begin{proof}
Opposite \(r\)-free nonunits in \(P_q\) and \(q\)-free nonunits
in \(P_r\) have common \(p\)-part at least \(q\), by
Lemma~\ref{red:complementary-degrees}.  Square divisibility bounds their
\(p\)-heights by two; since \(q>p\), all these heights are two.
Proposition~\ref{red:quintic-fixed-witness}, with \(t=p,u=q,s=r\), gives
an \(r\)-free nonunit \(X\in P_p\) fixed by \(r\)-inertia.
Its degree is \(pq^i\), by the principal height bound.
Pair it with a \(q\)-free witness in \(P_r\).
The normalized entry has mass \(p\), strict house below \(p\),
and residue \(p\) at every place over \(r\).
Since \((r-p)(p-1)>p\), Lemma~\ref{word:one-residue}
gives an unsigned \(r\)-word of length \(p<r-1\).
Lemma~\ref{word:tame-orbit} makes its \(r\)-inertia orbit
nontrivial.  Its field is contained in \(K_X\), a contradiction.
\end{proof}

\begin{lemma}[The middle quintic when its short words are exposed]
\label{red:middle-small}
An ordinary minimal table of dimension \(p^6q^5r^4\), for
odd primes \(p<q<r\), cannot satisfy
\[
 j=\lfloor q/p\rfloor,\qquad 1+j(j+1)<2q/3.
\]
\end{lemma}
\begin{proof}
Every \(p\)-free nonunit \(A\in P_r\) has \(q\)-height two.
If its height were one, pair it with an \(r\)-free member
of \(P_q\), of degree \(qp^i\).  The common mass is \(q\), and the two
positive residues at \(q,r\) force house \(q\) by
Lemma~\ref{word:two-residues}, contradicting strictness.

Suppose an \(r\)-free nonunit \(U\in P_p\) has \(q\)-height one.
Choose \(A\in P_r\) of degree \(q^2r\) as above and
\(B\in P_q\) of degree \(qr^k\).  On the single row \(U\),
the two normalized crossings have common mass \(q\).
The crossing with \(A\) has positive residues \(q\) at \(p,r\),
and is an \(r\)-word on at most \(1+j\) positions.
The exceptional alternative would require \(r=q+1\), impossible
for odd primes.  The crossing with \(B\) has residues \(q,0\)
at \(p,q\), so Lemma~\ref{word:prime-mass} makes it,
up to a scalar, a \(q\)-word on at most \(1+j\) positions.
After taking adjoints the two full character words are supported
on \(r\)-primary and \(q\)-primary difference sets of at most
\(1+j(j+1)\) positions.  Both are exposed by
Lemma~\ref{supp:exposure}; their intersection is \(\{1\}\).
A common nonunit constituent of \(AA^*\) and \(BB^*\)
(Lemma~\ref{supp:common}) would therefore have only phase \(1\)
on the row \(U\), contradicting strictness.
Consequently all \(r\)-free nonunits in \(P_p\) have \(q\)-height
two.

Proposition~\ref{red:quintic-fixed-witness}, now with \(t=q,u=p,s=r\),
gives an \(r\)-free nonunit \(X\in P_q\) fixed by \(r\)-inertia.
Here $d_X=qp^i$. Pair it with a $p$-free witness in $P_r$.
The normalized crossing $z$ has strict house below $q$, belongs to
$K_X$, and is congruent to $q$ at every place above $r$.
Since $K_X$ is unramified at $r$, $(z-q)/r$ is an algebraic integer.
Corollary~\ref{supp:first-ratio} gives $r>2q$, so every conjugate
of this integer has modulus less than $2q/r<1$. Its norm forces
$z=q$, contrary to strictness.
\end{proof}

\begin{proposition}[The middle quintic face]\label{red:middle-quintic}
There is no ordinary odd minimal table of dimension
\(p^6q^5r^4\), for \(p<q<r\).
\end{proposition}
\begin{proof}
Suppose such a table exists. By Corollary~\ref{supp:first-ratio},
$r>2q$. We distinguish the two possible orders of $2q$ and $p^2$.

If $2q<p^2$, put $j=\lfloor q/p\rfloor$. Then
$1\le j\le(p-1)/2$. For $j=1$ we have $1+j(j+1)<2q/3$ since
$q\ge5$. For $j\ge2$ we have $p\ge5$ and
$2p-3(j+1)\ge(p-3)/2\ge1$, so
\[
 3+3j(j+1)<2(pj+1)\le2q.
\]
In either case Lemma~\ref{red:middle-small} excludes the table.

If $2q>p^2$, choose a $p$-active twist fiber. It contains a
$q$-active label since $v_q(D)$ is odd
(Lemma~\ref{loc:fiber-activity}). Prime avoidance in $P_q$ modulo $r$
gives an $r$-free nonunit $X$ meeting this label. Let $Y$ be a
$q$-free nonunit in $P_r$. By Lemma~\ref{red:complementary-degrees},
the common $p$-part of $d_X,d_Y$ is at least $q>p$.
Both $p$-heights are at most three by square divisibility.
If $X$ had height three, $f_X$ would be a $p$-unit and its row would
miss every $p$-active label by Lemma~\ref{loc:weyl-line}.
Thus $X$ has height two, and the normalized $X,Y$ crossing has
weight $p^2$ and residue $p^2$ at $q$ and $r$. Since $r>2q>p^2$,
Lemma~\ref{word:two-residues} gives an $r$-primary word on at most
$1+\lfloor p^2/q\rfloor\le2$ positions: the signed exception
would require $r=p^2+1$, an even integer greater than two.
Rescaling to full character mass contradicts
Corollary~\ref{supp:two-position}.
\end{proof}

\section{The even proper divisors and the exact reduction}\label{sec:even-faces}
\label{red:dyadic-faces}

At dyadic valuation five the ordinary case already illustrates the obstruction.
Complementary witnesses leave a nonempty set of exceptional labels. Its size is at least
four by indicator integrality and at most three by orbit orthogonality.
For a slightly degenerate category one must recover comparable bounds on
fermion orbits, where the signed characters distinguish the two spins of interest.

For an integral minimal modular extension put $N=\dim\mathcal M=4\Delta$.
Use the spaces $A,B,C,H$ of Section~\ref{loc:fermion-dictionary}, and write
\[
 W=W_2,\qquad W_o=W_{2'},\qquad
 U=U_2=\lambda T_2,\quad\lambda=\gamma_2^{-1},\quad\lambda^8=1,
 \qquad K=W|_A.
\]
The positive and signed character orders and their codegrees are those
of Lemma~\ref{loc:fermion-orders}. All spins below are eigenvalues of the
actual twist $T_2$. We use $W=\rho_2(\begin{smallmatrix}0&-1\\1&0\end{smallmatrix})$
and $U=\rho_2(\begin{smallmatrix}1&1\\0&1\end{smallmatrix})$.
In particular $WUW=U^{-1}WU^{-1}$. Balancing for the modular $S$ uses
the inverse phase of each constituent in~\eqref{word:balancing};
complex conjugation fixes the monodromy sign of a fermion.
In ordinary data the same letters denote the corresponding operators
for that category's congruence normalization.

\paragraph{Indicator packets.}
If a complementary Weyl vacuum column satisfies
$|w_x|^2=(c/N)d_x^2\mathbf1_J(x)$ for a positive integer $c$
and a Galois-stable set $J$,
the exponent zero on the omitted local level and one on its
complement gives, by \eqref{loc:indicator},
\[
 b=\sum_V\nu_m(V)[V^*]=c\sum_{x\in J}e_x\in K_0(\mathcal M).
\]
Indeed the coefficients are algebraic integers, and Galois invariance
makes them rational. Thus $b=b^*$, $b^2=cb$, $b_1=1$, and
$\|b\|^2=c$; if $0\in J$, also $\dim b=c$.
For even $m$, shifting one fusion input by $g$ in the indicator
formula gives $\nu_m(gV)=\nu_m(V)$, because
$(\theta_{gX}/\theta_X)^m=(-\varepsilon_X)^m=1$.
Consequently its positive restriction
$b^+=\sum_{[V]\ {\rm even}}\nu_m(V)[V^*]$ satisfies, for every even $x$,
\begin{equation}\label{red:indicator-positive}
       \chi_x^+(b^+)=\frac{\chi_x(b)+\chi_{gx}(b)}4.
\end{equation}
The sum in the numerator cancels odd objects and doubles even
evaluations; the two even representatives have equal coefficients,
giving the second factor two.

\paragraph{The dyadic mass identity.}
Use the ordinary table, or the positive table on $A$ in an integral fermion
extension. Write $s_*$ for its normalized matrix, $v$ for its vacuum, and
$E$ for a full two-block projector (including the choice $E=I$).
Suppose that the full category dimension has odd dyadic valuation. For odd $k$ put
$l=1+4k$ and $M_E(k)=\sum_{Y\in E}|(s_*)_{0Y}|^2\theta_{Y,2}^{-2k}$.
Write $U(a)=U^a$, $L(a)=W U^{-a}W^{-1}$ and $H_l$ for the diagonal torus.
The projector $Q=s_* E s_*^{-1}$ commutes with the even upper and lower
unipotents. For the positive table this follows by restricting the full
block projector and the even powers to $A$. The torus $H_l$ acts on $v$
by $-1$, since its dyadic component is $5\bmod8$ and the odd components
are trivial. Taking the vacuum entry after $Q$ in
$U(2)L(2k)=L(2k/l)H_lU(2/l)$ gives
\begin{equation}\label{red:dyadic-mass-unified}
 \lambda^{2-2k}M_E(k)=-\lambda^{(2-2k)/l}M_E(k/l).
\end{equation}
Indeed $\langle v,Q L(2k)v\rangle=\lambda^{-2k}M_E(k)$.
The scalar powers agree modulo eight, and $k\mapsto k/(1+4k)$ permutes
the odd residues. Averaging equates the degree-square masses at
$\theta_{Y,2}^2=1$ and $\theta_{Y,2}^2=-1$; all higher orders have mean zero.

\subsection{The ordinary dyadic quintic}

\begin{lemma}[A dyadic defect-one row and two visible phases]
\label{red:ef:ordinary-dyadic-phases}
In an integral ordinary modular table with odd $v_2(D)$, a
defect-one row sees only actual dyadic phases $1,\eta$, for one
primitive fourth root $\eta$. Its full two-block has two labels
of equal degree, one of each phase. If an ordinary perfect
modular table has only those two dyadic phases globally, it is
impossible.
\end{lemma}
\begin{proof}
The dyadic conductor-node lemma gives at most two inertia orbits
in the block. Its local Weyl invariance and the complementary
Weyl's preservation of actual dyadic phase show that the block
contains phase one. Equation~\eqref{red:dyadic-mass-unified} equates the mass at orders
one and two together with the mass at order four.
The two occupied inertia orbits must therefore have phases one
and a single primitive fourth root. On this block $U(4),L(4)$
are scalar. The identity
$U(1)L(4)=L(4/5)H_5U(1/5)$ makes $H_5$ scalar, so units
$1\bmod4$ fix the labels. The unramified orbit is a singleton,
the other has length at most two, and equal square mass rules
out length two. Thus the block has two equal-degree labels.
The factorization $s=W_2W_{2'}$ gives the row-support assertion.

If these two phases are global, set $Z=T_2^2$, $F_0=sZs^{-1}$,
$P=(I+Z)/2$. The normalized dyadic level divides eight. Balancing gives
characters with dyadic conductor at most four, so $H_5=-I$.
The identity $U(2)L(2)=L(2/5)H_5U(2/5)$ gives $ZF_0=-F_0Z$.
Write $T_2=aI+cZ$, where $|a|=|c|=1/\sqrt2$.
Since $W=UWUWU$ and $PU=UP=\lambda P$,
\[
 PWP=\lambda^3 P(aI+cF_0)W^2P=\lambda^3 a W^2P .
\]
Here $W^2$ commutes with $Z$ and $PF_0P=0$. The complementary
vacuum vector $w=W_o v_0$ belongs to $P$ and $W^2w$ equals $w$
up to a sign. Since $Ww=sv_0$, the squared entries of $w$ are
$2d_j^2/D$ at the inactive labels and zero elsewhere.
The indicator at a CRT exponent zero at two and one on the odd
level gives an integral self-adjoint virtual element $b$ with
$b^2=2b$, unit coefficient one, dimension two and squared norm
two. Its sole nonunit coefficient is therefore $+1$ on an
invertible simple, contradicting perfectness.
\end{proof}

\begin{lemma}[At least four exceptional labels]
\label{red:ef:dyadic-mass}
Let $D=32n^2$, $n$ odd. Suppose every label whose dyadic phase
is outside $\{1,i\}$ has degree $n$, and at least one exists.
There are at least four such labels.
\end{lemma}
\begin{proof}
We choose $i$ to be the visible primitive fourth root; conjugation
interchanges the two choices. If the maximal dyadic twist order is $2^s\ge8$, let $L$ count
labels at that order. The exponent selecting the corresponding
highest-order involution gives the Gauss sum $n^2(32-2L)$.
The regular-object indicator is $n^2(32-2L)^2/32$.
Since $n$ is odd, indicator integrality implies
$(32-2L)^2/32\in\mathbb Z$, so $4\mid L$.
Otherwise let $u,v$ count the exceptional labels at $-1,-i$.
The Bruhat mass equality puts mass $D/2$ at orders one and two
and mass $D/2$ at order four. The dyadic Gauss sum is therefore
$n^2((16-2u)+i(16-2v))$. Its indicator gives
$8\mid u^2+v^2$, hence $u,v$ are even and their halves have
the same parity. Thus $4\mid u+v$, which is positive.
\end{proof}

\begin{proposition}[Ordinary exclusion at dyadic valuation five]
\label{red:ef:ordinary-five}
There is no ordinary perfect fusion-simple table of dimension
$2^5q^br^c$ for odd primes $q<r$ and arbitrary $b,c$.
\end{proposition}
\begin{proof}
Corollary~\ref{red:small-exponents} requires both companions to occur.
Every pair of opposite principal witnesses at $q,r$ has common
dyadic part at least $q$ and at most four. Thus each witness has
dyadic height two. Its defect is one. By
Lemma~\ref{red:ef:ordinary-dyadic-phases}, all witness rows see
$1$ and a common primitive fourth root: opposite principal
witnesses hit the equal-degree active representatives of each
other's blocks. Prime avoidance makes every $q$- or $r$-active
column hit a witness. Perfectness and the same lemma give an
occupied dyadic phase outside the common pair. On any such
fiber the complementary local groups act scalarly on each line,
so its labels have pure dyadic codegree. Consequently $b,c$ are
even; write $n=q^{b/2}r^{c/2}$.

An exceptional label has degree $2^j n$ and codegree
$2^{5-2j}$ for $j=0,1,2$. Its row is supported on $P_2$.
Reducing the trace-dimension identity modulo $q$ and $r$ gives
two distinct nonunit trace orbits of degrees $2q^h$ and $2r^l$,
by the principal height bound. This already rules out codegree
two. At codegree eight, the exceptional label and both witnesses
have dyadic height one; symmetry and principality make each
nonzero trace even. Their squared traces cost at least $4+4$
besides the vacuum, impossible. Thus every exceptional label
has degree $n$ and codegree $32$.

Take the $r$-free hit $U$ of degree $2q^h$ and a complementary
principal $r$-witness of degree $4r^l$. Their normalized crossing
is a nonscalar unsigned $r$-word of length two, so the
$r$-inertia orbit of $U$ has size $e\ge(r-1)/2\ge2$.
It hits every exceptional degree-$n$ label by two-principality.
These entries are fixed by $r$-inertia, and their LCM scale is
$2n$. Averaging the row norm and retaining the vacuum gives
$4eL<32$, where $L$ counts the exceptional labels, by
\eqref{arith:orbit-bound}. Thus $L\le3$, contrary to Lemma~\ref{red:ef:dyadic-mass}.
\end{proof}

\subsection{Positive and signed dyadic phases}

\begin{lemma}[Positive phases and signed pair transfer]
\label{red:fermion-phase-pair}
Assume $v_2(N)$ is odd.
\begin{enumerate}
\item A positive two-block containing a defect-one character
has squared dyadic twists $+1,-1$, with equal degree-square mass.
Its positive rows vanish on all other squared dyadic phases.
\item If all even actual dyadic twists have order dividing four,
the full dyadic level divides eight. Each block in (1) has two
equal-degree positive labels, both unramified at two. A purely
dyadic full character takes Gaussian-integer values on even objects.
\item Under this bound $U^4=-I$ on $C$. Above a block in
(1), of common positive degree $d$, there are exactly two odd
fermion orbits. Both are free, their spins have ratio $\pm i$,
and their degrees satisfy
\[
       u^2+v^2=2d^2,\qquad v_2(u)=v_2(v)=v_2(d).
\]
If all odd labels have only two dyadic spins of ratio $\pm i$,
$R_+$ has a nontrivial invertible basis element.
\end{enumerate}
\end{lemma}
\begin{proof}
Let $E$ be the block projector on $A$. It commutes with $K$.
Since $\widehat s^{-1}\mathbb C^E=W_o^{-1}\mathbb C^E$
has nonzero vacuum projection and $W_o$ preserves the squared
dyadic twist, the block contains squared twist one.
Equation~\eqref{red:dyadic-mass-unified} forces positive mass at squared
phase $-1$. By Lemma~\ref{arith:defect-one} there are at most
two inertia orbits. The two signs already occupy distinct orbits,
so no others occur. The equality of subspaces above then gives
the row-support assertion.

Assume now the even-twist bound. Balancing puts even-even entries
in $\mathbb Q(i,\zeta_m)$ for odd $m$. A dyadic Galois unit
$l\equiv1\pmod4$ fixes every even restricted character, hence
every even fermion pair. Its square fixes the actual twists, and
opposite twists distinguish the two representatives. Symmetry
then fixes the values of every character on even objects, so
every odd fermion orbit is fixed as a set. Taking $l=5$ and using
the square action on twists gives $\theta_{Y,2}^{24}=1$ for odd
$Y$. Thus the full level divides eight, and purely dyadic values
on even objects lie in $\mathbb Z[i]$. Positive inertia orbits
have length at most two. One branch in the defect-one block is
unramified; an orbit of length two on the other branch would make
the equal-mass identity an equality of an integer square with
twice an integer square. Hence both branches have length one
and equal degrees.

Put $Z=T_2^2|_A$ and $F_0=KZK^{-1}$. The dyadic torus $H_5$ acts as
$-I$ on $A$, since its permutation fixes each positive label and its sign
is negative. Modulo eight $2/5=2$, so $U(2)L(2)=L(2)H_5U(2)$ gives
$ZF_0=-F_0Z$. We have $Z^2=F_0^2=I$.
For completeness, the transport takes place within a \emph{full} residue
block. Modulo two the residue on the even subring determines every odd
value, since the square of an odd object is even and squaring in the
residue field is injective. Thus the positive and signed blocks are
intersections of the same full residue block with $A$ and $C$.
Its projector commutes with $W$ and $U$ by Theorem~\ref{loc:block}.
The unitary map $P_0=WU:A\to C$ therefore restricts to an isomorphism
on these intersections. Each positive block in (1) has rank two, so its
$C$-space has precisely two coordinates, one per odd fermion orbit,
including a single coordinate for a fixed orbit. This gives exactly two
odd orbits before making any assertion about their types or spins.
The map $P_0$ transports their spins to $A$.
Centrality of $W^2$ and $WUW=U^{-1}WU^{-1}$ give
\begin{align*}
 B_0=P_0^{-1}UP_0
 &=W^{-2}U^{-2}W|_A=\lambda^{-2} ZK^{-1},\\
 B_0^2&=\lambda^{-4} ZF_0 K^{-2},\qquad B_0^4=-I .
\end{align*}
Here $K^2$ commutes with $Z,F_0$ and $K^4=I$. On a positive block as
in (1) its inertia permutation fixes both labels, so $K^2$ is scalar.
Conjugation by $Z$ negates $B_0^2$ on this rank-two space. Its two
eigenvalues are $i,-i$, hence the two actual odd spins have ratio $\pm i$.

The two odd characters have equal codegree valuations. If inertia
interchanges them this is immediate. Otherwise, after passing to the strict henselization, their signed
local order is $\{(x,y)\in\mathcal O^2:x\equiv y\pmod{2^b}\}$;
its conductor has length $b$ on both branches, so trace duality
gives equal codegree valuations. The restriction-codegree formula
in Lemma~\ref{loc:fermion-orders} reads, with $a=v_2(N)$,
\[
 v_2(f_X^-)=
 \begin{cases}a-1-2v_2(d_X),&gX=X,\\
               a-2-2v_2(d_X),&gX\ne X.
 \end{cases}
\]
As $a$ is odd, equal defects cannot mix fixed and free types; within
one type they give equal degree heights. Full block orthogonality of
$1,g$ equates the odd and even squared-dimension masses, the latter
being $4d^2$. Two fixed labels would give $u^2+v^2=4d^2$;
equal heights make the left valuation odd, while the right valuation
is even. Thus both are free and $2u^2+2v^2=4d^2$.
Taking valuations in $u^2+v^2=2d^2$ gives the claimed height equality.

Finally suppose only two odd spins occur. On $C$ one has
$U=aI+bU^2$, with $|a|=|b|=1/\sqrt2$. Transport this relation to $A$ and multiply
$B_0=aI+bB_0^2$ on the left by $Z$ and on the right by $K^2$:
\[
 K=\lambda^{2}a ZK^2+\lambda^{-2} b F_0 .
\]
For $P=(1+Z)/2$, anticommutation gives
$PKP=\lambda^2a PK^2P$. The complementary vacuum vector
$w=W_ov_{[1]}$ satisfies $Pw=w$, and $K^2w$ is a scalar multiple
of $w$. Since $Kw=\widehat s v_{[1]}$, it follows that
\[
 |w_{[j]}|^2=(2d_j^2/\Delta)\mathbf1_{\theta_{j,2}^2=1}.
\]
In the full vacuum column exactly one even representative of each
such pair occurs, the one of actual dyadic twist one. At a CRT
exponent zero on the dyadic level and one on the odd level, the
indicator formula therefore says that the coefficients of
\[
               1+J,\qquad J=\sum_{[j]}\theta_{j,2}^2e_{[j]},
\]
are integers by \eqref{red:indicator-positive}. Now $J=J^*$, $J^2=1$, and
$\operatorname{FPdim}(J)=1$. Its based norm is one, so it is a
positive invertible basis element. Both signs occur, so it is
not the unit.
\end{proof}

\paragraph{A small dyadic sum.}
Let $E$ be a nonempty Galois-stable set of at most three indices.
Suppose $\alpha_{ij}$, $i,j\in E$, are integral dyadic units in a
two-power cyclotomic field, with symmetric squared moduli. Assume
Galois acts by permuting either index, up to signs:
$\sigma(\alpha_{ij})=\pm\alpha_{\sigma i,j}=\pm\alpha_{i,\sigma j}$.
Then $\sum_{i,j\in E}|\alpha_{ij}|^2$ cannot belong to $8\mathbb Z$.
Indeed reduction at the unique dyadic place gives $|E|^2\bmod2$,
so divisibility by eight would force $|E|=2$. If both indices are fixed,
each $\alpha_{ij}$ spans a rational Galois character line of order
at most two. The quadratic subfields of a dyadic cyclotomic field
are among $\mathbb Q(i),\mathbb Q(\sqrt2),\mathbb Q(\sqrt{-2})$.
Integral elements on the lines $\mathbb Q\sqrt{\pm2}$ are dyadic
nonunits. Thus each squared modulus is an odd integer square and the
sum is four modulo eight. If the indices are exchanged, covariance
gives the matrix of squared moduli
$\bigl(\begin{smallmatrix}y&\sigma y\\\sigma y&y\end{smallmatrix}\bigr)$.
The real dyadic integer $y$ has degree at most two, hence
$y=a+b\sqrt2$ with $a,b\in\mathbb Z$ and $a$ odd. If $b\ne0$ the
interchange acts by conjugation. In either case the sum is $4a$,
again four modulo eight.

\subsection{The fermion sextic exclusion}

\begin{theorem}[No minimal fermion category at dyadic valuation six]
\label{red:fermion-six}
There is no integral slightly degenerate category $\mathcal F$
with pointed subcategory $\mathrm{sVec}$, no other proper
nontrivial fusion subcategory, and dimension $2^6q^br^c$, where
$q<r$ are distinct odd primes. The exponents $b,c$ are arbitrary.
\end{theorem}
\begin{proof}
Put $\Delta=\dim\mathcal F/2$. Its positive quotient is perfect
and fusion-simple. Corollary~\ref{red:small-exponents} excludes a
missing companion prime. Opposite principal witnesses at $q,r$ have
degrees $2^i q^j,2^h r^l$. Their common dyadic weight is at least
$q$ by the two positive residues and at most four because
$v_2(\Delta)=5$. Hence $q=3$ and the common weight is four.
If $r>5$, Lemma~\ref{word:two-residues} gives $1+3\xi$, with $\xi$
a nontrivial $r$-power root; Corollary~\ref{supp:two-position} excludes it. Consequently $r=5$. Every complementary witness
has degree $4\cdot3^i$ or $4\cdot5^j$, and positive dyadic
defect one.

Choose the integral extension of Lemma~\ref{loc:extension}
and use its positive and signed pair coordinates.
Its dimension is $N=4\Delta$, of dyadic valuation seven.
We first bound the even twist orders. The signed characters then force the
odd orbits of exceptional spin to be free and of the same odd degree.
An orbit norm leaves at most three such orbits; their trace gives the
final parity contradiction.

\paragraph{The even twists have order at most four.}
By Lemma~\ref{red:fermion-phase-pair}(1), complementary witness
rows see only squared dyadic phases $+1,-1$. An even pair at
any other squared phase cannot be active at three or five:
principal prime avoidance modulo the opposite prime would force
a complementary hit. On such an extra phase fiber, the odd local
groups preserve the fiber and their upper unipotents act as the
identity. The conjugate unipotents generate, so those groups fix
the full coordinate lines. The complementary-codegree principle
therefore makes the positive codegree of each extra pair purely
dyadic. Suppose that at least one such pair exists. Its degree then
forces $b,c$ even; write
$n=3^{b/2}5^{c/2}$, so $\Delta=32n^2$.

For an extra pair $V$, $K v_V=\widehat s v_V$ has nonzero
vacuum coefficient, so $V\in P_2$. Its possible positive
codegrees are $32,8,2$. Only $32$ can occur. Indeed its character
vanishes modulo two on every supported even degree. Apply
dimension orthogonality after Galois trace and reduce once modulo
three and once modulo five. The vacuum contribution is one, so
there are two distinct prime-omitting nonunit orbits having
nonzero trace; both have even degree, since a nonunit degree
cannot be a prime power. Each trace is a nonzero even integer.
The Galois-averaged row norm therefore gives
\[
                       f_V\ge [K_V:\mathbb Q]+4+4>8.
\]
Thus $d_V=n$ for every extra pair.

Dimension orthogonality modulo five gives a five-free nonunit hit
$X$ on such a row. Pure-dyadic codegree puts it in $P_2$, whose
positive height bound gives $d_X=2\cdot3^h$. Crossing it with
a complementary $4\cdot5^j$ witness gives a nonscalar unsigned
five-word of length two. Its tame five-inertia orbit has length
$t\ge2$. Every extra pair is hit by this orbit, and its field is
five-unramified. If $L$ is the number of extra pairs, averaging
the $X$-row over this orbit and retaining the vacuum gives
\[
                              4tL<32.
\]
Hence $L\le3$.

Let $H_0$ be those whose squared dyadic
phase has maximal order $2^m$, $m\ge2$. On $A$ the corresponding
powered twist is $J=1-2P_{H_0}$. The group commutator
$JKJ^{-1}K^{-1}$ has algebraic-integer trace. Its scalar
normalization cancels, and expansion gives
\[
 \operatorname{Tr}(JKJ^{-1}K^{-1})
 =\dim A-4|H_0|+\frac18
                   \sum_{V,Y\in H_0}|\chi_V(Y)|^2.
\]
The last sum is rational by Galois invariance, so it belongs to
$8\mathbb Z$. The entries $\chi_V(Y)$ are dyadic units by principality.
Their squared moduli are symmetric since all these degrees are $n$.
The small dyadic sum above contradicts divisibility by eight, as
$1\le|H_0|\le3$. This proves the even-twist bound.

Lemma~\ref{red:fermion-phase-pair} now gives dyadic level at most
eight, Gaussian values on even objects for purely dyadic
characters, and two equal-degree positive labels in each
complementary block.

\paragraph{The odd spins leave a nonempty exceptional packet.}
Above each complementary block are two free odd orbits of spins
with ratio $\pm i$, and degrees $u,v$ satisfying
$u^2+v^2=2d^2$ and dyadic height two. If $d=4\cdot3^i$,
successive reduction modulo three gives $u=v=d$. If
$d=4\cdot5^j$, reduction modulo three gives $3\nmid uv$.
These odd mates provide a label of unit degree at either odd
prime for the opposite witness.

By Lemma~\ref{loc:principal-lift} each complementary witness has a
full principal lift at its odd prime. It hits both odd mates of every
opposite block because their degrees are units there.
The dyadic Weyl preserves its full
two-block, while the complementary Weyl preserves actual dyadic
spin. Therefore the same pair of spins, say $\delta,i\delta$,
occurs above every complementary block.

The full principal block at an odd prime contains only even
labels, by evaluation at $g$, and excludes $g$, by evaluation at
an odd label of prime-free degree. Principal prime avoidance
therefore makes every odd label active at three or five hit a
complementary witness. All such labels have spin $\delta$ or
$i\delta$.

There must be another odd spin, since otherwise the last assertion
of Lemma~\ref{red:fermion-phase-pair} contradicts perfectness of
$R_+$. Since $U^4$ is scalar on $C$, the other spins are
$-\delta,-i\delta$. Denote their odd fermion orbits by $E$.
The odd local groups fix their full lines by the same fiber
argument used above. Thus their full codegrees are purely
dyadic. Since $E$ is nonempty, $b,c$ are even. Put
$n=3^{b/2}5^{c/2}$. Each $V\in E$ is full two-principal because
$Wv_V=sv_V$ has nonzero vacuum entry. The full principal height
bound gives $d_V=2^j n$ for $j=0,1,2$.

\paragraph{A signed norm forces every exceptional orbit to be free and odd.}
The positive principal two-block has at least eight odd-degree
labels. Indeed $Z=T_2^2|_A$ and $KZK^{-1}$ anticommute, so its
rank is even. Its degree-square mass is divisible by $32$.
The positive principal height bound permits only odd degrees and
twice-odd degrees. If their counts are $k_0,k_1$, then
$4\mid k_0$, $2\mid k_1$, and
$k_0+4k_1\equiv0\pmod8$. The vacuum therefore gives $k_0\ge8$.
The signed character of $V\in E$ is nonzero on all these
odd-degree labels by principality; its Gaussian values contribute
at least eight to its norm.

Compare this signed character with an odd mate of a
$4\cdot5^j$ witness. Their restricted rows are orthogonal,
because their spins differ. Modulo three the vacuum survives,
and the three-bearing terms vanish on the three-free mate by
LCM integrality. There is therefore a three-free nonunit hit
$X_3$ on the $V$-row. Similarly there is a five-free hit $X_5$.
Pure-codegree support and the positive principal height bound give
\[
                     d_{X_3}=2\cdot5^l,
                     \qquad d_{X_5}=2\cdot3^h.
\]
Their nonzero Gaussian values vanish above two. An orbit of such
values costs at least four per label unless the squared modulus
is two. In the latter case $\chi_V(X)^2=\pm2i$, and
\[
                 \chi_X(V^2)=(d_V/d_X)^2\chi_V(X)^2.
\]
Since $V^2$ is even, $i$ lies in the positive character field of
$[X]$. Five-inertia fixes $i$, whereas a full Galois orbit of a character
whose field contains $i$ has a conjugate with $i$ replaced by $-i$.
Consequently its size is at least twice its five-inertia orbit size,
with the same squared moduli. Thus a norm-two value costs at least
four per five-inertia label, just as a value of norm at least four.
The five-inertia orbit of
$[X_5]$ has size at least two by its length-two crossing, while
the trivial lower bound for $[X_3]$ is one. These two disjoint
prime-omitting supports therefore cost at least eight and four.
The signed codegree is at least
\[
                               8+8+4=20.
\]
By Lemma~\ref{loc:fermion-orders} it is $32/4^j$ for a free
orbit and $64/4^j$ for a fixed one. Thus $j=0$.
A fermion-fixed odd label would have even $d_V^2$, because
$VV^*$ is $g$-invariant and all even constituents occur in free
pairs. Therefore every orbit in $E$ is free and has degree $n$.

\paragraph{The packet has fewer than four orbits.}
Choose $X=X_5$. Every $V\in E$ hits $X$ by two-principality;
LCM integrality makes $\chi_V(X)$ a nonzero multiple of two in
$\mathbb Z[i]$. Its squared modulus is at least four.
Let the five-inertia orbit of the positive label $[X]$ have
size $t\ge2$. It lifts to $t$ full representatives with a common
actual dyadic twist: five-inertia fixes their nonzero entries at
$V$, whereas replacing a representative by its fermion mate
negates those entries.

Write the lifts as $X_1,\ldots,X_t$, and put
\[
 b_i=\frac{v_{X_i}-v_{gX_i}}{\sqrt2},\qquad
 a_i=\frac{v_{X_i}+v_{gX_i}}{\sqrt2},\qquad
 v=\sum_{i=1}^t b_i,\quad a=\sum_{i=1}^t a_i.
\]
Thus $v\in B$ and $\|v\|^2=t$. Since $s$ interchanges $B,C$, its
image has zero components in $A,B,H$. On a free symmetric odd
coordinate $c_V=(v_V+v_{gV})/\sqrt2$, with $V\in E$, symmetry gives
\[
 (sv)_{c_V}=\frac{2}{\sqrt N}\sum_i S_{X_iV}
           =\frac{t\chi_V(X)}{\sqrt{32}},\qquad
 |(sv)_{c_V}|^2\ge\frac{t^2}{8}.
\]
Here $d_V=n$, $N=128n^2$, and the compatible lifts make the
$\chi_V(X_i)$ equal. Therefore
\[
             t=\|sv\|^2\ge |E|t^2/8.
\]
If $|E|=4$, then $t=2$ and equality forces every other $C$-coordinate
to vanish. The odd local groups fix each exceptional full line, so
$W_o sv=sv$. Since $W_o$ commutes with $s=WW_o$, this gives
$W_ov=v$. If $\eta$ is the common actual dyadic twist of the lifts,
then $Uv=\lambda\eta a$. Commutation of $U,W_o$ gives $W_oa=a$.
The labels $[X_i]$ lie in the positive principal two-block and $K$
preserves that block; hence $sa=Ka$ is supported there.
Choose a complementary witness $Y$ of degree $4\cdot5^j$, taking
its full principal lift at five. Its positive defect is one, so
$[Y]$ is outside the positive principal two-block. Nevertheless
\[
 (sa)_{a_Y}=\frac{2d_Y}{\sqrt N}\sum_{i=1}^2\chi_Y(X_i),
 \qquad
 \sum_{i=1}^2\chi_Y(X_i)\equiv2d_X\ne0\pmod{\mathfrak p_5}.
\]
The sum is nonzero, as is its prefactor, contradicting the support
of $sa$. Thus
\[
                              1\le|E|\le3.
\]

\paragraph{The trace contradiction.}
Work on the odd antisymmetric space $H$. Set $x=\lambda\delta$
and $T=U/x$. Its eigenvalues lie in $\{1,i,-1,-i\}$. The packet
projection satisfies
\[
 J=1-2P_E=\tfrac12((1-i)T+(1+i)T^{-1}).
\]
Unitarity and cyclicity give
\[
 \operatorname{Tr}_H(JWJW^{-1})
 =\operatorname{Im}\operatorname{Tr}_H(TWTW^{-1})
  +\operatorname{Re}\operatorname{Tr}_H(TWT^{-1}W^{-1}).
\]
The words in $U,W$ underlying these traces are represented modulo
eight by
\[
 h=\begin{pmatrix}0&1\\-1&1\end{pmatrix},\qquad
 k=\begin{pmatrix}2&1\\1&1\end{pmatrix}.
\]
Matrix multiplication gives $h^3=-I$ and $k^3=5I$ modulo eight, so
both have order dividing six. The Weyl conjugates $k$ to $k^{-1}$.
For $h$, let $P$ exchange the two coordinates. Then $PhP=h^{-1}$ and
$P(I+2h)$ has determinant $-7=1\bmod8$ and conjugates $h$ to its
inverse. The whole dyadic factor preserves $H$ since $U,W$ do.
The traces on $H$ are therefore real algebraic integers in
$\mathbb Q(\zeta_6)$, hence rational integers.
Replacing $U$ by $T$ multiplies the first trace by
$x^{-2}\in\mu_4$ and leaves the second unchanged. Thus the
displayed trace of $JWJW^{-1}$ is integral.
On $E$ the odd Weyls fix the full coordinate lines pointwise. Hence
$|W_{VY}|^2=|2S_{VY}/\sqrt N|^2=|\chi_V(Y)|^2/32$ on the
antisymmetric pair coordinates, since $d_V=n$ and $N=128n^2$.
Expanding $J=I-2P_E$ now gives
\[
 \operatorname{Tr}_H(JWJW^{-1})=\dim H-4|E|
              +\tfrac18\sum_{V,Y\in E}|\chi_V(Y)|^2 .
\]

The sum thus lies in $8\mathbb Z$. Every entry is a dyadic unit
in $\mathbb Q(\zeta_8)$ by principality and the level bound. Galois
covariance acts on either index up to the sign introduced by a
change of fermion representative. The degrees are equal, giving
symmetric squared moduli. The small dyadic sum excludes
$1\le|E|\le3$, as required.
\end{proof}

\subsection{The reduced dyadic quintic}

\begin{lemma}[A positive singleton and its odd mate]
\label{red:ef:singleton-mate}
A positive singleton two-block in an integral fermion extension
lifts to a full block consisting of two even mates and a free
odd pair of the same degree. Its row sees only the odd pair's
actual dyadic phase on odd columns.
For opposite complementary singleton witnesses at odd primes
$q,r$, all their odd mates have one common dyadic phase; every
odd $q$- or $r$-active column has that phase. Their positive
rows see only squared dyadic phase one.
\end{lemma}
\begin{proof}
The positive and signed orders have identical special fibers,
so the signed block has rank one. Hence there is one odd orbit.
Block orthogonality of the full columns $1,g$ equates its
degree-square mass with the even mass $2d_X^2$. A fixed orbit
would have degree $\sqrt2d_X$, impossible. The odd pair is
free of degree $d_X$. Local block invariance and complementary
phase preservation give its row's asserted odd support.
The same factorization on the positive space, whose vacuum
entry is nonzero, gives squared phase one on its positive row.

An opposite odd mate has degree prime to the witness's odd
prime. Choosing the witness's full principal lift makes it hit
that mate. Thus all opposite witnesses have the same visible
odd phase. The principal block $P_q$ contains only even labels,
since an odd label takes value $-1$ on $g$, whereas the vacuum
takes value $1$. Moreover the opposite odd mate $T$ has
$q\nmid d_T$, and $\chi_g(T)=-d_T\not\equiv d_T$
above $q$; hence $g\notin P_q$. The same argument with $q,r$
interchanged gives $g\notin P_r$. Prime avoidance modulo the other
odd prime therefore makes any odd active column hit a witness,
giving the last support assertion.
\end{proof}

\begin{lemma}[A single odd spin is impossible]
\label{red:ef:one-odd-spin}
Let $\mathcal M$ be an integral minimal modular extension of a category
$\mathcal F$ satisfying Lemma~\ref{loc:extension}. If every even dyadic
twist is a sign, at least two distinct odd dyadic twists occur.
\end{lemma}
\begin{proof}
Suppose $U|_C=cI$. Since $T_2^2|_A=I$, the Weyl
identity gives $(WU)^{-1}U(WU)|_A=\lambda^{-2}K^{-1}$.
Thus $K$ is scalar.
The complementary vacuum column is supported on the even
representatives $E$ with actual dyadic twist one and has
squared weights $4d_j^2/N$. The indicator construction gives an
integral self-adjoint element $b$ with $b^2=4b$ and dimension four.
Its positive restriction has every eigenvalue one by
\eqref{red:indicator-positive}, so its even part is $1+g$.
Hence $v=b-1-g$ has odd support and
\[
             gv=v,\qquad v^2=2(1+g),\qquad \dim v=2,
             \qquad \|v\|^2=2.
\]
A free fermion pair in its support would give an odd invertible:
its two coefficients agree and its total dimension is two. Lemma~\ref{loc:extension}
excludes such an invertible.
Otherwise positivity of dimensions forces $v=Y-Z$ for two
fermion-fixed self-dual odd simples. Now $vY$ has norm four
and coefficients one at $1,g$; its remaining part is an even
free pair $A'+gA'$, with positive sign by dimension. Thus
$d_Y=d_{A'}+1$. Reciprocity and $\langle vA',v\rangle=0$
give $vA'=Y+Z$; the norm of $vA'$ is two, so there are no
other terms. Hence $d_Z=d_{A'}-1$. Lemma~\ref{loc:extension} applies to the coprime
degrees $d_{A'}$ and $d_{A'}\pm1$, giving
$\chi_{A'}(Y)=\chi_{A'}(Z)=0$, contrary to
$\chi_{A'}(v)^2=4$.
\end{proof}

\begin{lemma}[Sign twists exclude a principal defect-two label]
\label{red:ef:sign-fs-square}
Let an integral fermion extension have dimension $N=64n^2$,
with $n$ odd. Suppose all even actual dyadic twists are signs.
Its positive principal two-block has no defect-two character.
\end{lemma}
\begin{proof}
Dyadic inertia fixes the even labels individually: balancing
makes their restricted values unramified at two, and opposite
twists distinguish the representatives. Since $N$ is a square,
the local torus sign is positive. Thus $W^2=I$ on $A,B,C$.
Put $J=W|_A$ and $P=WU:A\to C$. Restricting the Weyl relation
gives $UP=\lambda^{-2}PJ$. Since $J^2=I$, the odd dyadic
spins lie in a pair $\{\delta,-\delta\}$.

Choose a CRT exponent $k$ zero on the dyadic level and two
on the odd level. In the modular indicator formula for an even
$V$, the two fusion inputs are in the same fermion sector;
their dyadic twist ratio is a sign. Hence
$\nu_k(V)=\nu_2(V)$. For
$b=\sum_Y\nu_k(Y)[Y^*]$, the indicator diagonalization is
\[
                \chi_x(b)=f_x|(sT^ks^{-1})_{0x}|^2.
\]
Odd-local Bruhat at parameter $-2$ identifies the absolute
entries with a torus permutation of the odd-Weyl vacuum column.
Because $N$ is a square that column is rational. Complementary
operators preserve fermion parity and actual dyadic twist, and the
vacuum is even with twist one; its support is therefore even with
actual dyadic twist one. The
odd eigenvalues of $b$ vanish, so $gb=b$. An even pair has opposite
actual dyadic twists, and at most its twist-one representative can
have nonzero eigenvalue. Writing the relevant rational column entry
as $q_x$, the full eigenvalue is
\[
             \chi_x(b)=\frac{64n^2}{d_x^2}q_x^2
                         =\left(\frac{8nq_x}{d_x}\right)^2.
\]
In the positive quotient put
$W^+=\sum_{[V]}\nu_2(V)[V^*]$. Its coefficients are the second
indicators of even simples, hence integers. Formula~\eqref{red:indicator-positive}
now keeps both factors of two:
\[
 w'_{[x]}=\frac{\chi_x(b)+\chi_{gx}(b)}4
     \in\left\{0,\left(\frac{4nq_x}{d_x}\right)^2,
                    \left(\frac{4nq_{gx}}{d_x}\right)^2\right\}.
\]
At most one numerator term is nonzero, so these are rational squares.
They are also algebraic integers, hence ordinary integer squares. Modulo two, $(W^+)^2$ is the Casimir
$\sum_{[V]}[V][V^*]$: nonselfdual pairs cancel, and a selfdual
positive pair has a selfdual representative by opposite twists.
Thus $4\mid w'_x$ whenever its positive codegree is even.

Suppose $x$ is principal of defect two. Sign-only balancing
makes its positive field unramified at two, so
$v(V)=(\chi_x(V)-d_V)/2$ is locally integral. Orthogonality
gives
\[
 \sum_{[V]}v(V)v(V^*)=(f_x^++\Delta)/4\equiv1\pmod2,
 \qquad \Delta=16n^2.
\]
Cancel nonselfdual pairs. Frobenius in the residue field implies
$\sum_{[V]}\nu_2(V)v(V^*)\not\equiv0\pmod2$.
But this sum is $(w'_x-w'_0)/2$, divisible by two since both
eigenvalues are integer squares divisible by four. This is a
contradiction.
\end{proof}

\begin{proposition}[The reduced quintic face]
\label{red:ef:reduced-five}
There is no minimal slightly degenerate obstruction of dimension
$2^5q^br^c$ for odd primes $q<r$ and arbitrary $b,c$.
\end{proposition}
\begin{proof}
Initially the positive quotient has dimension $\Delta=16q^br^c$;
Corollary~\ref{red:small-exponents} requires $b,c>0$.
Opposite principal witnesses have common dyadic part at least
$q$ and at most four; each therefore has height two and
positive odd codegree. They are singleton two-blocks.
Their positive rows see only squared dyadic phase one, and
their odd rows see one common phase $\delta$, by
Lemma~\ref{red:ef:singleton-mate}. If one companion valuation were
odd, every occupied complementary dyadic fiber in each fermion
sector would contain an active label at that companion: its
nonsquare local torus acts by minus a permutation. Thus all even
dyadic twists would be signs and all odd spins would equal
$\delta$, contrary to Lemma~\ref{red:ef:one-odd-spin}. Hence $b,c$
are even. Put $n=q^{b/2}r^{c/2}$, so the integral extension has
dimension $64n^2$.
By Lemma~\ref{red:ef:singleton-mate}, any positive pair whose
squared dyadic twist is nontrivial is inactive at both $q,r$.
The complementary groups fix its full lines, so its positive
codegree is purely dyadic and its pair is two-principal.
The possibilities are $16,4,1$. Codegree one exhausts the row
norm; codegree four is excluded by the two distinct nonzero
even traces forced modulo $q,r$, each costing at least four.
Thus every such deep pair has degree $n$ and codegree $16$.

If there are $N_0$ deep pairs, dimension orthogonality on a
deep row, reduced modulo $r$, gives an $r$-free nonunit hit.
Pure-codegree support and the principal height bound give its
degree $2q^i$. Its crossing
with a complementary $4r^j$ witness gives an inertia orbit of
length $e\ge2$. The fixed entries on the $N_0$ deep pairs and the vacuum term in~\eqref{arith:orbit-bound} give $4eN_0<16$, so there is
exactly one deep pair $Z$. Its dyadic twist power acts on the
positive space as $R=I+(u-1)P_Z$, $u\ne1$. The trace of the
finite-group commutator $RKR^{-1}K^{-1}$ makes
\[
                    (2-u-u^{-1})|K_{ZZ}|^2
\]
an algebraic integer. Yet $|K_{ZZ}|^2=|\chi_Z(Z)|^2/16$
has dyadic valuation $-4$ by principality, whereas
$v_2(2-u-u^{-1})\le2$. This contradiction eliminates deep pairs.

All even dyadic twists are signs. Lemma~\ref{red:ef:sign-fs-square}
and the principal height bound now make every positive
two-principal degree odd. All complementary odd mates and all
odd $q$- or $r$-active columns have one common phase $\delta$.
Lemma~\ref{red:ef:one-odd-spin} gives another odd phase. Its fiber
is inactive at both odd primes, so any label $Y$ on it has
pure dyadic full codegree.

Take an odd mate $T$ of an $r$-free witness. Its degree is
$r$-free and its phase is $\delta$. The signed restrictions
of $T,Y$ are distinct and orthogonal. Reducing their
orthogonality modulo $r$, the vacuum forces a nonunit $r$-free
even pair $V$ hit by $Y$: the $r$-bearing terms vanish on $T$
by LCM integrality. Pure-codegree support places $V$ in the
positive principal two-block. Its degree is odd, hence a power
of $q$, impossible in the perfect fusion-simple positive table.
\end{proof}

\subsection{The dyadic sextic and middle quintic}

\begin{lemma}[Weight three with residues at two and three]
\label{red:ef:weight-three}
In an ordinary minimal table on dimension primes $2,3,r$, let
an LCM-normalized crossing have common weight three, residue one
at every place above two, and residue zero at every place above
three. Then it has the form $\lambda\sqrt{-3}$ for a dyadic
root of unity $\lambda$.
\end{lemma}
\begin{proof}
By Lemma~\ref{word:integrality} the two-odd-prime root configuration is unimodular; adjoining a
dyadic basis gives disjoint signed copies of it. Consequently
the crossing has an integral signed word of length at most three.
Erase its three-parts. The remaining word $w$ is divisible by
three in a three-unramified field and has house at most three.
If $w=0$, its terms are a cancelling pair, since the complementary
alphabet has no vanishing triangle. Thus the original word is
$\lambda(\alpha-\beta)$ for three-power roots $\alpha,\beta$.
If $w\ne0$, the algebraic norm and equality in the triangle
inequality give $\lambda(\alpha+\beta+\gamma)$.
In both cases inversion of the $r$-part, followed by reduction
above two, shows that $\lambda$ has trivial $r$-part: reduction
is faithful on odd-order roots, and both residues are one.
Thus $\lambda$ is dyadic.

Reduce the three-primary factors modulo two. The polynomial
$\Phi_{3^a}$ is irreducible over $\mathbb F_2$: the order of
two modulo $3^a$ is $2\cdot3^{a-1}$, as follows inductively
from $v_3(4^{3^j}-1)=j+1$. Thus every relation among the
$3^a$th roots is constant on their $\mu_3$-cosets. In the pair case,
the three terms $\alpha,\beta,1$ form the identity
$\mu_3$-coset, so $\{\alpha,\beta\}=\{\omega,\omega^{-1}\}$.
In the triple case the parity vector of
$\alpha,\beta,\gamma,1$ is a union of complete triples with
even cardinality at most four; it is therefore zero. The word
is $\lambda(1+2\xi)$ for a three-power root $\xi\ne1$.
If $\xi$ has order at least nine, its adjoint support
$\{1,\xi,\xi^{-1}\}$ has one point in each occupied
$\mu_3$-coset and is exposed. Its dilation by the three-adic
unit four meets it only at one. The common nonunit constituent
of the two adjoints would then have a scalar character value of
full degree, contradicting strictness. Thus $\xi=\omega^{\pm1}$.
Both alternatives equal a dyadic scalar times $\sqrt{-3}$.
\end{proof}

\begin{lemma}[Transfer from a two-principal ternary height-one row]
\label{red:ef:ternary-transfer}
An ordinary minimal table on primes $2,3,r$ with odd ternary
valuation has no two-principal row of degree $2^i3$.
\end{lemma}
\begin{proof}
Suppose $U$ is such a row. Let $\iota$ invert the three-power
roots, fixing the other primary factors, and put
$H=\chi_U+\chi_{U^\iota}$. These rows have the same ternary
twist. An odd-degree nonunit $B\in P_3$ has degree $3^h r^j$,
with $h,j>0$. Its normalized crossing with $U$ satisfies
Lemma~\ref{red:ef:weight-three}; it is nonzero and negated by
$\iota$. Thus $H(B)=0$, and $B^\iota\ne B$.
On every even degree, both summands of $H$ vanish above two.

If an odd-degree $J$ had a different ternary twist, principal
block orthogonality against $H$ would have vacuum term two.
Every remaining nonzero term has dyadic valuation greater than
one: its $H$ factor has positive valuation, and its $\chi_J$
factor is divisible by two by LCM integrality. This is impossible.
The vacuum therefore implies that all odd-degree labels are
ternary inactive. The odd members of $P_3$ number one modulo
two, since all but the vacuum pair under $\iota$. At odd ternary
valuation the principal mass identity is
\[
 \sum_{B\in P_3,\,\operatorname{ord}(\theta_{B,3})=3}d_B^2
       =2\sum_{B\in P_3,\,\theta_{B,3}=1}d_B^2.
\]
The left side is zero modulo four and the right side is two,
a contradiction.
\end{proof}

\begin{lemma}[A ternary phase outside a prescribed pair]
\label{red:ef:ternary-missing-phase}
In an ordinary minimal table of odd ternary valuation, the
ternary phases are not contained in $\{1,\eta\}$ for any
primitive cube root $\eta$.
\end{lemma}
\begin{proof}
Suppose they were. Let $P$ project onto the ternary inactive
coordinates and $w=W_{3'}e_0$. The nonsquare ternary torus
acts on the vacuum by its Legendre sign. The two nonzero
Bruhat identities are
\[
 PW_3U_3^tW_3^{-1}w
       =\left(\frac{-t}{3}\right)Pse_0\quad(t=1,2).
\]
The missing Fourier coefficient gives $w=-G Pse_0$ with
$|G|^2=3$. The indicator at a CRT exponent zero at three and
one on its complement therefore produces
\[
 b=3\sum_{\theta_{X,3}=1}e_X\in R,\qquad
 b=b^*,\quad b^2=3b,\quad b_1=1,\quad\dim b=3,
 \quad\|b\|^2=3.
\]
Here the coefficients are algebraic integers by the indicator
formula and rational by Galois stability of the inactive set.

Set $M=N_b$. It is symmetric integral with $M^2=3M$.
A row with diagonal one or two has exactly two off-diagonal
entries, each $\pm1$; rows with diagonal zero or three are
isolated. The unit component is a cycle. A cycle longer than
four has nonneighbors with a unique length-two path, contrary
to $M^2=3M$. A triangle has trace three, since the unit diagonal
is one and the trace is a multiple of three. It has rank one;
its positive dimension eigenvector forces all three dimensions
to equal one, contradicting perfectness.
On a four-cycle the vertex $A$ opposite the vacuum has diagonal
one and its two neighbors $Y,Z$ have diagonal two. Their
dimension equations give $d_Y,d_Z\in\{d_A-1,d_A+1\}$.
Coprime-degree vanishing gives $\chi_A(Y)=\chi_A(Z)=0$.
But $b=1\pm Y\pm Z$, so $\chi_A(b)=1$, whereas its only
possible eigenvalues are zero and three. This is impossible.
\end{proof}

\begin{theorem}[The dyadic sextic and middle quintic]
\label{red:ef:sextic-quintic-face}
There is no ordinary or reduced minimal obstruction of dimension
$2^6q^5r^4$, for odd primes $q<r$.
\end{theorem}
\begin{proof}
The reduced case is Theorem~\ref{red:fermion-six}. In the ordinary
case a dyadically active fiber contains a $q$-active column, by
Lemma~\ref{loc:fiber-activity}. Prime avoidance on $P_q$, modulo
$r$, supplies an $r$-free nonunit hitting that joint column. Its
degree is $2^iq$, with $i\le2$ by defect-zero vanishing. Opposite
principal witnesses force $2^i\ge q$, hence $q=3$ and $i=2$.

Call this degree-twelve witness $U$. Every three-free nonunit of
$P_r$ has degree $4r$ or $8r$. Its crossing with any such $U$ has
mass four and residues four at three and $r$. If $r\ge7$,
Lemma~\ref{word:two-residues} gives $1+3\xi$, forbidden by
Corollary~\ref{supp:two-position}. Thus $r=5$. The full-orbit part of
Lemma~\ref{word:first-range}, including its signed alternative,
now shows that every degree-twelve row in $P_3$ is five-ramified;
in particular the five-inertia orbit of $U$ has length $e\ge4$.

We locate an exceptional ternary phase. The five-free nonunits of
$P_2$ have ternary height two, by
Lemma~\ref{red:ef:ternary-transfer}. The odd nonunits of $P_5$
have degree $45$: height one would give a mass-three crossing with
positive residues at three and five, violating strictness and
Lemma~\ref{word:two-residues}. Both sets have ternary defect one.
Lemma~\ref{loc:visible-class} gives a common visible phase pair
$\{1,\eta\}$, with $\eta$ primitive cubic. Prime avoidance makes
every two-active or five-active column hit one of these witnesses.
Lemma~\ref{red:ef:ternary-missing-phase} supplies an occupied phase
outside the pair. The codegrees there are purely ternary, by
Lemma~\ref{red:visible-fiber}, so they are $3,27,$ or $243$.
Codegree three is excluded by Lemma~\ref{arith:small-codegrees}.

Suppose an exceptional row $Z$ has codegree $27$. All its hits lie
in $P_3$. Reducing its trace-dimension identity modulo five and
modulo two supplies separate nonzero-trace orbits of degrees
$2^i3$ and $3\cdot5^j$. Their ternary heights are one by the
principal bound. Since $d_Z=600$, symmetry and principality make
their nonzero absolute traces multiples of three. Each orbit costs
at least nine in the trace-norm bound. A third such orbit, a trace
larger than three in absolute value, or a relative field index
greater than one would exhaust $27$ before the positive vacuum
term. The two traces are therefore $\pm3$, the required orbits are
unique, and their full fields equal $K_Z$.

In particular the five-free trace witness is five-unramified.
Its dyadic height satisfies $2\le i\le3$ by complementary principal
degrees and square divisibility. Height two was just shown to be
five-ramified. At height three the trace-dimension identity modulo
five would give $1\pm3\cdot24\equiv0\pmod5$, also impossible.
Thus every exceptional row has codegree $243$ and degree $200$.

Lemma~\ref{red:quintic-mass}, with $t=3$ and $c=1$, supplies at
least nine exceptional labels. The row $U\in P_3$ hits all of them.
Their fields are five-unramified, and their LCM scale with $U$ is
three times their degree. Orbit averaging, with the vacuum term
retained in~\eqref{arith:orbit-bound}, gives $9eN_0<243$. But $eN_0\ge4\cdot9=36$, a
contradiction.
\end{proof}

\subsection{The surviving exponent pattern}
\label{red:ordered-even}

\begin{theorem}[Ordered exact reduction]\label{red:exact-reduction}
Let $p,q,r$ be distinct primes. If a braided fusion category of dimension
dividing $p^6q^6r^4$ is not weakly group-theoretical, then $r$ is the
largest prime. After interchanging $p$ and $q$, a divisor-minimal
counterexample is an ordinary integral perfect fusion-simple modular
category of exact dimension $p^6q^6r^4$, with $p<q<r$.
\end{theorem}
\begin{proof}
Choose a divisor-minimal counterexample. Lemma~\ref{loc:minimal}
leaves an ordinary modular table or a positive fermion quotient.
By Corollary~\ref{red:small-exponents}, all three primes occur,
each with valuation at least four in the category dimension.
Theorem~\ref{red:ordered-quartic} excludes a valuation at most four
at either of the two smaller primes. Hence the largest prime is $r$,
its valuation is four, and the other two valuations are five or six.

In odd dimension the table is ordinary.
Propositions~\ref{red:least-quintic} and~\ref{red:middle-quintic}
successively exclude valuation five at $p$ and $q$. In even dimension, dyadic valuation five is excluded in both
branches by Propositions~\ref{red:ef:ordinary-five} and
\ref{red:ef:reduced-five}. At dyadic valuation six, middle
valuation five is excluded by Theorem~\ref{red:ef:sextic-quintic-face}.
Thus both sextic exponents are attained. The remaining
slightly degenerate possibility is excluded by Theorem~\ref{red:fermion-six}. Only the stated ordinary exact dimension remains.
\end{proof}

\section{The even and odd dimension bounds}\label{sec:boundaries}

\subsection{Prime-power trace budgets}

\begin{proposition}[Prime-power trace budget]\label{arith:primepower-trace}
Suppose an ordinary perfect fusion-simple integral modular category has exactly three
dimension primes $p,u,v$, and $f_Z=p^{2a}$ for a nonunit $Z$. Then
\[
 Z\in P_p,\qquad p^{2a}\ge[K_Z:\mathbb Q]+2p^2.
\]
If $v_p(D)=2m$ and $v_p(d_Z)=m-a>0$, a nonzero absolute trace
at a $p$-bearing nonunit degree $d_X$ requires
\[
                    a\ge2,\qquad v_p(d_X)\le m-2.
\]
In particular, when $m=3,a=2$, there are nonzero-trace labels
of degrees $pu^i$ and $pv^j$.
\end{proposition}
\begin{proof}
At every prime other than $p$, the idempotent of $Z$ is
locally integral, so its full residue block is a singleton.
The complementary local Weyl factors therefore preserve its line.
Since the global $S$-column of $Z$ has a nonzero unit entry,
local block invariance forces its $p$-block to be principal
(Theorem~\ref{loc:block} and Lemma~\ref{loc:weyl-line}).

The trace-dimension identity \eqref{arith:trace-identities}, reduced
modulo $u$ and $v$, supplies two distinct nonzero-trace nonunit
orbits, respectively on $u$-free and $v$-free degrees. Their degrees
are $p$-bearing because prime-power dimensions are excluded.
Principality and \eqref{arith:all-places} make both integer traces
divisible by $p$. In the trace-norm identity they each cost at
least $p^2$, while the unit costs $[K_Z:\mathbb Q]$.

Put $h=m-a$ and $j=v_p(d_X)$. Principal support also puts $X$
in $P_p$ when its trace is nonzero. If $j\ge h$, write
\[
             \chi_Z(X)=p^{j-h}c\,\chi_X(Z),\qquad v_p(c)=0.
\]
The factor $c$ is rational and a $p$-adic unit. Since $p\mid d_Z$
and $X\in P_p$, each conjugate of $c\chi_X(Z)$ reduces to zero at a
place over $p$. Its trace is rational, locally integral and has
positive valuation, hence is divisible by $p$ in $\mathbb Z_p$.
For $j<h$, use $Z\in P_p$ and $p\mid d_X$ directly on $\chi_Z(X)$.
Thus the integer trace $T_X$ satisfies
$v_p(T_X)\ge1+\max(j-h,0)$.
The positive unit term in the trace-norm identity gives
$|T_X|<p^a$, hence $a\ge2$ and $j\le h+a-2=m-2$.
At $m=3,a=2$, the two selected degrees are therefore $pu^i,pv^j$.
\end{proof}

\subsection{The isolated dyadic sextic}

\begin{theorem}[No codegree sixteen at dyadic valuation six]
\label{red:codegree-sixteen}
Let $\mathcal C$ be an ordinary perfect fusion-simple integral
modular category with exactly three dimension primes
$2<q<r$ and $v_2(\dim\mathcal C)=6$. No simple has formal
codegree $16$. Consequently every purely dyadic formal
codegree is $64$.
\end{theorem}
\begin{proof}
Suppose $f_Z=16$. Then
\[
 D=64M^2,\qquad M=q^b r^c,\qquad b,c\ge1,
 \qquad d_Z=2M.
\]
The conductor bound makes $K_Z$ dyadic, and
Proposition~\ref{arith:primepower-trace} puts $Z$ in $P_2$.
Its row therefore meets every odd-degree label. Their number
is a positive multiple of four, since $D=\sum_Xd_X^2\equiv0\pmod4$
and the vacuum has odd degree. Each nonzero integral entry has
Galois mean square at least one. These labels contribute at least
four to the row norm.

The trace-dimension identity modulo $q$ and modulo $r$,
together with Proposition~\ref{arith:primepower-trace}, supplies
two distinct nonzero-trace even orbits, represented by
\[
                   d_X=2q^j,\qquad d_Y=2r^k.
\]
Their heights are one because the dyadic dimension valuation
is six and their nonzero traces have height at most one.
For either orbit the entry at $Z$ vanishes at the dyadic
place. A nonzero dyadic integer $\alpha$ vanishing at the dyadic place
has absolute norm at least two. If $h=[\QQ(\alpha):\QQ]$, the sum $E(\alpha)$
of the squared absolute values of its $h$ conjugates satisfies
\[
 E(\alpha)\ge h\,|N_{\QQ(\alpha)/\QQ}(\alpha)|^{2/h}
             \ge h\,2^{2/h}\ge4,
\]
with equality possible only for $h=1,2$.

Cross $X$ with a $q$-free nonunit of $P_r$. Such a label
has even degree, since nonunit prime-power degrees are
excluded. The normalized crossing has mass two, residue
two at $r$, and strict house below two. The large-residue
word lemma makes it the sum of two distinct $r$-primary
roots. Its tame stabilizer has order at most two, so the
$r$-inertia orbit of $X$ has size at least $(r-1)/2\ge2$.
The entry field $F=\mathbb Q(\chi_Z(X))$ is dyadic and hence
fixed by $r$-inertia. Thus $[K_X:F]\ge2$, and this orbit
spends at least eight in the $Z$-row norm: the exact sum
over the full $X$-orbit is
$[K_X:F]E(\chi_Z(X))$. The $Y$-orbit
spends at least four. The three contributions have disjoint supports. Therefore
\[
 16=f_Z\ \ge\
 \underbrace{4}_{d\ {\rm odd}}+
 \underbrace{8}_{\operatorname{Gal}\cdot X}+
 \underbrace{4}_{\operatorname{Gal}\cdot Y}=16 .
\]
Every inequality is an equality. There are exactly four
odd labels, every odd entry is a root of unity, and the only
even entries in the support are the two chosen orbits,
with total squared norms eight and four, respectively.
For the first orbit $[K_X:F]=2$ and its entry energy is four;
for the second the relative index is one and its entry
energy is four.

Equality in the energy bound makes both entry fields have degree
at most two. The two even orbits therefore have sizes at most four and
two. Cauchy's inequality bounds their sums of absolute
entry values by $\sqrt{32}$ and $\sqrt8$.

The modular inequality \eqref{arith:sts-bound} gives
\[
  \sqrt{f_Z}\,d_Z
       \le 1+\sum_{T\ne\one}d_T|\chi_Z(T)|.
\]
Each odd nonunit has degree at most $M$, the $X$-orbit has
degree at most $2M/r^c$, and the $Y$-orbit has degree at most
$2M/q^b$. Consequently
\[
 8M\le 1+3M+\frac{2M}{r^c}\sqrt{32}
                    +\frac{2M}{q^b}\sqrt8
      <1+\left(3+\frac{12}{5}+2\right)M
      =1+\frac{37}{5}M<8M.
\]
Here $q^b\ge3$, $r^c\ge5$, $M\ge15$, and
$\sqrt2<3/2$. This contradiction proves the theorem.
For the final assertion, a purely dyadic codegree makes
the other dimension valuations even; square divisibility
then leaves only $1,4,16,64$. The small-codegree exclusion
and the result just proved leave $64$.
\end{proof}

\begin{theorem}[Two cannot be isolated at sextic valuation]
\label{red:dyadic-isolation}
Let an ordinary minimal table have dimension primes
$2<q<r$ and $v_2(D)=6$. Some twist order is divisible by $2q$ or $2r$.
\end{theorem}
\begin{proof}
Suppose two is isolated. Proposition~\ref{loc:isolation}
gives $D=64n^2$ with $n$ odd.  Every dyadically active label
has pure dyadic codegree. Theorem~\ref{red:codegree-sixteen}
therefore makes every active degree $n$, of codegree $64$.
The isolated-column formula gives an integral vector

\begin{equation}\label{eq:dyadic-column}
u=8W_2v_0,\qquad \|u\|^2=64,\qquad
 v_2(u_X)=v_2(d_X)\quad\text{when }v_2(d_X)<3,
\end{equation}

and puts all these lower-height labels in \(P_2\).
The sign in the normalization of $W_2$ affects neither the norm
nor the coordinate valuations.

An \(r\)-free label \(X\) of degree \(2q^j\) has
\(r\)-inertia orbit length \(t\ge(r-1)/2\ge2\).
Indeed, pair it with a \(q\)-free nonunit in \(P_r\),
of degree \(2^h r^l\).  The normalized crossing has
mass two, strict house below two, and residue two at \(r\).
Since \(r>4\), Lemma~\ref{word:large-residue} gives two
distinct \(r\)-power roots.  Their tame stabilizer has order
at most two.  This also proves that \(K_X\) ramifies at \(r\).

Let $N$ count the active labels. Choose the greatest dyadic
twist order among them, and let $L$ be the number at that order.
An exponent can turn these phases into $-1$ and every smaller
phase into $1$, while killing the odd primary factors. The
Gauss sum is $n^2(64-2L)$ since the $L$ selected labels
all have degree $n$. The regular-object indicator is the squared
modulus of this sum divided by $D$, namely $n^2(64-2L)^2/64$.
It is a rational algebraic integer. Since $n$ is odd, $8\mid64-2L$,
so $4\mid L$ and $N\ge4$.
Prime avoidance gives an \(r\)-free even \(X\in P_2\)
of degree \(2^i q^j\), where \(i=1\) or \(2\).
It hits all \(N\) active labels.  Let \(t\) be the length
of its $r$-inertia orbit. On every active column the entry field is
dyadic, hence fixed by $r$-inertia, and the LCM of the two degrees is
$2^i n$. Applying \eqref{arith:orbit-bound} to these columns and
retaining the positive vacuum term gives

\begin{equation}\label{eq:dyadic-count}
4^i tN<64.
\end{equation}

This first forces \(i=1\), and then the orbit estimate gives
\(t\ge2\) and \(N\le7\).
The vacuum Gauss formula and \eqref{eq:dyadic-column} now yield
\[
 |u_0|=\frac18\left|64-\sum_{Z\ {\rm active}}
                       (1-\theta_{Z,2})\right|
       \ge\frac{64-2N}{8}>6.
\]
Thus \(u_0^2\ge49\).  The \(N\) odd active coordinates cost
at least four in \(\|u\|^2\); the \(t\ge2\) even
\(r\)-free coordinates cost at least eight; and an even
\(q\)-free member of \(P_2\) costs at least four.
The four disjoint contributions give $64\ge49+4+8+4=65$.
\end{proof}

\begin{corollary}[A crossing of common height two at two]
\label{red:dyadic-crossing}
In an ordinary minimal table of dimension \(2^6q^6r^4\),
some opposite \(r\)-free member of \(P_q\) and \(q\)-free
member of \(P_r\) have common dyadic part four.  Necessarily \(q=3\).
\end{corollary}
\begin{proof}
Complementarity and square divisibility make the common
dyadic part four or eight.  If every opposite pair had
part eight, every witness on both sides would have
odd codegree.  Every \(q\)-active column hits an
\(r\)-free witness in \(P_q\), by prime avoidance, and
would therefore be dyadically inactive by defect-zero
vanishing.  The same argument at \(r\) isolates two,
contrary to the theorem.  The complementary-degree bound
(Lemma~\ref{red:complementary-degrees}) gives $q\le4$, hence $q=3$.
\end{proof}

\begin{corollary}[The unconditional even-family exclusion]
\label{red:even-boundary}
In the ordered exact family with $p=2$, every member with $q>3$
or $r>5$ is weakly group-theoretical. Thus a minimal even
obstruction can have only the exceptional assignment $q=3,r=5$.
\end{corollary}
\begin{proof}
Corollary~\ref{red:dyadic-crossing} forces $q=3$ and a
compulsory crossing whose common weight is four. For $r\ge7$,
Lemma~\ref{word:two-residues} and strictness give its
actual word $z=1+3\xi$, where $\xi$ is a nonidentity
$r$-primary root. The exceptional signed branch would require
$r=5$. Since
\[
                         \chi_Y(X)=\frac{d_X}{4}z,
\]
this is a full positive two-position character word, of mass
$d_X$. Corollary~\ref{supp:two-position} excludes it at arbitrary primary
conductor.
\end{proof}

\subsection{The least odd prime and the transfer to short words}

\begin{proposition}[An isolated sextic budget]\label{supp:isolated-sextic}
Let $D=p^6q^6r^4$, with $p<q<r$ odd, be the dimension of
an ordinary perfect fusion-simple integral modular category.
If $p$ is isolated in the twist support, put
\[
 e_t=\frac{t-1}{\max\{a:a\mid t-1,\ a\le p\}},\qquad t=q,r.
\]
Then
\begin{equation}\label{supp:eq:sextic-budget}
                         (p-1)(e_q+e_r)<p^2.
\end{equation}
In particular $p$ cannot be isolated if $r>p^2$.
\end{proposition}
\begin{proof}
Write $D=p^6n^2$, $n=q^3r^2$. Proposition~\ref{loc:isolation}
and~\ref{arith:primepower-trace}, together with
Lemma~\ref{arith:principal-support}, imply that every
$p$-active $Z$ has degree $p^jn$, $j\in\{0,1\}$,
$p$-primary character field, and row supported on $P_p$.

Suppose first that some $Z$ has $d_Z=pn$. The trace-height
part of Proposition~\ref{arith:primepower-trace} supplies, for
each $t=q,r$, a $t$-free nonunit $X$ of degree $pu^i$,
where $u$ is the other companion, with nonzero absolute
entry trace $T_X$ divisible by $p$. Thus
\[
                         d_Z^2 T_X^2\ge p^4n^2.
\]
Cross $X$ with a $u$-free nonunit in $P_t$. The normalized
crossing has weight $p$ and residue $p$ at $t$.
Since $(t-p)(p-1)>p$, Lemma~\ref{word:one-residue}
makes its actual word an unsigned integral $t$-word
of weight $p$. Strictness supplies a nonidentity phase.
Its tame stabilizer has order dividing $t-1$ and at
most $p$, as it acts freely on the nonidentity phases.
Consequently $[K_X:\mathbb Q]\ge e_t$ and $K_X$
ramifies at $t$. Separation makes it unramified at $p$,
so the entry field $F_X\subseteq K_X\cap K_Z$ is
$\mathbb Q$.

The square Galois action on the active twist supplies
at least $(p-1)/2$ distinct twists in the orbit of $Z$.
Complex conjugation pairs distinct dual labels of the
same twist by Lemma~\ref{loc:odd-duality}; hence
$[K_Z:\mathbb Q]\ge p-1$. The two chosen $X$-orbits
are disjoint, since one is $q$-free and the other
$r$-free and prime-power nonunit degrees are impossible.
The trace norm of Lemma~\ref{arith:trace-energy},
multiplied by $d_Z^2$, therefore spends at least
\[
                      p^4n^2(p-1)(e_q+e_r)
\]
on those two orbits. The vacuum spends a positive amount,
so comparison with $D$ gives \eqref{supp:eq:sextic-budget}.

It remains that every active degree equals $n$. Let $N$
count these labels. Equation~\eqref{loc:isolated-vacuum} gives
\[
 K=\frac1{p^3}\sum_{Z\ p\text{-active}}(1-\theta_{Z,p})
                  \in\mathbb Z.
\]
It is positive. Pairing equal-twist dual labels shows
$p^3K\in2\mathcal O$, hence $K$ is even and $K\ge2$.
The normalized trace of $1-\theta_{Z,p}$ is at most
$p/(p-1)$, so
\begin{equation}\label{supp:eq:active-count}
                         N\ge2p^2(p-1).
\end{equation}
Every $X\in P_p$ meets every active $Z$, by principality
and $p\nmid n$. If such an $X$ had $p$-height two,
LCM integrality would give nonzero integral entries
$S_{XZ}/(p^2n)$. Their mean squared moduli are at least
one. The row norm, with its positive vacuum term, would
give $Np^4n^2<D$, contrary to \eqref{supp:eq:active-count}.

For each $t=q,r$, principal avoidance supplies a $t$-free
nonunit $X\in P_p$, necessarily of $p$-height one by
Corollary~\ref{arith:principal-height} and the preceding paragraph.
The same weight-$p$ crossing gives $[K_X:\mathbb Q]\ge e_t$
and makes $K_X$ unramified at $p$. On every active column
the integral entry $S_{XZ}/(pn)$ therefore lies in
$K_X\cap K_Z=\mathbb Q$ and is a nonzero integer.
The $[K_X:\QQ]$ distinct conjugate rows agree at these $N$
columns because the entries are rational. By orthogonality their
average has squared norm $D/[K_X:\QQ]$ in the unnormalized table.
Retaining the active entries, each of absolute value at least $pn$,
and the positive vacuum term in \eqref{arith:orbit-bound} gives
$N p^2 e_t\le N p^2 [K_X:\QQ]<p^6$. Together with
\eqref{supp:eq:active-count}, this yields
$(p-1)e_t<p^2/2$. Add the two inequalities to obtain
\eqref{supp:eq:sextic-budget}.

Finally, if $r>p^2$, the integer $e_r\ge(r-1)/p>p$
is at least $p+1$. With $e_q\ge1$, the left side of
\eqref{supp:eq:sextic-budget} is at least
$(p-1)(p+2)>p^2$.
\end{proof}

\begin{theorem}[Transfer from a short word]\label{supp:transfer}
For odd $p<q<r$, put
\[
 m=\left\lfloor\frac{\min(r,p^2)}q\right\rfloor,
       \qquad E=m(m+1)/2,
       \qquad H=(r-1)/2.
\]
If an ordinary perfect fusion-simple integral modular category
has dimension $p^6q^6r^4$, some nonunit label has, at a nonunit
row, a full positive character word on at most $m+1$ positions
whose phase ratios have $r$-power order. In particular $m\ge3$, and some
$1\le h\le E$ satisfies \eqref{supp:eq:scalar}.
\end{theorem}
\begin{proof}
If $r<p^2$, then $(p-1)q>p^2>r$ gives $m+1<p$.
Proposition~\ref{supp:quartic-seed} supplies the word.
A common scalar root in its phases disappears in the adjoint word.

Suppose $r>p^2$. Pair any $r$-free nonunit of $P_q$
with any $q$-free nonunit of $P_r$. Their degrees are
$p^iq^j,p^kr$ with $1\le i,k\le3$. Their normalized
crossing has two positive residues and strict house
less than its weight $p^{\min(i,k)}$.
Lemma~\ref{word:two-residues} excludes weight $p<q$.
A weight-$p^2$ crossing, since $p^2<r$, gives the
required word on at most $1+\lfloor p^2/q\rfloor$
positions. Its exceptional signed case would require
$r=p^2+1$, impossible for odd primes.

If no such crossing has weight $p^2$, every one of these
complementary witnesses has $p$-height three. On every
$q$-active column, principal avoidance on $P_q$ modulo
$r$ supplies an $r$-free hit of this type. Its $p$-unit
codegree forces that column to be $p$-inactive. The
same argument on $P_r$ separates $p$ from $r$.
Thus $p$ is isolated, contradicting
Proposition~\ref{supp:isolated-sextic}.

This proves existence of the word. One position
contradicts strictness. Two positions are impossible
since $r\ge7$. If $m=2$, then $r>2q\ge10$ and
Corollary~\ref{supp:three-position} applies.
Thus $m\ge3$. Automatic exposure and
Theorem~\ref{supp:descent} give the scalar condition.
\end{proof}

\subsection{The dimension bound and its arithmetic endpoint}

\begin{corollary}[The odd dimension bound]\label{supp:odd-frontier}
Every braided fusion category of odd dimension in the family
$p^6q^6r^4$, with distinct primes and any assignment of the
quartic prime, is weakly group-theoretical whenever its dimension
is less than
\[
 D_0=5^6\,7^6\,23^4
       =514\,422\,090\,765\,625=(22\,680\,875)^2.
\]
\end{corollary}
\begin{proof}
Theorem~\ref{red:exact-reduction} leaves
an ordinary perfect fusion-simple integral modular category
of exact dimension $p^6q^6r^4$ with $p<q<r$.
Theorem~\ref{supp:transfer} gives
$3q\le\min(r,p^2)$. For $p=3$ this is impossible.
For $p=5$, the inequalities $5<q\le25/3$ force $q=7$,
and $r\ge21$ forces the prime $r\ge23$.
For $p\ge7$ one has $q\ge11$ and $r\ge3q\ge33>23$.
Thus every surviving triple has dimension at least $D_0$.
\end{proof}

\begin{proposition}[Four positions at the odd boundary]\label{supp:eleven-mask}
At $(p,q,r)=(5,7,23)$ an explicit four-position word of mass $23$
realizes equality in the scalar condition \eqref{supp:eq:scalar} with
five nonzero differences up to sign. In a category of dimension $D_0$
satisfying the hypotheses of Theorem~\ref{supp:transfer}, every
four-position witness of that theorem has six distinct nonzero
differences up to sign.
\end{proposition}
\begin{proof}
Here $m=3$, $E=6$ and $H=11$. Write $\zeta=\exp(2\pi i/23)$.
The word
\[
                    2+7(\zeta^2+\zeta^8+\zeta^{16})
\]
has mass $23$, strict house below $23$, and residues $23$ at seven
and zero at $23$. Its support $\{0,2,8,16\}\subset\mathbb F_{23}$
has opposite differences represented by $2,6,7,8,9$. The powers of
$2$ identify $\mathbb Z/11$ with $\mathbb F_{23}^\times/\{\pm1\}$;
these five differences have exponents $A=\{1,3,4,5,9\}$.
They are precisely the nonzero squares modulo $11$. Thus
$c(t)=|A\cap(A+t)|$ is invariant under multiplication of $t$ by a
square, and $c(t)=c(-t)$. Since $-1$ is a nonsquare, $c(t)$ is
constant for $t\ne0$. Counting ordered distinct pairs gives
$c(t)=5\cdot4/10=2$. This is equality in
\eqref{supp:eq:scalar} at $h=5$, where $u(5)=2$.

Now fix a row carrying a witness of Theorem~\ref{supp:transfer}.
Take the root alphabet generated by the actual twists. It contains the
phases of the balancing witness, and its only possible primes besides $23$
are $5$ and $7$ by categorical Cauchy \cite[Theorem~3.9]{BNRWFinite}.
Minimize $h=(|S-S|-1)/2$ over all full words of nonunits at this row
with $23$-primary phase ratios, as in Theorem~\ref{supp:descent}.
The initial word gives $h\le6$; \eqref{supp:eq:scalar} leaves only
$h=5,6$. Suppose $h=5$, and put $D'=S-S$ for a minimizing support.
Write $h_i$ for the number of its five nonzero opposite pairs in the
$i$th tame orbit. The identity dilation contributes five to the overlap
sum of Theorem~\ref{supp:descent}, and each of the ten others
contributes at least two. Hence
\[
           25\le\sum_b k_b=\sum_i h_i^2\le25.
\]
All five pairs belong to a single tame orbit, each nonidentity overlap
has size two in opposite pairs, and all five pairs have the same
order $23^j$. Both $D'$ and every unit dilate of $D'$ are exposed by
Lemma~\ref{supp:exposure}, since $5|D'|=55<4\cdot23$.

If $j>1$, the unit $1+23$, extended by the identity at the
complementary primes, moves this tame orbit disjointly from itself.
Indeed, an intersection would give $1+23$ equal up to sign to a tame
unit modulo $23^j$. Its reduction modulo $23$ would determine that
signed tame unit to be $1$, which is impossible modulo $23^j$.
Thus the two adjoint supports meet only at zero.
Lemma~\ref{supp:common} gives a full one-position word for a nonunit,
contrary to strictness.

If $j=1$, index the opposite pairs by powers of $2$ as above.
Let $A'$ be their five exponents. Since $|A'\cap(A'+1)|=2$,
there are distinct $a,b$ with $a,a+1,b,b+1\in A'$.
The overlap for the nonzero shift $b-a$ has size two and contains
$b,b+1$, so consists exactly of these exponents. In the phase group
the corresponding adjoint supports intersect in
$\{0,\pm x,\pm2x\}$. All its nonzero differences up to sign are
among $x,2x,3x,4x$. Lemma~\ref{supp:common} now supplies a full
word for a nonunit with at most four opposite differences,
contradicting minimality.

It follows that $h=6$. A support on at most four positions has at most
six opposite differences, so the initial support must attain this
maximum. Thus its six unordered pairs give six distinct phase ratios
up to inversion. Existence at $D_0$ remains open.
\end{proof}

This conclusion is a necessary support condition at arbitrary
$23$-primary conductor. It does not reduce a six-difference witness
to conductor $23$, classify its coefficients, or construct modular data.
A finite search on $\mathbb F_{23}$ would therefore test only the
prime-level subcase and cannot settle the endpoint.

\begin{proof}[Proof of Theorem~\ref{thm:main}]
If a category on a proper divisor were a counterexample, one of
minimal divisor dimension would contradict
Theorem~\ref{red:exact-reduction}. All proper divisors are therefore
weakly group-theoretical. The same theorem and
Lemma~\ref{loc:minimal} give the assertion at the full dimension.
The even conclusion is Corollary~\ref{red:even-boundary}; the odd
conclusion is Corollary~\ref{supp:odd-frontier}.
\end{proof}

\begin{thebibliography}{99}
\small
\bibitem{BH}
M.~Beck and S.~Ho\c{s}ten, \href{https://doi.org/10.4310/mrl.2006.v13.n4.a10}{\emph{Cyclotomic polytopes and growth series of cyclotomic lattices}}, Math. Res. Lett. \textbf{13} (2006), 607--622.

\bibitem{BNRWFinite}
P.~Bruillard, S.-H.~Ng, E.~C.~Rowell, and Z.~Wang,
\href{https://doi.org/10.1090/jams/842}{\emph{Rank-finiteness for modular categories}},
J. Amer. Math. Soc. \textbf{29} (2016), no.~3, 857--881.

\bibitem{BPBurnside}
S.~Burciu and S.~Palcoux,
\href{https://doi.org/10.1142/S0219199726500070}{\emph{Burnside-type results for fusion categories}},
Commun. Contemp. Math., article 2650007.

\bibitem{DongWeak}
J.~Dong,
\href{https://doi.org/10.1007/s11401-026-0019-6}{\emph{Classification of certain weakly integral fusion categories}},
Chin. Ann. Math. Ser. B \textbf{47} (2026), 555--572.

\bibitem{DLNCongruence}
C.~Dong, X.~Lin, and S.-H.~Ng,
\href{https://doi.org/10.2140/ant.2015.9.2121}{\emph{Congruence property in conformal field theory}},
Algebra \& Number Theory \textbf{9} (2015), no.~9, 2121--2166.

\bibitem{DGNO}
V.~Drinfeld, S.~Gelaki, D.~Nikshych, and V.~Ostrik,
\href{https://doi.org/10.1007/s00029-010-0017-z}{\emph{On braided fusion categories I}},
Selecta Math. (N.S.) \textbf{16} (2010), 1--119.

\bibitem{EGNO}
P.~Etingof, S.~Gelaki, D.~Nikshych, and V.~Ostrik,
\href{https://math.mit.edu/~etingof/egnobookfinal.pdf}{\emph{Tensor Categories}},
Mathematical Surveys and Monographs, vol.~205,
American Mathematical Society, 2015.

\bibitem{ENOsolvable}
P.~Etingof, D.~Nikshych, and V.~Ostrik,
\href{https://doi.org/10.1016/j.aim.2010.06.009}{\emph{Weakly group-theoretical and solvable fusion categories}},
Adv. Math. \textbf{226} (2011), no.~1, 176--205.

\bibitem{GNnil}
S.~Gelaki and D.~Nikshych, \href{https://doi.org/10.1016/j.aim.2007.08.001}{\emph{Nilpotent fusion categories}}, Adv. Math. \textbf{217} (2008), 1053--1071.

\bibitem{JFR}
T.~Johnson-Freyd and D.~Reutter,
\href{https://doi.org/10.1090/jams/1023}{\emph{Minimal nondegenerate extensions}},
J. Amer. Math. Soc. \textbf{37} (2024), 81--150.

\bibitem{NSTV}
S.~Nasseh, S.~Sather-Wagstaff, R.~Takahashi and K.~VandeBogert,
\href{https://doi.org/10.1016/j.jpaa.2018.06.006}{\emph{Applications and homological properties of local rings with
decomposable maximal ideals}}, J. Pure Appl. Algebra \textbf{223}
(2019), 1272--1287.

\bibitem{NataleWGT}
S.~Natale,
\href{https://doi.org/10.4171/jncg/177}{\emph{On weakly group-theoretical non-degenerate braided fusion categories}},
J. Noncommut. Geom. \textbf{8} (2014), 1043--1060.

\bibitem{NataleCore}
S.~Natale,
\href{https://doi.org/10.1142/S0129167X1850012X}{\emph{The core of a weakly group-theoretical braided fusion category}},
Internat. J. Math. \textbf{29} (2018), no.~2, 1850012.

\bibitem{NPdegrees}
S.~Natale and J.~Plavnik,
\href{https://doi.org/10.2140/ant.2012.6.1171}{\emph{On fusion categories with few irreducible degrees}},
Algebra Number Theory \textbf{6} (2012), 1171--1197.

\bibitem{NSind}
S.-H.~Ng and P.~Schauenburg, \href{https://doi.org/10.1007/s00220-010-1096-6}{\emph{Congruence subgroups and generalized Frobenius--Schur indicators}}, Comm. Math. Phys. \textbf{300} (2010), 1--46.

\bibitem{YuSlightlyDegenerate}
Z.~Yu,
\href{https://doi.org/10.1016/j.jalgebra.2020.04.022}{\emph{On slightly degenerate fusion categories}},
J. Algebra \textbf{559} (2020), 408--431.
\end{thebibliography}
\end{document}
